% =====================================================================
% Preámbulo común — Diapositivas Magistrales, Econometría Avanzada
% =====================================================================
% Se incluye con % =====================================================================
% Preámbulo común — Diapositivas Magistrales, Econometría Avanzada
% =====================================================================
% Se incluye con % =====================================================================
% Preámbulo común — Diapositivas Magistrales, Econometría Avanzada
% =====================================================================
% Se incluye con % =====================================================================
% Preámbulo común — Diapositivas Magistrales, Econometría Avanzada
% =====================================================================
% Se incluye con \input{preambulo/preambulo-magistral} DESPUÉS de
% \documentclass[<tam>pt, xcolor=dvipsnames]{beamer} en cada archivo.
%
% Lo único que debe quedar en cada .tex individual es:
%   1) \documentclass[...]{beamer}      (el tamaño de fuente puede variar)
%   2) \input{preambulo/preambulo-magistral}
%   3) \subtitle{...}                   (título de la clase puntual)
% Todo lo demás (paquetes, macros, tema, título, autor, hypersetup, etc.)
% vive aquí para poder hacer cambios globales en un solo lugar.
% =====================================================================

% ---------------------------------------------------------------------
% 1. Codificación, idioma y fuentes
% ---------------------------------------------------------------------
\usepackage{etex}
\usepackage[utf8]{inputenc}
\usepackage[T1]{fontenc}
\usepackage[spanish,es-nodecimaldot]{babel}
\usepackage{lmodern}
\usepackage{type1cm}
\usepackage{textcomp}

% ---------------------------------------------------------------------
% 2. Matemáticas y símbolos
% ---------------------------------------------------------------------
\usepackage{amsmath,amssymb,amsthm,amsfonts}
\usepackage{bbm}
\usepackage{bm}

% ---------------------------------------------------------------------
% 3. Bibliografía y citas
% ---------------------------------------------------------------------
\usepackage[round]{natbib}
\bibliographystyle{erae}
\usepackage{bibentry}
\usepackage{url}
\let\htmlstyloaded=\relax

% Las citas se muestran en marrón para distinguirlas del contenido.
\let\oldcitet=\citet
\renewcommand{\citet}[1]{\textcolor{brown}{\oldcitet{#1}}}
\let\oldcitep=\citep
\renewcommand{\citep}[1]{\textcolor{brown}{\oldcitep{#1}}}
\let\oldcite=\cite
\renewcommand{\cite}[1]{\textcolor{brown}{\oldcite{#1}}}

% ---------------------------------------------------------------------
% 4. Figuras, tablas y colores
% ---------------------------------------------------------------------
\usepackage{float}
\usepackage{graphicx}
\usepackage{subcaption}
\usepackage[font={footnotesize}]{caption}
\usepackage{lscape}
\usepackage{longtable}
\usepackage{booktabs}
\usepackage{multirow}
\usepackage{array}
\usepackage[flushleft]{threeparttable}
\usepackage{color, colortbl}
\usepackage{eso-pic}
\usepackage{chngcntr}
\usepackage{adjustbox}
\usepackage{pdfpages}
\usepackage{media9}

% siunitx: usado para alinear/formatear números en tablas
\usepackage{siunitx}
\sisetup{
	detect-mode,
	tight-spacing		= true,
	group-digits		= false ,
	input-signs		= ,
	input-symbols		= ( ) [ ] - + *,
	input-open-uncertainty	= ,
	input-close-uncertainty	= ,
	table-align-text-post	= false
}

% ---------------------------------------------------------------------
% 5. TikZ / PGFPlots
% ---------------------------------------------------------------------
\usepackage{tikz}
\usepackage{pgfplots}
\usepgfplotslibrary{dateplot}
\usetikzlibrary{arrows,calc, patterns, positioning, shapes.geometric, decorations.pathreplacing,decorations.markings}

% ---------------------------------------------------------------------
% 6. Misceláneos de formato
% ---------------------------------------------------------------------
\usepackage{setspace}
\usepackage{lipsum}
\usepackage{framed}
\usepackage{mdframed}
\usepackage{fancyhdr}
\usepackage{beamerfoils}
\usepackage{pifont}
\usepackage{ragged2e}
\usepackage{xspace}
\usepackage{appendixnumberbeamer}
\usepackage[scale=2]{ccicons}
\usepackage[most]{tcolorbox}
\usepackage{enumitem}
\setbeamertemplate{note page}[plain]

% ---------------------------------------------------------------------
% 7. Hyperref (debe cargarse después de la mayoría de paquetes)
% ---------------------------------------------------------------------
\usepackage{hyperref}
\hypersetup {
	pdftitle={General},
	pdfauthor={Manuel Fern\'{a}ndez},
	colorlinks=true,
	citecolor=brown,
	linkcolor=black,
	urlcolor=blue
}
\makeatletter
\renewcommand{\hyper@natlinkstart}[1]{}
\renewcommand{\hyper@natlinkend}{}
\renewcommand{\hyper@natlinkbreak}[2]{#1}
\makeatother

% =====================================================================
% 8. Tema Beamer (metropolis)
% =====================================================================
\usetheme{metropolis}
\newcommand{\themename}{\textbf{\textsc{metropolis}}\xspace}
\metroset{sectionpage=none}

\setbeamerfont{citep}{family=\rm}
\setbeamerfont{caption}{size=\scriptsize}
\setbeamerfont{normal text}{size=7pt}

%\setcounter{tocdepth}{1}
\setbeamertemplate{section in toc}[sections numbered]
\AtBeginSection[]{
  \begin{frame}
  \vfill
  \centering
  \begin{beamercolorbox}[sep=8pt,center,shadow=true,rounded=true]{title}
    \usebeamerfont{title}\insertsectionhead\par%
  \end{beamercolorbox}
  \vfill
  \end{frame}
}
% NOTA: \setbeamercolor{button}, \setbeamercovered{invisible} y
% \AtBeginSubsection NO están aquí porque sólo algunas clases los usaban
% en el original (afectan botones de hipervínculo, el efecto de \pause,
% y los separadores de \subsection). Se preservan como "ajustes locales"
% en los .tex que ya los tenían, justo después de \input{preambulo}.

% =====================================================================
% 9. Macros de Estout (tablas de resultados)
% =====================================================================
\newcommand{\sym}[1]{\rlap{#1}}% Thanks to David Carlisle
\let\estinput=\input% define a new input command so that we can still flatten the document
\newcommand{\estwide}[3]{
	\vspace{.75ex}{
		\begin{tabular*}
			{\textwidth}{@{\hskip\tabcolsep\extracolsep\fill}l*{#2}{#3}}
			\toprule
			\estinput{#1}
			\bottomrule
			\addlinespace[.75ex]
		\end{tabular*}
	}
}
\newcommand{\estauto}[3]{
	\vspace{.75ex}{
		\begin{tabular}{l*{#2}{#3}}
			\toprule
			\estinput{#1}
			\bottomrule
			\addlinespace[.75ex]
		\end{tabular}
	}
}
% Allow line breaks with \\ in specialcells
\newcommand{\specialcell}[2][c]{%
	\begin{tabular}[#1]{@{}c@{}}#2\end{tabular}}

% =====================================================================
% 10. Notas y fuentes al pie de figuras/tablas
% =====================================================================
% Note/Source/Text after Tables
\newcommand{\figtext}[1]{
	\vspace{-1.9ex}
	\captionsetup{justification=justified,font=footnotesize}
	\caption*{\hspace{6pt}\hangindent=1.5em #1}
}
\newcommand{\fignote}[1]{\figtext{\emph{Note:~}~#1}}
\newcommand{\figsource}[1]{\figtext{\emph{Source:~}~#1}}
% Add significance note with \starnote
\newcommand{\starnote}{\figtext{* p $<$ 0.10 ** p $<$ 0.05 *** p $<$ 0.01. }}

% =====================================================================
% 11. Entornos tipo teorema (en español)
% =====================================================================
\newtheorem{Prop}{Proposition}[section]
\newtheorem{Def}{Definition}
%\newtheorem{Assum}{Assumption}
%\newtheorem{Question}{Question}
%\newtheorem{Comment}{Comment}
%\theoremstyle{plain}
%\numberwithin{Assum}{section}
%\numberwithin{figure}{section}
%\numberwithin{equation}{section}
%\numberwithin{Def}{section}

\newenvironment{Teorema}[1][Teorema.]{\begin{trivlist}
		\item[\hskip \labelsep {\bfseries #1}]}{\end{trivlist}}
\newenvironment{Corolario}[1][Corolario.]{\begin{trivlist}
		\item[\hskip \labelsep {\bfseries #1}]}{\end{trivlist}}
\newenvironment{Definicion}[1][Definición.]{\begin{trivlist}
		\item[\hskip \labelsep {\bfseries #1}]}{\end{trivlist}}
\newenvironment{Lema}[1][Lema.]{\begin{trivlist}
		\item[\hskip \labelsep {\bfseries #1}]}{\end{trivlist}}
\newenvironment{Ejemplo}[1][Ejemplo.]{\begin{trivlist}
		\item[\hskip \labelsep {\bfseries #1}]}{\end{trivlist}}

\newenvironment{myindentpar}[1]%
{\begin{list}{}%
		{\setlength{\leftmargin}{#1}}%
		\item[]%
	}
	{\end{list}}

% =====================================================================
% 12. Notación matemática y econométrica
% =====================================================================
\newcommand{\plim}{\operatornamewithlimits{plim}}
\newcommand{\Var}{\operatorname{Var}}
\newcommand{\pderiv}[2]{\frac{\partial #1}{\partial #2}}
\newcommand{\spderiv}[2]{\frac{\partial^2 #1}{\partial #2^2}}
\newcommand{\xderiv}[3]{\frac{\partial^2 #1}{\partial #2 \partial #3}}
\newcommand{\argmin}{\operatornamewithlimits{arg\,min}}
\newcommand{\argmax}{\operatornamewithlimits{arg\,max}}
\newcommand{\amax}{\operatornamewithlimits{max}}
\newcommand{\amin}{\operatornamewithlimits{min}}
\newcommand*\circled[1]{\tikz[baseline=(char.base)]{\node[shape=circle,draw,inner sep=1pt] (char) {#1};}}
\newcommand{\indep}{\perp \!\!\! \perp}
\makeatletter
\newcommand{\distas}[1]{\mathbin{\overset{#1}{\kern\z@\sim}}}%
\newsavebox{\mybox}\newsavebox{\mysim}
\newcommand{\distras}[1]{%
	\savebox{\mybox}{\hbox{\kern3pt$\scriptstyle#1$\kern3pt}}%
	\savebox{\mysim}{\hbox{$\sim$}}%
	\mathbin{\overset{#1}{\kern\z@\resizebox{\wd\mybox}{\ht\mysim}{$\sim$}}}%
}
\makeatother

\DeclareMathOperator{\EX}{\mathbb{E}}% expected value
\newcommand{\Prob}[1]{\text{Pr}\left[#1\right]}
\newcommand{\RK}[1]{\mathbb{R}^{#1}}
\newcommand{\HT}[1]{\mathbb{H}_{#1}}

% Vectores y matrices en negrilla (romana)
\newcommand{\xB}{\textbf{x}}
\newcommand{\yB}{\textbf{y}}
\newcommand{\zB}{\textbf{z}}
\newcommand{\wB}{\textbf{w}}
\newcommand{\eB}{\textbf{e}}
\newcommand{\vB}{\textbf{v}}
\newcommand{\AB}{\textbf{A}}
\newcommand{\XB}{\textbf{X}}
\newcommand{\YB}{\textbf{Y}}
\newcommand{\ZB}{\textbf{Z}}
\newcommand{\QB}{\textbf{Q}}
\newcommand{\WB}{\textbf{W}}

% Vectores griegos en negrilla
\newcommand{\betaB}{\boldsymbol{\beta}}
\newcommand{\gammaB}{\boldsymbol{\gamma}}
\newcommand{\alphaB}{\boldsymbol{\alpha}}
\newcommand{\thetaB}{\boldsymbol{\theta}}
\newcommand{\deltaB}{\boldsymbol{\delta}}
\newcommand{\phiB}{\boldsymbol{\phi}}
\newcommand{\epsilonB}{\boldsymbol{\epsilon}}
\newcommand{\etaB}{\boldsymbol{\eta}}
\newcommand{\betaBH}{\boldsymbol{\hat{\beta}}}
\newcommand{\betaH}{\hat{\beta}}

% Colores de énfasis
\newcommand{\brown}[1]{\textcolor{brown}{#1}}
\newcommand{\blue}[1]{\textcolor{blue}{#1}}
\newcommand{\red}[1]{\textcolor{red}{#1}}

% Espaciados usados en itemize/enumerate
\newcommand{\smallspace}{\vspace{0.1cm}}
\newcommand{\mediumspace}{\vspace{0.3cm}}
\newcommand{\largespace}{\vspace{6cm}}

% Tamaños de fuente auxiliares
\newcommand\fontvia{\fontsize{4}{6}\selectfont}
\newcommand\fontvib{\fontsize{5}{8}\selectfont}
\newcommand\fontvic{\fontsize{7}{10}\selectfont}
\newcommand\fontvid{\fontsize{8}{11}\selectfont}
\newcommand\fontvie{\fontsize{10}{11}\selectfont}

% Celdas de mapa de calor estilizado para la intuición de shift-share
% Monotone (brown) cell: value 0..100 -> brown intensity
\newcommand{\sharecell}[3]{%
  \pgfmathtruncatemacro{\vv}{#3}%
  \fill[brown!\vv!white] (#1,#2) rectangle ({#1+1},{#2+1});%
  \draw[gray!40] (#1,#2) rectangle ({#1+1},{#2+1});%
}
% Diverging (blue/red) cell: value -50..50 -> blue (neg) or red (pos)
\newcommand{\divcell}[3]{%
  \pgfmathtruncatemacro{\absvv}{abs(#3)*2}%
  \pgfmathtruncatemacro{\negflag}{(#3)<0 ? 1 : 0}%
  \ifnum\negflag=1
    \fill[blue!\absvv!white] (#1,#2) rectangle ({#1+1},{#2+1});%
  \else
    \fill[red!\absvv!white] (#1,#2) rectangle ({#1+1},{#2+1});%
  \fi
  \draw[gray!40] (#1,#2) rectangle ({#1+1},{#2+1});%
}
% Uniform cell: same color regardless of (col,row)
\newcommand{\unifcell}[3]{%
  \pgfmathtruncatemacro{\vv}{#3}%
  \fill[red!\vv!white] (#1,#2) rectangle ({#1+1},{#2+1});%
  \draw[gray!40] (#1,#2) rectangle ({#1+1},{#2+1});%
}

% =====================================================================
% 13. Datos de portada (comunes a todas las clases)
% =====================================================================
\title{Econometría Avanzada}
\author{Manuel Fernández}
\date{\today}
%\titlegraphic{\hfill\includegraphics[height=1.5cm]{logo.pdf}}
 DESPUÉS de
% \documentclass[<tam>pt, xcolor=dvipsnames]{beamer} en cada archivo.
%
% Lo único que debe quedar en cada .tex individual es:
%   1) \documentclass[...]{beamer}      (el tamaño de fuente puede variar)
%   2) % =====================================================================
% Preámbulo común — Diapositivas Magistrales, Econometría Avanzada
% =====================================================================
% Se incluye con \input{preambulo/preambulo-magistral} DESPUÉS de
% \documentclass[<tam>pt, xcolor=dvipsnames]{beamer} en cada archivo.
%
% Lo único que debe quedar en cada .tex individual es:
%   1) \documentclass[...]{beamer}      (el tamaño de fuente puede variar)
%   2) \input{preambulo/preambulo-magistral}
%   3) \subtitle{...}                   (título de la clase puntual)
% Todo lo demás (paquetes, macros, tema, título, autor, hypersetup, etc.)
% vive aquí para poder hacer cambios globales en un solo lugar.
% =====================================================================

% ---------------------------------------------------------------------
% 1. Codificación, idioma y fuentes
% ---------------------------------------------------------------------
\usepackage{etex}
\usepackage[utf8]{inputenc}
\usepackage[T1]{fontenc}
\usepackage[spanish,es-nodecimaldot]{babel}
\usepackage{lmodern}
\usepackage{type1cm}
\usepackage{textcomp}

% ---------------------------------------------------------------------
% 2. Matemáticas y símbolos
% ---------------------------------------------------------------------
\usepackage{amsmath,amssymb,amsthm,amsfonts}
\usepackage{bbm}
\usepackage{bm}

% ---------------------------------------------------------------------
% 3. Bibliografía y citas
% ---------------------------------------------------------------------
\usepackage[round]{natbib}
\bibliographystyle{erae}
\usepackage{bibentry}
\usepackage{url}
\let\htmlstyloaded=\relax

% Las citas se muestran en marrón para distinguirlas del contenido.
\let\oldcitet=\citet
\renewcommand{\citet}[1]{\textcolor{brown}{\oldcitet{#1}}}
\let\oldcitep=\citep
\renewcommand{\citep}[1]{\textcolor{brown}{\oldcitep{#1}}}
\let\oldcite=\cite
\renewcommand{\cite}[1]{\textcolor{brown}{\oldcite{#1}}}

% ---------------------------------------------------------------------
% 4. Figuras, tablas y colores
% ---------------------------------------------------------------------
\usepackage{float}
\usepackage{graphicx}
\usepackage{subcaption}
\usepackage[font={footnotesize}]{caption}
\usepackage{lscape}
\usepackage{longtable}
\usepackage{booktabs}
\usepackage{multirow}
\usepackage{array}
\usepackage[flushleft]{threeparttable}
\usepackage{color, colortbl}
\usepackage{eso-pic}
\usepackage{chngcntr}
\usepackage{adjustbox}
\usepackage{pdfpages}
\usepackage{media9}

% siunitx: usado para alinear/formatear números en tablas
\usepackage{siunitx}
\sisetup{
	detect-mode,
	tight-spacing		= true,
	group-digits		= false ,
	input-signs		= ,
	input-symbols		= ( ) [ ] - + *,
	input-open-uncertainty	= ,
	input-close-uncertainty	= ,
	table-align-text-post	= false
}

% ---------------------------------------------------------------------
% 5. TikZ / PGFPlots
% ---------------------------------------------------------------------
\usepackage{tikz}
\usepackage{pgfplots}
\usepgfplotslibrary{dateplot}
\usetikzlibrary{arrows,calc, patterns, positioning, shapes.geometric, decorations.pathreplacing,decorations.markings}

% ---------------------------------------------------------------------
% 6. Misceláneos de formato
% ---------------------------------------------------------------------
\usepackage{setspace}
\usepackage{lipsum}
\usepackage{framed}
\usepackage{mdframed}
\usepackage{fancyhdr}
\usepackage{beamerfoils}
\usepackage{pifont}
\usepackage{ragged2e}
\usepackage{xspace}
\usepackage{appendixnumberbeamer}
\usepackage[scale=2]{ccicons}
\usepackage[most]{tcolorbox}
\usepackage{enumitem}
\setbeamertemplate{note page}[plain]

% ---------------------------------------------------------------------
% 7. Hyperref (debe cargarse después de la mayoría de paquetes)
% ---------------------------------------------------------------------
\usepackage{hyperref}
\hypersetup {
	pdftitle={General},
	pdfauthor={Manuel Fern\'{a}ndez},
	colorlinks=true,
	citecolor=brown,
	linkcolor=black,
	urlcolor=blue
}
\makeatletter
\renewcommand{\hyper@natlinkstart}[1]{}
\renewcommand{\hyper@natlinkend}{}
\renewcommand{\hyper@natlinkbreak}[2]{#1}
\makeatother

% =====================================================================
% 8. Tema Beamer (metropolis)
% =====================================================================
\usetheme{metropolis}
\newcommand{\themename}{\textbf{\textsc{metropolis}}\xspace}
\metroset{sectionpage=none}

\setbeamerfont{citep}{family=\rm}
\setbeamerfont{caption}{size=\scriptsize}
\setbeamerfont{normal text}{size=7pt}

%\setcounter{tocdepth}{1}
\setbeamertemplate{section in toc}[sections numbered]
\AtBeginSection[]{
  \begin{frame}
  \vfill
  \centering
  \begin{beamercolorbox}[sep=8pt,center,shadow=true,rounded=true]{title}
    \usebeamerfont{title}\insertsectionhead\par%
  \end{beamercolorbox}
  \vfill
  \end{frame}
}
% NOTA: \setbeamercolor{button}, \setbeamercovered{invisible} y
% \AtBeginSubsection NO están aquí porque sólo algunas clases los usaban
% en el original (afectan botones de hipervínculo, el efecto de \pause,
% y los separadores de \subsection). Se preservan como "ajustes locales"
% en los .tex que ya los tenían, justo después de \input{preambulo}.

% =====================================================================
% 9. Macros de Estout (tablas de resultados)
% =====================================================================
\newcommand{\sym}[1]{\rlap{#1}}% Thanks to David Carlisle
\let\estinput=\input% define a new input command so that we can still flatten the document
\newcommand{\estwide}[3]{
	\vspace{.75ex}{
		\begin{tabular*}
			{\textwidth}{@{\hskip\tabcolsep\extracolsep\fill}l*{#2}{#3}}
			\toprule
			\estinput{#1}
			\bottomrule
			\addlinespace[.75ex]
		\end{tabular*}
	}
}
\newcommand{\estauto}[3]{
	\vspace{.75ex}{
		\begin{tabular}{l*{#2}{#3}}
			\toprule
			\estinput{#1}
			\bottomrule
			\addlinespace[.75ex]
		\end{tabular}
	}
}
% Allow line breaks with \\ in specialcells
\newcommand{\specialcell}[2][c]{%
	\begin{tabular}[#1]{@{}c@{}}#2\end{tabular}}

% =====================================================================
% 10. Notas y fuentes al pie de figuras/tablas
% =====================================================================
% Note/Source/Text after Tables
\newcommand{\figtext}[1]{
	\vspace{-1.9ex}
	\captionsetup{justification=justified,font=footnotesize}
	\caption*{\hspace{6pt}\hangindent=1.5em #1}
}
\newcommand{\fignote}[1]{\figtext{\emph{Note:~}~#1}}
\newcommand{\figsource}[1]{\figtext{\emph{Source:~}~#1}}
% Add significance note with \starnote
\newcommand{\starnote}{\figtext{* p $<$ 0.10 ** p $<$ 0.05 *** p $<$ 0.01. }}

% =====================================================================
% 11. Entornos tipo teorema (en español)
% =====================================================================
\newtheorem{Prop}{Proposition}[section]
\newtheorem{Def}{Definition}
%\newtheorem{Assum}{Assumption}
%\newtheorem{Question}{Question}
%\newtheorem{Comment}{Comment}
%\theoremstyle{plain}
%\numberwithin{Assum}{section}
%\numberwithin{figure}{section}
%\numberwithin{equation}{section}
%\numberwithin{Def}{section}

\newenvironment{Teorema}[1][Teorema.]{\begin{trivlist}
		\item[\hskip \labelsep {\bfseries #1}]}{\end{trivlist}}
\newenvironment{Corolario}[1][Corolario.]{\begin{trivlist}
		\item[\hskip \labelsep {\bfseries #1}]}{\end{trivlist}}
\newenvironment{Definicion}[1][Definición.]{\begin{trivlist}
		\item[\hskip \labelsep {\bfseries #1}]}{\end{trivlist}}
\newenvironment{Lema}[1][Lema.]{\begin{trivlist}
		\item[\hskip \labelsep {\bfseries #1}]}{\end{trivlist}}
\newenvironment{Ejemplo}[1][Ejemplo.]{\begin{trivlist}
		\item[\hskip \labelsep {\bfseries #1}]}{\end{trivlist}}

\newenvironment{myindentpar}[1]%
{\begin{list}{}%
		{\setlength{\leftmargin}{#1}}%
		\item[]%
	}
	{\end{list}}

% =====================================================================
% 12. Notación matemática y econométrica
% =====================================================================
\newcommand{\plim}{\operatornamewithlimits{plim}}
\newcommand{\Var}{\operatorname{Var}}
\newcommand{\pderiv}[2]{\frac{\partial #1}{\partial #2}}
\newcommand{\spderiv}[2]{\frac{\partial^2 #1}{\partial #2^2}}
\newcommand{\xderiv}[3]{\frac{\partial^2 #1}{\partial #2 \partial #3}}
\newcommand{\argmin}{\operatornamewithlimits{arg\,min}}
\newcommand{\argmax}{\operatornamewithlimits{arg\,max}}
\newcommand{\amax}{\operatornamewithlimits{max}}
\newcommand{\amin}{\operatornamewithlimits{min}}
\newcommand*\circled[1]{\tikz[baseline=(char.base)]{\node[shape=circle,draw,inner sep=1pt] (char) {#1};}}
\newcommand{\indep}{\perp \!\!\! \perp}
\makeatletter
\newcommand{\distas}[1]{\mathbin{\overset{#1}{\kern\z@\sim}}}%
\newsavebox{\mybox}\newsavebox{\mysim}
\newcommand{\distras}[1]{%
	\savebox{\mybox}{\hbox{\kern3pt$\scriptstyle#1$\kern3pt}}%
	\savebox{\mysim}{\hbox{$\sim$}}%
	\mathbin{\overset{#1}{\kern\z@\resizebox{\wd\mybox}{\ht\mysim}{$\sim$}}}%
}
\makeatother

\DeclareMathOperator{\EX}{\mathbb{E}}% expected value
\newcommand{\Prob}[1]{\text{Pr}\left[#1\right]}
\newcommand{\RK}[1]{\mathbb{R}^{#1}}
\newcommand{\HT}[1]{\mathbb{H}_{#1}}

% Vectores y matrices en negrilla (romana)
\newcommand{\xB}{\textbf{x}}
\newcommand{\yB}{\textbf{y}}
\newcommand{\zB}{\textbf{z}}
\newcommand{\wB}{\textbf{w}}
\newcommand{\eB}{\textbf{e}}
\newcommand{\vB}{\textbf{v}}
\newcommand{\AB}{\textbf{A}}
\newcommand{\XB}{\textbf{X}}
\newcommand{\YB}{\textbf{Y}}
\newcommand{\ZB}{\textbf{Z}}
\newcommand{\QB}{\textbf{Q}}
\newcommand{\WB}{\textbf{W}}

% Vectores griegos en negrilla
\newcommand{\betaB}{\boldsymbol{\beta}}
\newcommand{\gammaB}{\boldsymbol{\gamma}}
\newcommand{\alphaB}{\boldsymbol{\alpha}}
\newcommand{\thetaB}{\boldsymbol{\theta}}
\newcommand{\deltaB}{\boldsymbol{\delta}}
\newcommand{\phiB}{\boldsymbol{\phi}}
\newcommand{\epsilonB}{\boldsymbol{\epsilon}}
\newcommand{\etaB}{\boldsymbol{\eta}}
\newcommand{\betaBH}{\boldsymbol{\hat{\beta}}}
\newcommand{\betaH}{\hat{\beta}}

% Colores de énfasis
\newcommand{\brown}[1]{\textcolor{brown}{#1}}
\newcommand{\blue}[1]{\textcolor{blue}{#1}}
\newcommand{\red}[1]{\textcolor{red}{#1}}

% Espaciados usados en itemize/enumerate
\newcommand{\smallspace}{\vspace{0.1cm}}
\newcommand{\mediumspace}{\vspace{0.3cm}}
\newcommand{\largespace}{\vspace{6cm}}

% Tamaños de fuente auxiliares
\newcommand\fontvia{\fontsize{4}{6}\selectfont}
\newcommand\fontvib{\fontsize{5}{8}\selectfont}
\newcommand\fontvic{\fontsize{7}{10}\selectfont}
\newcommand\fontvid{\fontsize{8}{11}\selectfont}
\newcommand\fontvie{\fontsize{10}{11}\selectfont}

% Celdas de mapa de calor estilizado para la intuición de shift-share
% Monotone (brown) cell: value 0..100 -> brown intensity
\newcommand{\sharecell}[3]{%
  \pgfmathtruncatemacro{\vv}{#3}%
  \fill[brown!\vv!white] (#1,#2) rectangle ({#1+1},{#2+1});%
  \draw[gray!40] (#1,#2) rectangle ({#1+1},{#2+1});%
}
% Diverging (blue/red) cell: value -50..50 -> blue (neg) or red (pos)
\newcommand{\divcell}[3]{%
  \pgfmathtruncatemacro{\absvv}{abs(#3)*2}%
  \pgfmathtruncatemacro{\negflag}{(#3)<0 ? 1 : 0}%
  \ifnum\negflag=1
    \fill[blue!\absvv!white] (#1,#2) rectangle ({#1+1},{#2+1});%
  \else
    \fill[red!\absvv!white] (#1,#2) rectangle ({#1+1},{#2+1});%
  \fi
  \draw[gray!40] (#1,#2) rectangle ({#1+1},{#2+1});%
}
% Uniform cell: same color regardless of (col,row)
\newcommand{\unifcell}[3]{%
  \pgfmathtruncatemacro{\vv}{#3}%
  \fill[red!\vv!white] (#1,#2) rectangle ({#1+1},{#2+1});%
  \draw[gray!40] (#1,#2) rectangle ({#1+1},{#2+1});%
}

% =====================================================================
% 13. Datos de portada (comunes a todas las clases)
% =====================================================================
\title{Econometría Avanzada}
\author{Manuel Fernández}
\date{\today}
%\titlegraphic{\hfill\includegraphics[height=1.5cm]{logo.pdf}}

%   3) \subtitle{...}                   (título de la clase puntual)
% Todo lo demás (paquetes, macros, tema, título, autor, hypersetup, etc.)
% vive aquí para poder hacer cambios globales en un solo lugar.
% =====================================================================

% ---------------------------------------------------------------------
% 1. Codificación, idioma y fuentes
% ---------------------------------------------------------------------
\usepackage{etex}
\usepackage[utf8]{inputenc}
\usepackage[T1]{fontenc}
\usepackage[spanish,es-nodecimaldot]{babel}
\usepackage{lmodern}
\usepackage{type1cm}
\usepackage{textcomp}

% ---------------------------------------------------------------------
% 2. Matemáticas y símbolos
% ---------------------------------------------------------------------
\usepackage{amsmath,amssymb,amsthm,amsfonts}
\usepackage{bbm}
\usepackage{bm}

% ---------------------------------------------------------------------
% 3. Bibliografía y citas
% ---------------------------------------------------------------------
\usepackage[round]{natbib}
\bibliographystyle{erae}
\usepackage{bibentry}
\usepackage{url}
\let\htmlstyloaded=\relax

% Las citas se muestran en marrón para distinguirlas del contenido.
\let\oldcitet=\citet
\renewcommand{\citet}[1]{\textcolor{brown}{\oldcitet{#1}}}
\let\oldcitep=\citep
\renewcommand{\citep}[1]{\textcolor{brown}{\oldcitep{#1}}}
\let\oldcite=\cite
\renewcommand{\cite}[1]{\textcolor{brown}{\oldcite{#1}}}

% ---------------------------------------------------------------------
% 4. Figuras, tablas y colores
% ---------------------------------------------------------------------
\usepackage{float}
\usepackage{graphicx}
\usepackage{subcaption}
\usepackage[font={footnotesize}]{caption}
\usepackage{lscape}
\usepackage{longtable}
\usepackage{booktabs}
\usepackage{multirow}
\usepackage{array}
\usepackage[flushleft]{threeparttable}
\usepackage{color, colortbl}
\usepackage{eso-pic}
\usepackage{chngcntr}
\usepackage{adjustbox}
\usepackage{pdfpages}
\usepackage{media9}

% siunitx: usado para alinear/formatear números en tablas
\usepackage{siunitx}
\sisetup{
	detect-mode,
	tight-spacing		= true,
	group-digits		= false ,
	input-signs		= ,
	input-symbols		= ( ) [ ] - + *,
	input-open-uncertainty	= ,
	input-close-uncertainty	= ,
	table-align-text-post	= false
}

% ---------------------------------------------------------------------
% 5. TikZ / PGFPlots
% ---------------------------------------------------------------------
\usepackage{tikz}
\usepackage{pgfplots}
\usepgfplotslibrary{dateplot}
\usetikzlibrary{arrows,calc, patterns, positioning, shapes.geometric, decorations.pathreplacing,decorations.markings}

% ---------------------------------------------------------------------
% 6. Misceláneos de formato
% ---------------------------------------------------------------------
\usepackage{setspace}
\usepackage{lipsum}
\usepackage{framed}
\usepackage{mdframed}
\usepackage{fancyhdr}
\usepackage{beamerfoils}
\usepackage{pifont}
\usepackage{ragged2e}
\usepackage{xspace}
\usepackage{appendixnumberbeamer}
\usepackage[scale=2]{ccicons}
\usepackage[most]{tcolorbox}
\usepackage{enumitem}
\setbeamertemplate{note page}[plain]

% ---------------------------------------------------------------------
% 7. Hyperref (debe cargarse después de la mayoría de paquetes)
% ---------------------------------------------------------------------
\usepackage{hyperref}
\hypersetup {
	pdftitle={General},
	pdfauthor={Manuel Fern\'{a}ndez},
	colorlinks=true,
	citecolor=brown,
	linkcolor=black,
	urlcolor=blue
}
\makeatletter
\renewcommand{\hyper@natlinkstart}[1]{}
\renewcommand{\hyper@natlinkend}{}
\renewcommand{\hyper@natlinkbreak}[2]{#1}
\makeatother

% =====================================================================
% 8. Tema Beamer (metropolis)
% =====================================================================
\usetheme{metropolis}
\newcommand{\themename}{\textbf{\textsc{metropolis}}\xspace}
\metroset{sectionpage=none}

\setbeamerfont{citep}{family=\rm}
\setbeamerfont{caption}{size=\scriptsize}
\setbeamerfont{normal text}{size=7pt}

%\setcounter{tocdepth}{1}
\setbeamertemplate{section in toc}[sections numbered]
\AtBeginSection[]{
  \begin{frame}
  \vfill
  \centering
  \begin{beamercolorbox}[sep=8pt,center,shadow=true,rounded=true]{title}
    \usebeamerfont{title}\insertsectionhead\par%
  \end{beamercolorbox}
  \vfill
  \end{frame}
}
% NOTA: \setbeamercolor{button}, \setbeamercovered{invisible} y
% \AtBeginSubsection NO están aquí porque sólo algunas clases los usaban
% en el original (afectan botones de hipervínculo, el efecto de \pause,
% y los separadores de \subsection). Se preservan como "ajustes locales"
% en los .tex que ya los tenían, justo después de \input{preambulo}.

% =====================================================================
% 9. Macros de Estout (tablas de resultados)
% =====================================================================
\newcommand{\sym}[1]{\rlap{#1}}% Thanks to David Carlisle
\let\estinput=\input% define a new input command so that we can still flatten the document
\newcommand{\estwide}[3]{
	\vspace{.75ex}{
		\begin{tabular*}
			{\textwidth}{@{\hskip\tabcolsep\extracolsep\fill}l*{#2}{#3}}
			\toprule
			\estinput{#1}
			\bottomrule
			\addlinespace[.75ex]
		\end{tabular*}
	}
}
\newcommand{\estauto}[3]{
	\vspace{.75ex}{
		\begin{tabular}{l*{#2}{#3}}
			\toprule
			\estinput{#1}
			\bottomrule
			\addlinespace[.75ex]
		\end{tabular}
	}
}
% Allow line breaks with \\ in specialcells
\newcommand{\specialcell}[2][c]{%
	\begin{tabular}[#1]{@{}c@{}}#2\end{tabular}}

% =====================================================================
% 10. Notas y fuentes al pie de figuras/tablas
% =====================================================================
% Note/Source/Text after Tables
\newcommand{\figtext}[1]{
	\vspace{-1.9ex}
	\captionsetup{justification=justified,font=footnotesize}
	\caption*{\hspace{6pt}\hangindent=1.5em #1}
}
\newcommand{\fignote}[1]{\figtext{\emph{Note:~}~#1}}
\newcommand{\figsource}[1]{\figtext{\emph{Source:~}~#1}}
% Add significance note with \starnote
\newcommand{\starnote}{\figtext{* p $<$ 0.10 ** p $<$ 0.05 *** p $<$ 0.01. }}

% =====================================================================
% 11. Entornos tipo teorema (en español)
% =====================================================================
\newtheorem{Prop}{Proposition}[section]
\newtheorem{Def}{Definition}
%\newtheorem{Assum}{Assumption}
%\newtheorem{Question}{Question}
%\newtheorem{Comment}{Comment}
%\theoremstyle{plain}
%\numberwithin{Assum}{section}
%\numberwithin{figure}{section}
%\numberwithin{equation}{section}
%\numberwithin{Def}{section}

\newenvironment{Teorema}[1][Teorema.]{\begin{trivlist}
		\item[\hskip \labelsep {\bfseries #1}]}{\end{trivlist}}
\newenvironment{Corolario}[1][Corolario.]{\begin{trivlist}
		\item[\hskip \labelsep {\bfseries #1}]}{\end{trivlist}}
\newenvironment{Definicion}[1][Definición.]{\begin{trivlist}
		\item[\hskip \labelsep {\bfseries #1}]}{\end{trivlist}}
\newenvironment{Lema}[1][Lema.]{\begin{trivlist}
		\item[\hskip \labelsep {\bfseries #1}]}{\end{trivlist}}
\newenvironment{Ejemplo}[1][Ejemplo.]{\begin{trivlist}
		\item[\hskip \labelsep {\bfseries #1}]}{\end{trivlist}}

\newenvironment{myindentpar}[1]%
{\begin{list}{}%
		{\setlength{\leftmargin}{#1}}%
		\item[]%
	}
	{\end{list}}

% =====================================================================
% 12. Notación matemática y econométrica
% =====================================================================
\newcommand{\plim}{\operatornamewithlimits{plim}}
\newcommand{\Var}{\operatorname{Var}}
\newcommand{\pderiv}[2]{\frac{\partial #1}{\partial #2}}
\newcommand{\spderiv}[2]{\frac{\partial^2 #1}{\partial #2^2}}
\newcommand{\xderiv}[3]{\frac{\partial^2 #1}{\partial #2 \partial #3}}
\newcommand{\argmin}{\operatornamewithlimits{arg\,min}}
\newcommand{\argmax}{\operatornamewithlimits{arg\,max}}
\newcommand{\amax}{\operatornamewithlimits{max}}
\newcommand{\amin}{\operatornamewithlimits{min}}
\newcommand*\circled[1]{\tikz[baseline=(char.base)]{\node[shape=circle,draw,inner sep=1pt] (char) {#1};}}
\newcommand{\indep}{\perp \!\!\! \perp}
\makeatletter
\newcommand{\distas}[1]{\mathbin{\overset{#1}{\kern\z@\sim}}}%
\newsavebox{\mybox}\newsavebox{\mysim}
\newcommand{\distras}[1]{%
	\savebox{\mybox}{\hbox{\kern3pt$\scriptstyle#1$\kern3pt}}%
	\savebox{\mysim}{\hbox{$\sim$}}%
	\mathbin{\overset{#1}{\kern\z@\resizebox{\wd\mybox}{\ht\mysim}{$\sim$}}}%
}
\makeatother

\DeclareMathOperator{\EX}{\mathbb{E}}% expected value
\newcommand{\Prob}[1]{\text{Pr}\left[#1\right]}
\newcommand{\RK}[1]{\mathbb{R}^{#1}}
\newcommand{\HT}[1]{\mathbb{H}_{#1}}

% Vectores y matrices en negrilla (romana)
\newcommand{\xB}{\textbf{x}}
\newcommand{\yB}{\textbf{y}}
\newcommand{\zB}{\textbf{z}}
\newcommand{\wB}{\textbf{w}}
\newcommand{\eB}{\textbf{e}}
\newcommand{\vB}{\textbf{v}}
\newcommand{\AB}{\textbf{A}}
\newcommand{\XB}{\textbf{X}}
\newcommand{\YB}{\textbf{Y}}
\newcommand{\ZB}{\textbf{Z}}
\newcommand{\QB}{\textbf{Q}}
\newcommand{\WB}{\textbf{W}}

% Vectores griegos en negrilla
\newcommand{\betaB}{\boldsymbol{\beta}}
\newcommand{\gammaB}{\boldsymbol{\gamma}}
\newcommand{\alphaB}{\boldsymbol{\alpha}}
\newcommand{\thetaB}{\boldsymbol{\theta}}
\newcommand{\deltaB}{\boldsymbol{\delta}}
\newcommand{\phiB}{\boldsymbol{\phi}}
\newcommand{\epsilonB}{\boldsymbol{\epsilon}}
\newcommand{\etaB}{\boldsymbol{\eta}}
\newcommand{\betaBH}{\boldsymbol{\hat{\beta}}}
\newcommand{\betaH}{\hat{\beta}}

% Colores de énfasis
\newcommand{\brown}[1]{\textcolor{brown}{#1}}
\newcommand{\blue}[1]{\textcolor{blue}{#1}}
\newcommand{\red}[1]{\textcolor{red}{#1}}

% Espaciados usados en itemize/enumerate
\newcommand{\smallspace}{\vspace{0.1cm}}
\newcommand{\mediumspace}{\vspace{0.3cm}}
\newcommand{\largespace}{\vspace{6cm}}

% Tamaños de fuente auxiliares
\newcommand\fontvia{\fontsize{4}{6}\selectfont}
\newcommand\fontvib{\fontsize{5}{8}\selectfont}
\newcommand\fontvic{\fontsize{7}{10}\selectfont}
\newcommand\fontvid{\fontsize{8}{11}\selectfont}
\newcommand\fontvie{\fontsize{10}{11}\selectfont}

% Celdas de mapa de calor estilizado para la intuición de shift-share
% Monotone (brown) cell: value 0..100 -> brown intensity
\newcommand{\sharecell}[3]{%
  \pgfmathtruncatemacro{\vv}{#3}%
  \fill[brown!\vv!white] (#1,#2) rectangle ({#1+1},{#2+1});%
  \draw[gray!40] (#1,#2) rectangle ({#1+1},{#2+1});%
}
% Diverging (blue/red) cell: value -50..50 -> blue (neg) or red (pos)
\newcommand{\divcell}[3]{%
  \pgfmathtruncatemacro{\absvv}{abs(#3)*2}%
  \pgfmathtruncatemacro{\negflag}{(#3)<0 ? 1 : 0}%
  \ifnum\negflag=1
    \fill[blue!\absvv!white] (#1,#2) rectangle ({#1+1},{#2+1});%
  \else
    \fill[red!\absvv!white] (#1,#2) rectangle ({#1+1},{#2+1});%
  \fi
  \draw[gray!40] (#1,#2) rectangle ({#1+1},{#2+1});%
}
% Uniform cell: same color regardless of (col,row)
\newcommand{\unifcell}[3]{%
  \pgfmathtruncatemacro{\vv}{#3}%
  \fill[red!\vv!white] (#1,#2) rectangle ({#1+1},{#2+1});%
  \draw[gray!40] (#1,#2) rectangle ({#1+1},{#2+1});%
}

% =====================================================================
% 13. Datos de portada (comunes a todas las clases)
% =====================================================================
\title{Econometría Avanzada}
\author{Manuel Fernández}
\date{\today}
%\titlegraphic{\hfill\includegraphics[height=1.5cm]{logo.pdf}}
 DESPUÉS de
% \documentclass[<tam>pt, xcolor=dvipsnames]{beamer} en cada archivo.
%
% Lo único que debe quedar en cada .tex individual es:
%   1) \documentclass[...]{beamer}      (el tamaño de fuente puede variar)
%   2) % =====================================================================
% Preámbulo común — Diapositivas Magistrales, Econometría Avanzada
% =====================================================================
% Se incluye con % =====================================================================
% Preámbulo común — Diapositivas Magistrales, Econometría Avanzada
% =====================================================================
% Se incluye con \input{preambulo/preambulo-magistral} DESPUÉS de
% \documentclass[<tam>pt, xcolor=dvipsnames]{beamer} en cada archivo.
%
% Lo único que debe quedar en cada .tex individual es:
%   1) \documentclass[...]{beamer}      (el tamaño de fuente puede variar)
%   2) \input{preambulo/preambulo-magistral}
%   3) \subtitle{...}                   (título de la clase puntual)
% Todo lo demás (paquetes, macros, tema, título, autor, hypersetup, etc.)
% vive aquí para poder hacer cambios globales en un solo lugar.
% =====================================================================

% ---------------------------------------------------------------------
% 1. Codificación, idioma y fuentes
% ---------------------------------------------------------------------
\usepackage{etex}
\usepackage[utf8]{inputenc}
\usepackage[T1]{fontenc}
\usepackage[spanish,es-nodecimaldot]{babel}
\usepackage{lmodern}
\usepackage{type1cm}
\usepackage{textcomp}

% ---------------------------------------------------------------------
% 2. Matemáticas y símbolos
% ---------------------------------------------------------------------
\usepackage{amsmath,amssymb,amsthm,amsfonts}
\usepackage{bbm}
\usepackage{bm}

% ---------------------------------------------------------------------
% 3. Bibliografía y citas
% ---------------------------------------------------------------------
\usepackage[round]{natbib}
\bibliographystyle{erae}
\usepackage{bibentry}
\usepackage{url}
\let\htmlstyloaded=\relax

% Las citas se muestran en marrón para distinguirlas del contenido.
\let\oldcitet=\citet
\renewcommand{\citet}[1]{\textcolor{brown}{\oldcitet{#1}}}
\let\oldcitep=\citep
\renewcommand{\citep}[1]{\textcolor{brown}{\oldcitep{#1}}}
\let\oldcite=\cite
\renewcommand{\cite}[1]{\textcolor{brown}{\oldcite{#1}}}

% ---------------------------------------------------------------------
% 4. Figuras, tablas y colores
% ---------------------------------------------------------------------
\usepackage{float}
\usepackage{graphicx}
\usepackage{subcaption}
\usepackage[font={footnotesize}]{caption}
\usepackage{lscape}
\usepackage{longtable}
\usepackage{booktabs}
\usepackage{multirow}
\usepackage{array}
\usepackage[flushleft]{threeparttable}
\usepackage{color, colortbl}
\usepackage{eso-pic}
\usepackage{chngcntr}
\usepackage{adjustbox}
\usepackage{pdfpages}
\usepackage{media9}

% siunitx: usado para alinear/formatear números en tablas
\usepackage{siunitx}
\sisetup{
	detect-mode,
	tight-spacing		= true,
	group-digits		= false ,
	input-signs		= ,
	input-symbols		= ( ) [ ] - + *,
	input-open-uncertainty	= ,
	input-close-uncertainty	= ,
	table-align-text-post	= false
}

% ---------------------------------------------------------------------
% 5. TikZ / PGFPlots
% ---------------------------------------------------------------------
\usepackage{tikz}
\usepackage{pgfplots}
\usepgfplotslibrary{dateplot}
\usetikzlibrary{arrows,calc, patterns, positioning, shapes.geometric, decorations.pathreplacing,decorations.markings}

% ---------------------------------------------------------------------
% 6. Misceláneos de formato
% ---------------------------------------------------------------------
\usepackage{setspace}
\usepackage{lipsum}
\usepackage{framed}
\usepackage{mdframed}
\usepackage{fancyhdr}
\usepackage{beamerfoils}
\usepackage{pifont}
\usepackage{ragged2e}
\usepackage{xspace}
\usepackage{appendixnumberbeamer}
\usepackage[scale=2]{ccicons}
\usepackage[most]{tcolorbox}
\usepackage{enumitem}
\setbeamertemplate{note page}[plain]

% ---------------------------------------------------------------------
% 7. Hyperref (debe cargarse después de la mayoría de paquetes)
% ---------------------------------------------------------------------
\usepackage{hyperref}
\hypersetup {
	pdftitle={General},
	pdfauthor={Manuel Fern\'{a}ndez},
	colorlinks=true,
	citecolor=brown,
	linkcolor=black,
	urlcolor=blue
}
\makeatletter
\renewcommand{\hyper@natlinkstart}[1]{}
\renewcommand{\hyper@natlinkend}{}
\renewcommand{\hyper@natlinkbreak}[2]{#1}
\makeatother

% =====================================================================
% 8. Tema Beamer (metropolis)
% =====================================================================
\usetheme{metropolis}
\newcommand{\themename}{\textbf{\textsc{metropolis}}\xspace}
\metroset{sectionpage=none}

\setbeamerfont{citep}{family=\rm}
\setbeamerfont{caption}{size=\scriptsize}
\setbeamerfont{normal text}{size=7pt}

%\setcounter{tocdepth}{1}
\setbeamertemplate{section in toc}[sections numbered]
\AtBeginSection[]{
  \begin{frame}
  \vfill
  \centering
  \begin{beamercolorbox}[sep=8pt,center,shadow=true,rounded=true]{title}
    \usebeamerfont{title}\insertsectionhead\par%
  \end{beamercolorbox}
  \vfill
  \end{frame}
}
% NOTA: \setbeamercolor{button}, \setbeamercovered{invisible} y
% \AtBeginSubsection NO están aquí porque sólo algunas clases los usaban
% en el original (afectan botones de hipervínculo, el efecto de \pause,
% y los separadores de \subsection). Se preservan como "ajustes locales"
% en los .tex que ya los tenían, justo después de \input{preambulo}.

% =====================================================================
% 9. Macros de Estout (tablas de resultados)
% =====================================================================
\newcommand{\sym}[1]{\rlap{#1}}% Thanks to David Carlisle
\let\estinput=\input% define a new input command so that we can still flatten the document
\newcommand{\estwide}[3]{
	\vspace{.75ex}{
		\begin{tabular*}
			{\textwidth}{@{\hskip\tabcolsep\extracolsep\fill}l*{#2}{#3}}
			\toprule
			\estinput{#1}
			\bottomrule
			\addlinespace[.75ex]
		\end{tabular*}
	}
}
\newcommand{\estauto}[3]{
	\vspace{.75ex}{
		\begin{tabular}{l*{#2}{#3}}
			\toprule
			\estinput{#1}
			\bottomrule
			\addlinespace[.75ex]
		\end{tabular}
	}
}
% Allow line breaks with \\ in specialcells
\newcommand{\specialcell}[2][c]{%
	\begin{tabular}[#1]{@{}c@{}}#2\end{tabular}}

% =====================================================================
% 10. Notas y fuentes al pie de figuras/tablas
% =====================================================================
% Note/Source/Text after Tables
\newcommand{\figtext}[1]{
	\vspace{-1.9ex}
	\captionsetup{justification=justified,font=footnotesize}
	\caption*{\hspace{6pt}\hangindent=1.5em #1}
}
\newcommand{\fignote}[1]{\figtext{\emph{Note:~}~#1}}
\newcommand{\figsource}[1]{\figtext{\emph{Source:~}~#1}}
% Add significance note with \starnote
\newcommand{\starnote}{\figtext{* p $<$ 0.10 ** p $<$ 0.05 *** p $<$ 0.01. }}

% =====================================================================
% 11. Entornos tipo teorema (en español)
% =====================================================================
\newtheorem{Prop}{Proposition}[section]
\newtheorem{Def}{Definition}
%\newtheorem{Assum}{Assumption}
%\newtheorem{Question}{Question}
%\newtheorem{Comment}{Comment}
%\theoremstyle{plain}
%\numberwithin{Assum}{section}
%\numberwithin{figure}{section}
%\numberwithin{equation}{section}
%\numberwithin{Def}{section}

\newenvironment{Teorema}[1][Teorema.]{\begin{trivlist}
		\item[\hskip \labelsep {\bfseries #1}]}{\end{trivlist}}
\newenvironment{Corolario}[1][Corolario.]{\begin{trivlist}
		\item[\hskip \labelsep {\bfseries #1}]}{\end{trivlist}}
\newenvironment{Definicion}[1][Definición.]{\begin{trivlist}
		\item[\hskip \labelsep {\bfseries #1}]}{\end{trivlist}}
\newenvironment{Lema}[1][Lema.]{\begin{trivlist}
		\item[\hskip \labelsep {\bfseries #1}]}{\end{trivlist}}
\newenvironment{Ejemplo}[1][Ejemplo.]{\begin{trivlist}
		\item[\hskip \labelsep {\bfseries #1}]}{\end{trivlist}}

\newenvironment{myindentpar}[1]%
{\begin{list}{}%
		{\setlength{\leftmargin}{#1}}%
		\item[]%
	}
	{\end{list}}

% =====================================================================
% 12. Notación matemática y econométrica
% =====================================================================
\newcommand{\plim}{\operatornamewithlimits{plim}}
\newcommand{\Var}{\operatorname{Var}}
\newcommand{\pderiv}[2]{\frac{\partial #1}{\partial #2}}
\newcommand{\spderiv}[2]{\frac{\partial^2 #1}{\partial #2^2}}
\newcommand{\xderiv}[3]{\frac{\partial^2 #1}{\partial #2 \partial #3}}
\newcommand{\argmin}{\operatornamewithlimits{arg\,min}}
\newcommand{\argmax}{\operatornamewithlimits{arg\,max}}
\newcommand{\amax}{\operatornamewithlimits{max}}
\newcommand{\amin}{\operatornamewithlimits{min}}
\newcommand*\circled[1]{\tikz[baseline=(char.base)]{\node[shape=circle,draw,inner sep=1pt] (char) {#1};}}
\newcommand{\indep}{\perp \!\!\! \perp}
\makeatletter
\newcommand{\distas}[1]{\mathbin{\overset{#1}{\kern\z@\sim}}}%
\newsavebox{\mybox}\newsavebox{\mysim}
\newcommand{\distras}[1]{%
	\savebox{\mybox}{\hbox{\kern3pt$\scriptstyle#1$\kern3pt}}%
	\savebox{\mysim}{\hbox{$\sim$}}%
	\mathbin{\overset{#1}{\kern\z@\resizebox{\wd\mybox}{\ht\mysim}{$\sim$}}}%
}
\makeatother

\DeclareMathOperator{\EX}{\mathbb{E}}% expected value
\newcommand{\Prob}[1]{\text{Pr}\left[#1\right]}
\newcommand{\RK}[1]{\mathbb{R}^{#1}}
\newcommand{\HT}[1]{\mathbb{H}_{#1}}

% Vectores y matrices en negrilla (romana)
\newcommand{\xB}{\textbf{x}}
\newcommand{\yB}{\textbf{y}}
\newcommand{\zB}{\textbf{z}}
\newcommand{\wB}{\textbf{w}}
\newcommand{\eB}{\textbf{e}}
\newcommand{\vB}{\textbf{v}}
\newcommand{\AB}{\textbf{A}}
\newcommand{\XB}{\textbf{X}}
\newcommand{\YB}{\textbf{Y}}
\newcommand{\ZB}{\textbf{Z}}
\newcommand{\QB}{\textbf{Q}}
\newcommand{\WB}{\textbf{W}}

% Vectores griegos en negrilla
\newcommand{\betaB}{\boldsymbol{\beta}}
\newcommand{\gammaB}{\boldsymbol{\gamma}}
\newcommand{\alphaB}{\boldsymbol{\alpha}}
\newcommand{\thetaB}{\boldsymbol{\theta}}
\newcommand{\deltaB}{\boldsymbol{\delta}}
\newcommand{\phiB}{\boldsymbol{\phi}}
\newcommand{\epsilonB}{\boldsymbol{\epsilon}}
\newcommand{\etaB}{\boldsymbol{\eta}}
\newcommand{\betaBH}{\boldsymbol{\hat{\beta}}}
\newcommand{\betaH}{\hat{\beta}}

% Colores de énfasis
\newcommand{\brown}[1]{\textcolor{brown}{#1}}
\newcommand{\blue}[1]{\textcolor{blue}{#1}}
\newcommand{\red}[1]{\textcolor{red}{#1}}

% Espaciados usados en itemize/enumerate
\newcommand{\smallspace}{\vspace{0.1cm}}
\newcommand{\mediumspace}{\vspace{0.3cm}}
\newcommand{\largespace}{\vspace{6cm}}

% Tamaños de fuente auxiliares
\newcommand\fontvia{\fontsize{4}{6}\selectfont}
\newcommand\fontvib{\fontsize{5}{8}\selectfont}
\newcommand\fontvic{\fontsize{7}{10}\selectfont}
\newcommand\fontvid{\fontsize{8}{11}\selectfont}
\newcommand\fontvie{\fontsize{10}{11}\selectfont}

% Celdas de mapa de calor estilizado para la intuición de shift-share
% Monotone (brown) cell: value 0..100 -> brown intensity
\newcommand{\sharecell}[3]{%
  \pgfmathtruncatemacro{\vv}{#3}%
  \fill[brown!\vv!white] (#1,#2) rectangle ({#1+1},{#2+1});%
  \draw[gray!40] (#1,#2) rectangle ({#1+1},{#2+1});%
}
% Diverging (blue/red) cell: value -50..50 -> blue (neg) or red (pos)
\newcommand{\divcell}[3]{%
  \pgfmathtruncatemacro{\absvv}{abs(#3)*2}%
  \pgfmathtruncatemacro{\negflag}{(#3)<0 ? 1 : 0}%
  \ifnum\negflag=1
    \fill[blue!\absvv!white] (#1,#2) rectangle ({#1+1},{#2+1});%
  \else
    \fill[red!\absvv!white] (#1,#2) rectangle ({#1+1},{#2+1});%
  \fi
  \draw[gray!40] (#1,#2) rectangle ({#1+1},{#2+1});%
}
% Uniform cell: same color regardless of (col,row)
\newcommand{\unifcell}[3]{%
  \pgfmathtruncatemacro{\vv}{#3}%
  \fill[red!\vv!white] (#1,#2) rectangle ({#1+1},{#2+1});%
  \draw[gray!40] (#1,#2) rectangle ({#1+1},{#2+1});%
}

% =====================================================================
% 13. Datos de portada (comunes a todas las clases)
% =====================================================================
\title{Econometría Avanzada}
\author{Manuel Fernández}
\date{\today}
%\titlegraphic{\hfill\includegraphics[height=1.5cm]{logo.pdf}}
 DESPUÉS de
% \documentclass[<tam>pt, xcolor=dvipsnames]{beamer} en cada archivo.
%
% Lo único que debe quedar en cada .tex individual es:
%   1) \documentclass[...]{beamer}      (el tamaño de fuente puede variar)
%   2) % =====================================================================
% Preámbulo común — Diapositivas Magistrales, Econometría Avanzada
% =====================================================================
% Se incluye con \input{preambulo/preambulo-magistral} DESPUÉS de
% \documentclass[<tam>pt, xcolor=dvipsnames]{beamer} en cada archivo.
%
% Lo único que debe quedar en cada .tex individual es:
%   1) \documentclass[...]{beamer}      (el tamaño de fuente puede variar)
%   2) \input{preambulo/preambulo-magistral}
%   3) \subtitle{...}                   (título de la clase puntual)
% Todo lo demás (paquetes, macros, tema, título, autor, hypersetup, etc.)
% vive aquí para poder hacer cambios globales en un solo lugar.
% =====================================================================

% ---------------------------------------------------------------------
% 1. Codificación, idioma y fuentes
% ---------------------------------------------------------------------
\usepackage{etex}
\usepackage[utf8]{inputenc}
\usepackage[T1]{fontenc}
\usepackage[spanish,es-nodecimaldot]{babel}
\usepackage{lmodern}
\usepackage{type1cm}
\usepackage{textcomp}

% ---------------------------------------------------------------------
% 2. Matemáticas y símbolos
% ---------------------------------------------------------------------
\usepackage{amsmath,amssymb,amsthm,amsfonts}
\usepackage{bbm}
\usepackage{bm}

% ---------------------------------------------------------------------
% 3. Bibliografía y citas
% ---------------------------------------------------------------------
\usepackage[round]{natbib}
\bibliographystyle{erae}
\usepackage{bibentry}
\usepackage{url}
\let\htmlstyloaded=\relax

% Las citas se muestran en marrón para distinguirlas del contenido.
\let\oldcitet=\citet
\renewcommand{\citet}[1]{\textcolor{brown}{\oldcitet{#1}}}
\let\oldcitep=\citep
\renewcommand{\citep}[1]{\textcolor{brown}{\oldcitep{#1}}}
\let\oldcite=\cite
\renewcommand{\cite}[1]{\textcolor{brown}{\oldcite{#1}}}

% ---------------------------------------------------------------------
% 4. Figuras, tablas y colores
% ---------------------------------------------------------------------
\usepackage{float}
\usepackage{graphicx}
\usepackage{subcaption}
\usepackage[font={footnotesize}]{caption}
\usepackage{lscape}
\usepackage{longtable}
\usepackage{booktabs}
\usepackage{multirow}
\usepackage{array}
\usepackage[flushleft]{threeparttable}
\usepackage{color, colortbl}
\usepackage{eso-pic}
\usepackage{chngcntr}
\usepackage{adjustbox}
\usepackage{pdfpages}
\usepackage{media9}

% siunitx: usado para alinear/formatear números en tablas
\usepackage{siunitx}
\sisetup{
	detect-mode,
	tight-spacing		= true,
	group-digits		= false ,
	input-signs		= ,
	input-symbols		= ( ) [ ] - + *,
	input-open-uncertainty	= ,
	input-close-uncertainty	= ,
	table-align-text-post	= false
}

% ---------------------------------------------------------------------
% 5. TikZ / PGFPlots
% ---------------------------------------------------------------------
\usepackage{tikz}
\usepackage{pgfplots}
\usepgfplotslibrary{dateplot}
\usetikzlibrary{arrows,calc, patterns, positioning, shapes.geometric, decorations.pathreplacing,decorations.markings}

% ---------------------------------------------------------------------
% 6. Misceláneos de formato
% ---------------------------------------------------------------------
\usepackage{setspace}
\usepackage{lipsum}
\usepackage{framed}
\usepackage{mdframed}
\usepackage{fancyhdr}
\usepackage{beamerfoils}
\usepackage{pifont}
\usepackage{ragged2e}
\usepackage{xspace}
\usepackage{appendixnumberbeamer}
\usepackage[scale=2]{ccicons}
\usepackage[most]{tcolorbox}
\usepackage{enumitem}
\setbeamertemplate{note page}[plain]

% ---------------------------------------------------------------------
% 7. Hyperref (debe cargarse después de la mayoría de paquetes)
% ---------------------------------------------------------------------
\usepackage{hyperref}
\hypersetup {
	pdftitle={General},
	pdfauthor={Manuel Fern\'{a}ndez},
	colorlinks=true,
	citecolor=brown,
	linkcolor=black,
	urlcolor=blue
}
\makeatletter
\renewcommand{\hyper@natlinkstart}[1]{}
\renewcommand{\hyper@natlinkend}{}
\renewcommand{\hyper@natlinkbreak}[2]{#1}
\makeatother

% =====================================================================
% 8. Tema Beamer (metropolis)
% =====================================================================
\usetheme{metropolis}
\newcommand{\themename}{\textbf{\textsc{metropolis}}\xspace}
\metroset{sectionpage=none}

\setbeamerfont{citep}{family=\rm}
\setbeamerfont{caption}{size=\scriptsize}
\setbeamerfont{normal text}{size=7pt}

%\setcounter{tocdepth}{1}
\setbeamertemplate{section in toc}[sections numbered]
\AtBeginSection[]{
  \begin{frame}
  \vfill
  \centering
  \begin{beamercolorbox}[sep=8pt,center,shadow=true,rounded=true]{title}
    \usebeamerfont{title}\insertsectionhead\par%
  \end{beamercolorbox}
  \vfill
  \end{frame}
}
% NOTA: \setbeamercolor{button}, \setbeamercovered{invisible} y
% \AtBeginSubsection NO están aquí porque sólo algunas clases los usaban
% en el original (afectan botones de hipervínculo, el efecto de \pause,
% y los separadores de \subsection). Se preservan como "ajustes locales"
% en los .tex que ya los tenían, justo después de \input{preambulo}.

% =====================================================================
% 9. Macros de Estout (tablas de resultados)
% =====================================================================
\newcommand{\sym}[1]{\rlap{#1}}% Thanks to David Carlisle
\let\estinput=\input% define a new input command so that we can still flatten the document
\newcommand{\estwide}[3]{
	\vspace{.75ex}{
		\begin{tabular*}
			{\textwidth}{@{\hskip\tabcolsep\extracolsep\fill}l*{#2}{#3}}
			\toprule
			\estinput{#1}
			\bottomrule
			\addlinespace[.75ex]
		\end{tabular*}
	}
}
\newcommand{\estauto}[3]{
	\vspace{.75ex}{
		\begin{tabular}{l*{#2}{#3}}
			\toprule
			\estinput{#1}
			\bottomrule
			\addlinespace[.75ex]
		\end{tabular}
	}
}
% Allow line breaks with \\ in specialcells
\newcommand{\specialcell}[2][c]{%
	\begin{tabular}[#1]{@{}c@{}}#2\end{tabular}}

% =====================================================================
% 10. Notas y fuentes al pie de figuras/tablas
% =====================================================================
% Note/Source/Text after Tables
\newcommand{\figtext}[1]{
	\vspace{-1.9ex}
	\captionsetup{justification=justified,font=footnotesize}
	\caption*{\hspace{6pt}\hangindent=1.5em #1}
}
\newcommand{\fignote}[1]{\figtext{\emph{Note:~}~#1}}
\newcommand{\figsource}[1]{\figtext{\emph{Source:~}~#1}}
% Add significance note with \starnote
\newcommand{\starnote}{\figtext{* p $<$ 0.10 ** p $<$ 0.05 *** p $<$ 0.01. }}

% =====================================================================
% 11. Entornos tipo teorema (en español)
% =====================================================================
\newtheorem{Prop}{Proposition}[section]
\newtheorem{Def}{Definition}
%\newtheorem{Assum}{Assumption}
%\newtheorem{Question}{Question}
%\newtheorem{Comment}{Comment}
%\theoremstyle{plain}
%\numberwithin{Assum}{section}
%\numberwithin{figure}{section}
%\numberwithin{equation}{section}
%\numberwithin{Def}{section}

\newenvironment{Teorema}[1][Teorema.]{\begin{trivlist}
		\item[\hskip \labelsep {\bfseries #1}]}{\end{trivlist}}
\newenvironment{Corolario}[1][Corolario.]{\begin{trivlist}
		\item[\hskip \labelsep {\bfseries #1}]}{\end{trivlist}}
\newenvironment{Definicion}[1][Definición.]{\begin{trivlist}
		\item[\hskip \labelsep {\bfseries #1}]}{\end{trivlist}}
\newenvironment{Lema}[1][Lema.]{\begin{trivlist}
		\item[\hskip \labelsep {\bfseries #1}]}{\end{trivlist}}
\newenvironment{Ejemplo}[1][Ejemplo.]{\begin{trivlist}
		\item[\hskip \labelsep {\bfseries #1}]}{\end{trivlist}}

\newenvironment{myindentpar}[1]%
{\begin{list}{}%
		{\setlength{\leftmargin}{#1}}%
		\item[]%
	}
	{\end{list}}

% =====================================================================
% 12. Notación matemática y econométrica
% =====================================================================
\newcommand{\plim}{\operatornamewithlimits{plim}}
\newcommand{\Var}{\operatorname{Var}}
\newcommand{\pderiv}[2]{\frac{\partial #1}{\partial #2}}
\newcommand{\spderiv}[2]{\frac{\partial^2 #1}{\partial #2^2}}
\newcommand{\xderiv}[3]{\frac{\partial^2 #1}{\partial #2 \partial #3}}
\newcommand{\argmin}{\operatornamewithlimits{arg\,min}}
\newcommand{\argmax}{\operatornamewithlimits{arg\,max}}
\newcommand{\amax}{\operatornamewithlimits{max}}
\newcommand{\amin}{\operatornamewithlimits{min}}
\newcommand*\circled[1]{\tikz[baseline=(char.base)]{\node[shape=circle,draw,inner sep=1pt] (char) {#1};}}
\newcommand{\indep}{\perp \!\!\! \perp}
\makeatletter
\newcommand{\distas}[1]{\mathbin{\overset{#1}{\kern\z@\sim}}}%
\newsavebox{\mybox}\newsavebox{\mysim}
\newcommand{\distras}[1]{%
	\savebox{\mybox}{\hbox{\kern3pt$\scriptstyle#1$\kern3pt}}%
	\savebox{\mysim}{\hbox{$\sim$}}%
	\mathbin{\overset{#1}{\kern\z@\resizebox{\wd\mybox}{\ht\mysim}{$\sim$}}}%
}
\makeatother

\DeclareMathOperator{\EX}{\mathbb{E}}% expected value
\newcommand{\Prob}[1]{\text{Pr}\left[#1\right]}
\newcommand{\RK}[1]{\mathbb{R}^{#1}}
\newcommand{\HT}[1]{\mathbb{H}_{#1}}

% Vectores y matrices en negrilla (romana)
\newcommand{\xB}{\textbf{x}}
\newcommand{\yB}{\textbf{y}}
\newcommand{\zB}{\textbf{z}}
\newcommand{\wB}{\textbf{w}}
\newcommand{\eB}{\textbf{e}}
\newcommand{\vB}{\textbf{v}}
\newcommand{\AB}{\textbf{A}}
\newcommand{\XB}{\textbf{X}}
\newcommand{\YB}{\textbf{Y}}
\newcommand{\ZB}{\textbf{Z}}
\newcommand{\QB}{\textbf{Q}}
\newcommand{\WB}{\textbf{W}}

% Vectores griegos en negrilla
\newcommand{\betaB}{\boldsymbol{\beta}}
\newcommand{\gammaB}{\boldsymbol{\gamma}}
\newcommand{\alphaB}{\boldsymbol{\alpha}}
\newcommand{\thetaB}{\boldsymbol{\theta}}
\newcommand{\deltaB}{\boldsymbol{\delta}}
\newcommand{\phiB}{\boldsymbol{\phi}}
\newcommand{\epsilonB}{\boldsymbol{\epsilon}}
\newcommand{\etaB}{\boldsymbol{\eta}}
\newcommand{\betaBH}{\boldsymbol{\hat{\beta}}}
\newcommand{\betaH}{\hat{\beta}}

% Colores de énfasis
\newcommand{\brown}[1]{\textcolor{brown}{#1}}
\newcommand{\blue}[1]{\textcolor{blue}{#1}}
\newcommand{\red}[1]{\textcolor{red}{#1}}

% Espaciados usados en itemize/enumerate
\newcommand{\smallspace}{\vspace{0.1cm}}
\newcommand{\mediumspace}{\vspace{0.3cm}}
\newcommand{\largespace}{\vspace{6cm}}

% Tamaños de fuente auxiliares
\newcommand\fontvia{\fontsize{4}{6}\selectfont}
\newcommand\fontvib{\fontsize{5}{8}\selectfont}
\newcommand\fontvic{\fontsize{7}{10}\selectfont}
\newcommand\fontvid{\fontsize{8}{11}\selectfont}
\newcommand\fontvie{\fontsize{10}{11}\selectfont}

% Celdas de mapa de calor estilizado para la intuición de shift-share
% Monotone (brown) cell: value 0..100 -> brown intensity
\newcommand{\sharecell}[3]{%
  \pgfmathtruncatemacro{\vv}{#3}%
  \fill[brown!\vv!white] (#1,#2) rectangle ({#1+1},{#2+1});%
  \draw[gray!40] (#1,#2) rectangle ({#1+1},{#2+1});%
}
% Diverging (blue/red) cell: value -50..50 -> blue (neg) or red (pos)
\newcommand{\divcell}[3]{%
  \pgfmathtruncatemacro{\absvv}{abs(#3)*2}%
  \pgfmathtruncatemacro{\negflag}{(#3)<0 ? 1 : 0}%
  \ifnum\negflag=1
    \fill[blue!\absvv!white] (#1,#2) rectangle ({#1+1},{#2+1});%
  \else
    \fill[red!\absvv!white] (#1,#2) rectangle ({#1+1},{#2+1});%
  \fi
  \draw[gray!40] (#1,#2) rectangle ({#1+1},{#2+1});%
}
% Uniform cell: same color regardless of (col,row)
\newcommand{\unifcell}[3]{%
  \pgfmathtruncatemacro{\vv}{#3}%
  \fill[red!\vv!white] (#1,#2) rectangle ({#1+1},{#2+1});%
  \draw[gray!40] (#1,#2) rectangle ({#1+1},{#2+1});%
}

% =====================================================================
% 13. Datos de portada (comunes a todas las clases)
% =====================================================================
\title{Econometría Avanzada}
\author{Manuel Fernández}
\date{\today}
%\titlegraphic{\hfill\includegraphics[height=1.5cm]{logo.pdf}}

%   3) \subtitle{...}                   (título de la clase puntual)
% Todo lo demás (paquetes, macros, tema, título, autor, hypersetup, etc.)
% vive aquí para poder hacer cambios globales en un solo lugar.
% =====================================================================

% ---------------------------------------------------------------------
% 1. Codificación, idioma y fuentes
% ---------------------------------------------------------------------
\usepackage{etex}
\usepackage[utf8]{inputenc}
\usepackage[T1]{fontenc}
\usepackage[spanish,es-nodecimaldot]{babel}
\usepackage{lmodern}
\usepackage{type1cm}
\usepackage{textcomp}

% ---------------------------------------------------------------------
% 2. Matemáticas y símbolos
% ---------------------------------------------------------------------
\usepackage{amsmath,amssymb,amsthm,amsfonts}
\usepackage{bbm}
\usepackage{bm}

% ---------------------------------------------------------------------
% 3. Bibliografía y citas
% ---------------------------------------------------------------------
\usepackage[round]{natbib}
\bibliographystyle{erae}
\usepackage{bibentry}
\usepackage{url}
\let\htmlstyloaded=\relax

% Las citas se muestran en marrón para distinguirlas del contenido.
\let\oldcitet=\citet
\renewcommand{\citet}[1]{\textcolor{brown}{\oldcitet{#1}}}
\let\oldcitep=\citep
\renewcommand{\citep}[1]{\textcolor{brown}{\oldcitep{#1}}}
\let\oldcite=\cite
\renewcommand{\cite}[1]{\textcolor{brown}{\oldcite{#1}}}

% ---------------------------------------------------------------------
% 4. Figuras, tablas y colores
% ---------------------------------------------------------------------
\usepackage{float}
\usepackage{graphicx}
\usepackage{subcaption}
\usepackage[font={footnotesize}]{caption}
\usepackage{lscape}
\usepackage{longtable}
\usepackage{booktabs}
\usepackage{multirow}
\usepackage{array}
\usepackage[flushleft]{threeparttable}
\usepackage{color, colortbl}
\usepackage{eso-pic}
\usepackage{chngcntr}
\usepackage{adjustbox}
\usepackage{pdfpages}
\usepackage{media9}

% siunitx: usado para alinear/formatear números en tablas
\usepackage{siunitx}
\sisetup{
	detect-mode,
	tight-spacing		= true,
	group-digits		= false ,
	input-signs		= ,
	input-symbols		= ( ) [ ] - + *,
	input-open-uncertainty	= ,
	input-close-uncertainty	= ,
	table-align-text-post	= false
}

% ---------------------------------------------------------------------
% 5. TikZ / PGFPlots
% ---------------------------------------------------------------------
\usepackage{tikz}
\usepackage{pgfplots}
\usepgfplotslibrary{dateplot}
\usetikzlibrary{arrows,calc, patterns, positioning, shapes.geometric, decorations.pathreplacing,decorations.markings}

% ---------------------------------------------------------------------
% 6. Misceláneos de formato
% ---------------------------------------------------------------------
\usepackage{setspace}
\usepackage{lipsum}
\usepackage{framed}
\usepackage{mdframed}
\usepackage{fancyhdr}
\usepackage{beamerfoils}
\usepackage{pifont}
\usepackage{ragged2e}
\usepackage{xspace}
\usepackage{appendixnumberbeamer}
\usepackage[scale=2]{ccicons}
\usepackage[most]{tcolorbox}
\usepackage{enumitem}
\setbeamertemplate{note page}[plain]

% ---------------------------------------------------------------------
% 7. Hyperref (debe cargarse después de la mayoría de paquetes)
% ---------------------------------------------------------------------
\usepackage{hyperref}
\hypersetup {
	pdftitle={General},
	pdfauthor={Manuel Fern\'{a}ndez},
	colorlinks=true,
	citecolor=brown,
	linkcolor=black,
	urlcolor=blue
}
\makeatletter
\renewcommand{\hyper@natlinkstart}[1]{}
\renewcommand{\hyper@natlinkend}{}
\renewcommand{\hyper@natlinkbreak}[2]{#1}
\makeatother

% =====================================================================
% 8. Tema Beamer (metropolis)
% =====================================================================
\usetheme{metropolis}
\newcommand{\themename}{\textbf{\textsc{metropolis}}\xspace}
\metroset{sectionpage=none}

\setbeamerfont{citep}{family=\rm}
\setbeamerfont{caption}{size=\scriptsize}
\setbeamerfont{normal text}{size=7pt}

%\setcounter{tocdepth}{1}
\setbeamertemplate{section in toc}[sections numbered]
\AtBeginSection[]{
  \begin{frame}
  \vfill
  \centering
  \begin{beamercolorbox}[sep=8pt,center,shadow=true,rounded=true]{title}
    \usebeamerfont{title}\insertsectionhead\par%
  \end{beamercolorbox}
  \vfill
  \end{frame}
}
% NOTA: \setbeamercolor{button}, \setbeamercovered{invisible} y
% \AtBeginSubsection NO están aquí porque sólo algunas clases los usaban
% en el original (afectan botones de hipervínculo, el efecto de \pause,
% y los separadores de \subsection). Se preservan como "ajustes locales"
% en los .tex que ya los tenían, justo después de \input{preambulo}.

% =====================================================================
% 9. Macros de Estout (tablas de resultados)
% =====================================================================
\newcommand{\sym}[1]{\rlap{#1}}% Thanks to David Carlisle
\let\estinput=\input% define a new input command so that we can still flatten the document
\newcommand{\estwide}[3]{
	\vspace{.75ex}{
		\begin{tabular*}
			{\textwidth}{@{\hskip\tabcolsep\extracolsep\fill}l*{#2}{#3}}
			\toprule
			\estinput{#1}
			\bottomrule
			\addlinespace[.75ex]
		\end{tabular*}
	}
}
\newcommand{\estauto}[3]{
	\vspace{.75ex}{
		\begin{tabular}{l*{#2}{#3}}
			\toprule
			\estinput{#1}
			\bottomrule
			\addlinespace[.75ex]
		\end{tabular}
	}
}
% Allow line breaks with \\ in specialcells
\newcommand{\specialcell}[2][c]{%
	\begin{tabular}[#1]{@{}c@{}}#2\end{tabular}}

% =====================================================================
% 10. Notas y fuentes al pie de figuras/tablas
% =====================================================================
% Note/Source/Text after Tables
\newcommand{\figtext}[1]{
	\vspace{-1.9ex}
	\captionsetup{justification=justified,font=footnotesize}
	\caption*{\hspace{6pt}\hangindent=1.5em #1}
}
\newcommand{\fignote}[1]{\figtext{\emph{Note:~}~#1}}
\newcommand{\figsource}[1]{\figtext{\emph{Source:~}~#1}}
% Add significance note with \starnote
\newcommand{\starnote}{\figtext{* p $<$ 0.10 ** p $<$ 0.05 *** p $<$ 0.01. }}

% =====================================================================
% 11. Entornos tipo teorema (en español)
% =====================================================================
\newtheorem{Prop}{Proposition}[section]
\newtheorem{Def}{Definition}
%\newtheorem{Assum}{Assumption}
%\newtheorem{Question}{Question}
%\newtheorem{Comment}{Comment}
%\theoremstyle{plain}
%\numberwithin{Assum}{section}
%\numberwithin{figure}{section}
%\numberwithin{equation}{section}
%\numberwithin{Def}{section}

\newenvironment{Teorema}[1][Teorema.]{\begin{trivlist}
		\item[\hskip \labelsep {\bfseries #1}]}{\end{trivlist}}
\newenvironment{Corolario}[1][Corolario.]{\begin{trivlist}
		\item[\hskip \labelsep {\bfseries #1}]}{\end{trivlist}}
\newenvironment{Definicion}[1][Definición.]{\begin{trivlist}
		\item[\hskip \labelsep {\bfseries #1}]}{\end{trivlist}}
\newenvironment{Lema}[1][Lema.]{\begin{trivlist}
		\item[\hskip \labelsep {\bfseries #1}]}{\end{trivlist}}
\newenvironment{Ejemplo}[1][Ejemplo.]{\begin{trivlist}
		\item[\hskip \labelsep {\bfseries #1}]}{\end{trivlist}}

\newenvironment{myindentpar}[1]%
{\begin{list}{}%
		{\setlength{\leftmargin}{#1}}%
		\item[]%
	}
	{\end{list}}

% =====================================================================
% 12. Notación matemática y econométrica
% =====================================================================
\newcommand{\plim}{\operatornamewithlimits{plim}}
\newcommand{\Var}{\operatorname{Var}}
\newcommand{\pderiv}[2]{\frac{\partial #1}{\partial #2}}
\newcommand{\spderiv}[2]{\frac{\partial^2 #1}{\partial #2^2}}
\newcommand{\xderiv}[3]{\frac{\partial^2 #1}{\partial #2 \partial #3}}
\newcommand{\argmin}{\operatornamewithlimits{arg\,min}}
\newcommand{\argmax}{\operatornamewithlimits{arg\,max}}
\newcommand{\amax}{\operatornamewithlimits{max}}
\newcommand{\amin}{\operatornamewithlimits{min}}
\newcommand*\circled[1]{\tikz[baseline=(char.base)]{\node[shape=circle,draw,inner sep=1pt] (char) {#1};}}
\newcommand{\indep}{\perp \!\!\! \perp}
\makeatletter
\newcommand{\distas}[1]{\mathbin{\overset{#1}{\kern\z@\sim}}}%
\newsavebox{\mybox}\newsavebox{\mysim}
\newcommand{\distras}[1]{%
	\savebox{\mybox}{\hbox{\kern3pt$\scriptstyle#1$\kern3pt}}%
	\savebox{\mysim}{\hbox{$\sim$}}%
	\mathbin{\overset{#1}{\kern\z@\resizebox{\wd\mybox}{\ht\mysim}{$\sim$}}}%
}
\makeatother

\DeclareMathOperator{\EX}{\mathbb{E}}% expected value
\newcommand{\Prob}[1]{\text{Pr}\left[#1\right]}
\newcommand{\RK}[1]{\mathbb{R}^{#1}}
\newcommand{\HT}[1]{\mathbb{H}_{#1}}

% Vectores y matrices en negrilla (romana)
\newcommand{\xB}{\textbf{x}}
\newcommand{\yB}{\textbf{y}}
\newcommand{\zB}{\textbf{z}}
\newcommand{\wB}{\textbf{w}}
\newcommand{\eB}{\textbf{e}}
\newcommand{\vB}{\textbf{v}}
\newcommand{\AB}{\textbf{A}}
\newcommand{\XB}{\textbf{X}}
\newcommand{\YB}{\textbf{Y}}
\newcommand{\ZB}{\textbf{Z}}
\newcommand{\QB}{\textbf{Q}}
\newcommand{\WB}{\textbf{W}}

% Vectores griegos en negrilla
\newcommand{\betaB}{\boldsymbol{\beta}}
\newcommand{\gammaB}{\boldsymbol{\gamma}}
\newcommand{\alphaB}{\boldsymbol{\alpha}}
\newcommand{\thetaB}{\boldsymbol{\theta}}
\newcommand{\deltaB}{\boldsymbol{\delta}}
\newcommand{\phiB}{\boldsymbol{\phi}}
\newcommand{\epsilonB}{\boldsymbol{\epsilon}}
\newcommand{\etaB}{\boldsymbol{\eta}}
\newcommand{\betaBH}{\boldsymbol{\hat{\beta}}}
\newcommand{\betaH}{\hat{\beta}}

% Colores de énfasis
\newcommand{\brown}[1]{\textcolor{brown}{#1}}
\newcommand{\blue}[1]{\textcolor{blue}{#1}}
\newcommand{\red}[1]{\textcolor{red}{#1}}

% Espaciados usados en itemize/enumerate
\newcommand{\smallspace}{\vspace{0.1cm}}
\newcommand{\mediumspace}{\vspace{0.3cm}}
\newcommand{\largespace}{\vspace{6cm}}

% Tamaños de fuente auxiliares
\newcommand\fontvia{\fontsize{4}{6}\selectfont}
\newcommand\fontvib{\fontsize{5}{8}\selectfont}
\newcommand\fontvic{\fontsize{7}{10}\selectfont}
\newcommand\fontvid{\fontsize{8}{11}\selectfont}
\newcommand\fontvie{\fontsize{10}{11}\selectfont}

% Celdas de mapa de calor estilizado para la intuición de shift-share
% Monotone (brown) cell: value 0..100 -> brown intensity
\newcommand{\sharecell}[3]{%
  \pgfmathtruncatemacro{\vv}{#3}%
  \fill[brown!\vv!white] (#1,#2) rectangle ({#1+1},{#2+1});%
  \draw[gray!40] (#1,#2) rectangle ({#1+1},{#2+1});%
}
% Diverging (blue/red) cell: value -50..50 -> blue (neg) or red (pos)
\newcommand{\divcell}[3]{%
  \pgfmathtruncatemacro{\absvv}{abs(#3)*2}%
  \pgfmathtruncatemacro{\negflag}{(#3)<0 ? 1 : 0}%
  \ifnum\negflag=1
    \fill[blue!\absvv!white] (#1,#2) rectangle ({#1+1},{#2+1});%
  \else
    \fill[red!\absvv!white] (#1,#2) rectangle ({#1+1},{#2+1});%
  \fi
  \draw[gray!40] (#1,#2) rectangle ({#1+1},{#2+1});%
}
% Uniform cell: same color regardless of (col,row)
\newcommand{\unifcell}[3]{%
  \pgfmathtruncatemacro{\vv}{#3}%
  \fill[red!\vv!white] (#1,#2) rectangle ({#1+1},{#2+1});%
  \draw[gray!40] (#1,#2) rectangle ({#1+1},{#2+1});%
}

% =====================================================================
% 13. Datos de portada (comunes a todas las clases)
% =====================================================================
\title{Econometría Avanzada}
\author{Manuel Fernández}
\date{\today}
%\titlegraphic{\hfill\includegraphics[height=1.5cm]{logo.pdf}}

%   3) \subtitle{...}                   (título de la clase puntual)
% Todo lo demás (paquetes, macros, tema, título, autor, hypersetup, etc.)
% vive aquí para poder hacer cambios globales en un solo lugar.
% =====================================================================

% ---------------------------------------------------------------------
% 1. Codificación, idioma y fuentes
% ---------------------------------------------------------------------
\usepackage{etex}
\usepackage[utf8]{inputenc}
\usepackage[T1]{fontenc}
\usepackage[spanish,es-nodecimaldot]{babel}
\usepackage{lmodern}
\usepackage{type1cm}
\usepackage{textcomp}

% ---------------------------------------------------------------------
% 2. Matemáticas y símbolos
% ---------------------------------------------------------------------
\usepackage{amsmath,amssymb,amsthm,amsfonts}
\usepackage{bbm}
\usepackage{bm}

% ---------------------------------------------------------------------
% 3. Bibliografía y citas
% ---------------------------------------------------------------------
\usepackage[round]{natbib}
\bibliographystyle{erae}
\usepackage{bibentry}
\usepackage{url}
\let\htmlstyloaded=\relax

% Las citas se muestran en marrón para distinguirlas del contenido.
\let\oldcitet=\citet
\renewcommand{\citet}[1]{\textcolor{brown}{\oldcitet{#1}}}
\let\oldcitep=\citep
\renewcommand{\citep}[1]{\textcolor{brown}{\oldcitep{#1}}}
\let\oldcite=\cite
\renewcommand{\cite}[1]{\textcolor{brown}{\oldcite{#1}}}

% ---------------------------------------------------------------------
% 4. Figuras, tablas y colores
% ---------------------------------------------------------------------
\usepackage{float}
\usepackage{graphicx}
\usepackage{subcaption}
\usepackage[font={footnotesize}]{caption}
\usepackage{lscape}
\usepackage{longtable}
\usepackage{booktabs}
\usepackage{multirow}
\usepackage{array}
\usepackage[flushleft]{threeparttable}
\usepackage{color, colortbl}
\usepackage{eso-pic}
\usepackage{chngcntr}
\usepackage{adjustbox}
\usepackage{pdfpages}
\usepackage{media9}

% siunitx: usado para alinear/formatear números en tablas
\usepackage{siunitx}
\sisetup{
	detect-mode,
	tight-spacing		= true,
	group-digits		= false ,
	input-signs		= ,
	input-symbols		= ( ) [ ] - + *,
	input-open-uncertainty	= ,
	input-close-uncertainty	= ,
	table-align-text-post	= false
}

% ---------------------------------------------------------------------
% 5. TikZ / PGFPlots
% ---------------------------------------------------------------------
\usepackage{tikz}
\usepackage{pgfplots}
\usepgfplotslibrary{dateplot}
\usetikzlibrary{arrows,calc, patterns, positioning, shapes.geometric, decorations.pathreplacing,decorations.markings}

% ---------------------------------------------------------------------
% 6. Misceláneos de formato
% ---------------------------------------------------------------------
\usepackage{setspace}
\usepackage{lipsum}
\usepackage{framed}
\usepackage{mdframed}
\usepackage{fancyhdr}
\usepackage{beamerfoils}
\usepackage{pifont}
\usepackage{ragged2e}
\usepackage{xspace}
\usepackage{appendixnumberbeamer}
\usepackage[scale=2]{ccicons}
\usepackage[most]{tcolorbox}
\usepackage{enumitem}
\setbeamertemplate{note page}[plain]

% ---------------------------------------------------------------------
% 7. Hyperref (debe cargarse después de la mayoría de paquetes)
% ---------------------------------------------------------------------
\usepackage{hyperref}
\hypersetup {
	pdftitle={General},
	pdfauthor={Manuel Fern\'{a}ndez},
	colorlinks=true,
	citecolor=brown,
	linkcolor=black,
	urlcolor=blue
}
\makeatletter
\renewcommand{\hyper@natlinkstart}[1]{}
\renewcommand{\hyper@natlinkend}{}
\renewcommand{\hyper@natlinkbreak}[2]{#1}
\makeatother

% =====================================================================
% 8. Tema Beamer (metropolis)
% =====================================================================
\usetheme{metropolis}
\newcommand{\themename}{\textbf{\textsc{metropolis}}\xspace}
\metroset{sectionpage=none}

\setbeamerfont{citep}{family=\rm}
\setbeamerfont{caption}{size=\scriptsize}
\setbeamerfont{normal text}{size=7pt}

%\setcounter{tocdepth}{1}
\setbeamertemplate{section in toc}[sections numbered]
\AtBeginSection[]{
  \begin{frame}
  \vfill
  \centering
  \begin{beamercolorbox}[sep=8pt,center,shadow=true,rounded=true]{title}
    \usebeamerfont{title}\insertsectionhead\par%
  \end{beamercolorbox}
  \vfill
  \end{frame}
}
% NOTA: \setbeamercolor{button}, \setbeamercovered{invisible} y
% \AtBeginSubsection NO están aquí porque sólo algunas clases los usaban
% en el original (afectan botones de hipervínculo, el efecto de \pause,
% y los separadores de \subsection). Se preservan como "ajustes locales"
% en los .tex que ya los tenían, justo después de \input{preambulo}.

% =====================================================================
% 9. Macros de Estout (tablas de resultados)
% =====================================================================
\newcommand{\sym}[1]{\rlap{#1}}% Thanks to David Carlisle
\let\estinput=\input% define a new input command so that we can still flatten the document
\newcommand{\estwide}[3]{
	\vspace{.75ex}{
		\begin{tabular*}
			{\textwidth}{@{\hskip\tabcolsep\extracolsep\fill}l*{#2}{#3}}
			\toprule
			\estinput{#1}
			\bottomrule
			\addlinespace[.75ex]
		\end{tabular*}
	}
}
\newcommand{\estauto}[3]{
	\vspace{.75ex}{
		\begin{tabular}{l*{#2}{#3}}
			\toprule
			\estinput{#1}
			\bottomrule
			\addlinespace[.75ex]
		\end{tabular}
	}
}
% Allow line breaks with \\ in specialcells
\newcommand{\specialcell}[2][c]{%
	\begin{tabular}[#1]{@{}c@{}}#2\end{tabular}}

% =====================================================================
% 10. Notas y fuentes al pie de figuras/tablas
% =====================================================================
% Note/Source/Text after Tables
\newcommand{\figtext}[1]{
	\vspace{-1.9ex}
	\captionsetup{justification=justified,font=footnotesize}
	\caption*{\hspace{6pt}\hangindent=1.5em #1}
}
\newcommand{\fignote}[1]{\figtext{\emph{Note:~}~#1}}
\newcommand{\figsource}[1]{\figtext{\emph{Source:~}~#1}}
% Add significance note with \starnote
\newcommand{\starnote}{\figtext{* p $<$ 0.10 ** p $<$ 0.05 *** p $<$ 0.01. }}

% =====================================================================
% 11. Entornos tipo teorema (en español)
% =====================================================================
\newtheorem{Prop}{Proposition}[section]
\newtheorem{Def}{Definition}
%\newtheorem{Assum}{Assumption}
%\newtheorem{Question}{Question}
%\newtheorem{Comment}{Comment}
%\theoremstyle{plain}
%\numberwithin{Assum}{section}
%\numberwithin{figure}{section}
%\numberwithin{equation}{section}
%\numberwithin{Def}{section}

\newenvironment{Teorema}[1][Teorema.]{\begin{trivlist}
		\item[\hskip \labelsep {\bfseries #1}]}{\end{trivlist}}
\newenvironment{Corolario}[1][Corolario.]{\begin{trivlist}
		\item[\hskip \labelsep {\bfseries #1}]}{\end{trivlist}}
\newenvironment{Definicion}[1][Definición.]{\begin{trivlist}
		\item[\hskip \labelsep {\bfseries #1}]}{\end{trivlist}}
\newenvironment{Lema}[1][Lema.]{\begin{trivlist}
		\item[\hskip \labelsep {\bfseries #1}]}{\end{trivlist}}
\newenvironment{Ejemplo}[1][Ejemplo.]{\begin{trivlist}
		\item[\hskip \labelsep {\bfseries #1}]}{\end{trivlist}}

\newenvironment{myindentpar}[1]%
{\begin{list}{}%
		{\setlength{\leftmargin}{#1}}%
		\item[]%
	}
	{\end{list}}

% =====================================================================
% 12. Notación matemática y econométrica
% =====================================================================
\newcommand{\plim}{\operatornamewithlimits{plim}}
\newcommand{\Var}{\operatorname{Var}}
\newcommand{\pderiv}[2]{\frac{\partial #1}{\partial #2}}
\newcommand{\spderiv}[2]{\frac{\partial^2 #1}{\partial #2^2}}
\newcommand{\xderiv}[3]{\frac{\partial^2 #1}{\partial #2 \partial #3}}
\newcommand{\argmin}{\operatornamewithlimits{arg\,min}}
\newcommand{\argmax}{\operatornamewithlimits{arg\,max}}
\newcommand{\amax}{\operatornamewithlimits{max}}
\newcommand{\amin}{\operatornamewithlimits{min}}
\newcommand*\circled[1]{\tikz[baseline=(char.base)]{\node[shape=circle,draw,inner sep=1pt] (char) {#1};}}
\newcommand{\indep}{\perp \!\!\! \perp}
\makeatletter
\newcommand{\distas}[1]{\mathbin{\overset{#1}{\kern\z@\sim}}}%
\newsavebox{\mybox}\newsavebox{\mysim}
\newcommand{\distras}[1]{%
	\savebox{\mybox}{\hbox{\kern3pt$\scriptstyle#1$\kern3pt}}%
	\savebox{\mysim}{\hbox{$\sim$}}%
	\mathbin{\overset{#1}{\kern\z@\resizebox{\wd\mybox}{\ht\mysim}{$\sim$}}}%
}
\makeatother

\DeclareMathOperator{\EX}{\mathbb{E}}% expected value
\newcommand{\Prob}[1]{\text{Pr}\left[#1\right]}
\newcommand{\RK}[1]{\mathbb{R}^{#1}}
\newcommand{\HT}[1]{\mathbb{H}_{#1}}

% Vectores y matrices en negrilla (romana)
\newcommand{\xB}{\textbf{x}}
\newcommand{\yB}{\textbf{y}}
\newcommand{\zB}{\textbf{z}}
\newcommand{\wB}{\textbf{w}}
\newcommand{\eB}{\textbf{e}}
\newcommand{\vB}{\textbf{v}}
\newcommand{\AB}{\textbf{A}}
\newcommand{\XB}{\textbf{X}}
\newcommand{\YB}{\textbf{Y}}
\newcommand{\ZB}{\textbf{Z}}
\newcommand{\QB}{\textbf{Q}}
\newcommand{\WB}{\textbf{W}}

% Vectores griegos en negrilla
\newcommand{\betaB}{\boldsymbol{\beta}}
\newcommand{\gammaB}{\boldsymbol{\gamma}}
\newcommand{\alphaB}{\boldsymbol{\alpha}}
\newcommand{\thetaB}{\boldsymbol{\theta}}
\newcommand{\deltaB}{\boldsymbol{\delta}}
\newcommand{\phiB}{\boldsymbol{\phi}}
\newcommand{\epsilonB}{\boldsymbol{\epsilon}}
\newcommand{\etaB}{\boldsymbol{\eta}}
\newcommand{\betaBH}{\boldsymbol{\hat{\beta}}}
\newcommand{\betaH}{\hat{\beta}}

% Colores de énfasis
\newcommand{\brown}[1]{\textcolor{brown}{#1}}
\newcommand{\blue}[1]{\textcolor{blue}{#1}}
\newcommand{\red}[1]{\textcolor{red}{#1}}

% Espaciados usados en itemize/enumerate
\newcommand{\smallspace}{\vspace{0.1cm}}
\newcommand{\mediumspace}{\vspace{0.3cm}}
\newcommand{\largespace}{\vspace{6cm}}

% Tamaños de fuente auxiliares
\newcommand\fontvia{\fontsize{4}{6}\selectfont}
\newcommand\fontvib{\fontsize{5}{8}\selectfont}
\newcommand\fontvic{\fontsize{7}{10}\selectfont}
\newcommand\fontvid{\fontsize{8}{11}\selectfont}
\newcommand\fontvie{\fontsize{10}{11}\selectfont}

% Celdas de mapa de calor estilizado para la intuición de shift-share
% Monotone (brown) cell: value 0..100 -> brown intensity
\newcommand{\sharecell}[3]{%
  \pgfmathtruncatemacro{\vv}{#3}%
  \fill[brown!\vv!white] (#1,#2) rectangle ({#1+1},{#2+1});%
  \draw[gray!40] (#1,#2) rectangle ({#1+1},{#2+1});%
}
% Diverging (blue/red) cell: value -50..50 -> blue (neg) or red (pos)
\newcommand{\divcell}[3]{%
  \pgfmathtruncatemacro{\absvv}{abs(#3)*2}%
  \pgfmathtruncatemacro{\negflag}{(#3)<0 ? 1 : 0}%
  \ifnum\negflag=1
    \fill[blue!\absvv!white] (#1,#2) rectangle ({#1+1},{#2+1});%
  \else
    \fill[red!\absvv!white] (#1,#2) rectangle ({#1+1},{#2+1});%
  \fi
  \draw[gray!40] (#1,#2) rectangle ({#1+1},{#2+1});%
}
% Uniform cell: same color regardless of (col,row)
\newcommand{\unifcell}[3]{%
  \pgfmathtruncatemacro{\vv}{#3}%
  \fill[red!\vv!white] (#1,#2) rectangle ({#1+1},{#2+1});%
  \draw[gray!40] (#1,#2) rectangle ({#1+1},{#2+1});%
}

% =====================================================================
% 13. Datos de portada (comunes a todas las clases)
% =====================================================================
\title{Econometría Avanzada}
\author{Manuel Fernández}
\date{\today}
%\titlegraphic{\hfill\includegraphics[height=1.5cm]{logo.pdf}}
 DESPUÉS de
% \documentclass[<tam>pt, xcolor=dvipsnames]{beamer} en cada archivo.
%
% Lo único que debe quedar en cada .tex individual es:
%   1) \documentclass[...]{beamer}      (el tamaño de fuente puede variar)
%   2) % =====================================================================
% Preámbulo común — Diapositivas Magistrales, Econometría Avanzada
% =====================================================================
% Se incluye con % =====================================================================
% Preámbulo común — Diapositivas Magistrales, Econometría Avanzada
% =====================================================================
% Se incluye con % =====================================================================
% Preámbulo común — Diapositivas Magistrales, Econometría Avanzada
% =====================================================================
% Se incluye con \input{preambulo/preambulo-magistral} DESPUÉS de
% \documentclass[<tam>pt, xcolor=dvipsnames]{beamer} en cada archivo.
%
% Lo único que debe quedar en cada .tex individual es:
%   1) \documentclass[...]{beamer}      (el tamaño de fuente puede variar)
%   2) \input{preambulo/preambulo-magistral}
%   3) \subtitle{...}                   (título de la clase puntual)
% Todo lo demás (paquetes, macros, tema, título, autor, hypersetup, etc.)
% vive aquí para poder hacer cambios globales en un solo lugar.
% =====================================================================

% ---------------------------------------------------------------------
% 1. Codificación, idioma y fuentes
% ---------------------------------------------------------------------
\usepackage{etex}
\usepackage[utf8]{inputenc}
\usepackage[T1]{fontenc}
\usepackage[spanish,es-nodecimaldot]{babel}
\usepackage{lmodern}
\usepackage{type1cm}
\usepackage{textcomp}

% ---------------------------------------------------------------------
% 2. Matemáticas y símbolos
% ---------------------------------------------------------------------
\usepackage{amsmath,amssymb,amsthm,amsfonts}
\usepackage{bbm}
\usepackage{bm}

% ---------------------------------------------------------------------
% 3. Bibliografía y citas
% ---------------------------------------------------------------------
\usepackage[round]{natbib}
\bibliographystyle{erae}
\usepackage{bibentry}
\usepackage{url}
\let\htmlstyloaded=\relax

% Las citas se muestran en marrón para distinguirlas del contenido.
\let\oldcitet=\citet
\renewcommand{\citet}[1]{\textcolor{brown}{\oldcitet{#1}}}
\let\oldcitep=\citep
\renewcommand{\citep}[1]{\textcolor{brown}{\oldcitep{#1}}}
\let\oldcite=\cite
\renewcommand{\cite}[1]{\textcolor{brown}{\oldcite{#1}}}

% ---------------------------------------------------------------------
% 4. Figuras, tablas y colores
% ---------------------------------------------------------------------
\usepackage{float}
\usepackage{graphicx}
\usepackage{subcaption}
\usepackage[font={footnotesize}]{caption}
\usepackage{lscape}
\usepackage{longtable}
\usepackage{booktabs}
\usepackage{multirow}
\usepackage{array}
\usepackage[flushleft]{threeparttable}
\usepackage{color, colortbl}
\usepackage{eso-pic}
\usepackage{chngcntr}
\usepackage{adjustbox}
\usepackage{pdfpages}
\usepackage{media9}

% siunitx: usado para alinear/formatear números en tablas
\usepackage{siunitx}
\sisetup{
	detect-mode,
	tight-spacing		= true,
	group-digits		= false ,
	input-signs		= ,
	input-symbols		= ( ) [ ] - + *,
	input-open-uncertainty	= ,
	input-close-uncertainty	= ,
	table-align-text-post	= false
}

% ---------------------------------------------------------------------
% 5. TikZ / PGFPlots
% ---------------------------------------------------------------------
\usepackage{tikz}
\usepackage{pgfplots}
\usepgfplotslibrary{dateplot}
\usetikzlibrary{arrows,calc, patterns, positioning, shapes.geometric, decorations.pathreplacing,decorations.markings}

% ---------------------------------------------------------------------
% 6. Misceláneos de formato
% ---------------------------------------------------------------------
\usepackage{setspace}
\usepackage{lipsum}
\usepackage{framed}
\usepackage{mdframed}
\usepackage{fancyhdr}
\usepackage{beamerfoils}
\usepackage{pifont}
\usepackage{ragged2e}
\usepackage{xspace}
\usepackage{appendixnumberbeamer}
\usepackage[scale=2]{ccicons}
\usepackage[most]{tcolorbox}
\usepackage{enumitem}
\setbeamertemplate{note page}[plain]

% ---------------------------------------------------------------------
% 7. Hyperref (debe cargarse después de la mayoría de paquetes)
% ---------------------------------------------------------------------
\usepackage{hyperref}
\hypersetup {
	pdftitle={General},
	pdfauthor={Manuel Fern\'{a}ndez},
	colorlinks=true,
	citecolor=brown,
	linkcolor=black,
	urlcolor=blue
}
\makeatletter
\renewcommand{\hyper@natlinkstart}[1]{}
\renewcommand{\hyper@natlinkend}{}
\renewcommand{\hyper@natlinkbreak}[2]{#1}
\makeatother

% =====================================================================
% 8. Tema Beamer (metropolis)
% =====================================================================
\usetheme{metropolis}
\newcommand{\themename}{\textbf{\textsc{metropolis}}\xspace}
\metroset{sectionpage=none}

\setbeamerfont{citep}{family=\rm}
\setbeamerfont{caption}{size=\scriptsize}
\setbeamerfont{normal text}{size=7pt}

%\setcounter{tocdepth}{1}
\setbeamertemplate{section in toc}[sections numbered]
\AtBeginSection[]{
  \begin{frame}
  \vfill
  \centering
  \begin{beamercolorbox}[sep=8pt,center,shadow=true,rounded=true]{title}
    \usebeamerfont{title}\insertsectionhead\par%
  \end{beamercolorbox}
  \vfill
  \end{frame}
}
% NOTA: \setbeamercolor{button}, \setbeamercovered{invisible} y
% \AtBeginSubsection NO están aquí porque sólo algunas clases los usaban
% en el original (afectan botones de hipervínculo, el efecto de \pause,
% y los separadores de \subsection). Se preservan como "ajustes locales"
% en los .tex que ya los tenían, justo después de \input{preambulo}.

% =====================================================================
% 9. Macros de Estout (tablas de resultados)
% =====================================================================
\newcommand{\sym}[1]{\rlap{#1}}% Thanks to David Carlisle
\let\estinput=\input% define a new input command so that we can still flatten the document
\newcommand{\estwide}[3]{
	\vspace{.75ex}{
		\begin{tabular*}
			{\textwidth}{@{\hskip\tabcolsep\extracolsep\fill}l*{#2}{#3}}
			\toprule
			\estinput{#1}
			\bottomrule
			\addlinespace[.75ex]
		\end{tabular*}
	}
}
\newcommand{\estauto}[3]{
	\vspace{.75ex}{
		\begin{tabular}{l*{#2}{#3}}
			\toprule
			\estinput{#1}
			\bottomrule
			\addlinespace[.75ex]
		\end{tabular}
	}
}
% Allow line breaks with \\ in specialcells
\newcommand{\specialcell}[2][c]{%
	\begin{tabular}[#1]{@{}c@{}}#2\end{tabular}}

% =====================================================================
% 10. Notas y fuentes al pie de figuras/tablas
% =====================================================================
% Note/Source/Text after Tables
\newcommand{\figtext}[1]{
	\vspace{-1.9ex}
	\captionsetup{justification=justified,font=footnotesize}
	\caption*{\hspace{6pt}\hangindent=1.5em #1}
}
\newcommand{\fignote}[1]{\figtext{\emph{Note:~}~#1}}
\newcommand{\figsource}[1]{\figtext{\emph{Source:~}~#1}}
% Add significance note with \starnote
\newcommand{\starnote}{\figtext{* p $<$ 0.10 ** p $<$ 0.05 *** p $<$ 0.01. }}

% =====================================================================
% 11. Entornos tipo teorema (en español)
% =====================================================================
\newtheorem{Prop}{Proposition}[section]
\newtheorem{Def}{Definition}
%\newtheorem{Assum}{Assumption}
%\newtheorem{Question}{Question}
%\newtheorem{Comment}{Comment}
%\theoremstyle{plain}
%\numberwithin{Assum}{section}
%\numberwithin{figure}{section}
%\numberwithin{equation}{section}
%\numberwithin{Def}{section}

\newenvironment{Teorema}[1][Teorema.]{\begin{trivlist}
		\item[\hskip \labelsep {\bfseries #1}]}{\end{trivlist}}
\newenvironment{Corolario}[1][Corolario.]{\begin{trivlist}
		\item[\hskip \labelsep {\bfseries #1}]}{\end{trivlist}}
\newenvironment{Definicion}[1][Definición.]{\begin{trivlist}
		\item[\hskip \labelsep {\bfseries #1}]}{\end{trivlist}}
\newenvironment{Lema}[1][Lema.]{\begin{trivlist}
		\item[\hskip \labelsep {\bfseries #1}]}{\end{trivlist}}
\newenvironment{Ejemplo}[1][Ejemplo.]{\begin{trivlist}
		\item[\hskip \labelsep {\bfseries #1}]}{\end{trivlist}}

\newenvironment{myindentpar}[1]%
{\begin{list}{}%
		{\setlength{\leftmargin}{#1}}%
		\item[]%
	}
	{\end{list}}

% =====================================================================
% 12. Notación matemática y econométrica
% =====================================================================
\newcommand{\plim}{\operatornamewithlimits{plim}}
\newcommand{\Var}{\operatorname{Var}}
\newcommand{\pderiv}[2]{\frac{\partial #1}{\partial #2}}
\newcommand{\spderiv}[2]{\frac{\partial^2 #1}{\partial #2^2}}
\newcommand{\xderiv}[3]{\frac{\partial^2 #1}{\partial #2 \partial #3}}
\newcommand{\argmin}{\operatornamewithlimits{arg\,min}}
\newcommand{\argmax}{\operatornamewithlimits{arg\,max}}
\newcommand{\amax}{\operatornamewithlimits{max}}
\newcommand{\amin}{\operatornamewithlimits{min}}
\newcommand*\circled[1]{\tikz[baseline=(char.base)]{\node[shape=circle,draw,inner sep=1pt] (char) {#1};}}
\newcommand{\indep}{\perp \!\!\! \perp}
\makeatletter
\newcommand{\distas}[1]{\mathbin{\overset{#1}{\kern\z@\sim}}}%
\newsavebox{\mybox}\newsavebox{\mysim}
\newcommand{\distras}[1]{%
	\savebox{\mybox}{\hbox{\kern3pt$\scriptstyle#1$\kern3pt}}%
	\savebox{\mysim}{\hbox{$\sim$}}%
	\mathbin{\overset{#1}{\kern\z@\resizebox{\wd\mybox}{\ht\mysim}{$\sim$}}}%
}
\makeatother

\DeclareMathOperator{\EX}{\mathbb{E}}% expected value
\newcommand{\Prob}[1]{\text{Pr}\left[#1\right]}
\newcommand{\RK}[1]{\mathbb{R}^{#1}}
\newcommand{\HT}[1]{\mathbb{H}_{#1}}

% Vectores y matrices en negrilla (romana)
\newcommand{\xB}{\textbf{x}}
\newcommand{\yB}{\textbf{y}}
\newcommand{\zB}{\textbf{z}}
\newcommand{\wB}{\textbf{w}}
\newcommand{\eB}{\textbf{e}}
\newcommand{\vB}{\textbf{v}}
\newcommand{\AB}{\textbf{A}}
\newcommand{\XB}{\textbf{X}}
\newcommand{\YB}{\textbf{Y}}
\newcommand{\ZB}{\textbf{Z}}
\newcommand{\QB}{\textbf{Q}}
\newcommand{\WB}{\textbf{W}}

% Vectores griegos en negrilla
\newcommand{\betaB}{\boldsymbol{\beta}}
\newcommand{\gammaB}{\boldsymbol{\gamma}}
\newcommand{\alphaB}{\boldsymbol{\alpha}}
\newcommand{\thetaB}{\boldsymbol{\theta}}
\newcommand{\deltaB}{\boldsymbol{\delta}}
\newcommand{\phiB}{\boldsymbol{\phi}}
\newcommand{\epsilonB}{\boldsymbol{\epsilon}}
\newcommand{\etaB}{\boldsymbol{\eta}}
\newcommand{\betaBH}{\boldsymbol{\hat{\beta}}}
\newcommand{\betaH}{\hat{\beta}}

% Colores de énfasis
\newcommand{\brown}[1]{\textcolor{brown}{#1}}
\newcommand{\blue}[1]{\textcolor{blue}{#1}}
\newcommand{\red}[1]{\textcolor{red}{#1}}

% Espaciados usados en itemize/enumerate
\newcommand{\smallspace}{\vspace{0.1cm}}
\newcommand{\mediumspace}{\vspace{0.3cm}}
\newcommand{\largespace}{\vspace{6cm}}

% Tamaños de fuente auxiliares
\newcommand\fontvia{\fontsize{4}{6}\selectfont}
\newcommand\fontvib{\fontsize{5}{8}\selectfont}
\newcommand\fontvic{\fontsize{7}{10}\selectfont}
\newcommand\fontvid{\fontsize{8}{11}\selectfont}
\newcommand\fontvie{\fontsize{10}{11}\selectfont}

% Celdas de mapa de calor estilizado para la intuición de shift-share
% Monotone (brown) cell: value 0..100 -> brown intensity
\newcommand{\sharecell}[3]{%
  \pgfmathtruncatemacro{\vv}{#3}%
  \fill[brown!\vv!white] (#1,#2) rectangle ({#1+1},{#2+1});%
  \draw[gray!40] (#1,#2) rectangle ({#1+1},{#2+1});%
}
% Diverging (blue/red) cell: value -50..50 -> blue (neg) or red (pos)
\newcommand{\divcell}[3]{%
  \pgfmathtruncatemacro{\absvv}{abs(#3)*2}%
  \pgfmathtruncatemacro{\negflag}{(#3)<0 ? 1 : 0}%
  \ifnum\negflag=1
    \fill[blue!\absvv!white] (#1,#2) rectangle ({#1+1},{#2+1});%
  \else
    \fill[red!\absvv!white] (#1,#2) rectangle ({#1+1},{#2+1});%
  \fi
  \draw[gray!40] (#1,#2) rectangle ({#1+1},{#2+1});%
}
% Uniform cell: same color regardless of (col,row)
\newcommand{\unifcell}[3]{%
  \pgfmathtruncatemacro{\vv}{#3}%
  \fill[red!\vv!white] (#1,#2) rectangle ({#1+1},{#2+1});%
  \draw[gray!40] (#1,#2) rectangle ({#1+1},{#2+1});%
}

% =====================================================================
% 13. Datos de portada (comunes a todas las clases)
% =====================================================================
\title{Econometría Avanzada}
\author{Manuel Fernández}
\date{\today}
%\titlegraphic{\hfill\includegraphics[height=1.5cm]{logo.pdf}}
 DESPUÉS de
% \documentclass[<tam>pt, xcolor=dvipsnames]{beamer} en cada archivo.
%
% Lo único que debe quedar en cada .tex individual es:
%   1) \documentclass[...]{beamer}      (el tamaño de fuente puede variar)
%   2) % =====================================================================
% Preámbulo común — Diapositivas Magistrales, Econometría Avanzada
% =====================================================================
% Se incluye con \input{preambulo/preambulo-magistral} DESPUÉS de
% \documentclass[<tam>pt, xcolor=dvipsnames]{beamer} en cada archivo.
%
% Lo único que debe quedar en cada .tex individual es:
%   1) \documentclass[...]{beamer}      (el tamaño de fuente puede variar)
%   2) \input{preambulo/preambulo-magistral}
%   3) \subtitle{...}                   (título de la clase puntual)
% Todo lo demás (paquetes, macros, tema, título, autor, hypersetup, etc.)
% vive aquí para poder hacer cambios globales en un solo lugar.
% =====================================================================

% ---------------------------------------------------------------------
% 1. Codificación, idioma y fuentes
% ---------------------------------------------------------------------
\usepackage{etex}
\usepackage[utf8]{inputenc}
\usepackage[T1]{fontenc}
\usepackage[spanish,es-nodecimaldot]{babel}
\usepackage{lmodern}
\usepackage{type1cm}
\usepackage{textcomp}

% ---------------------------------------------------------------------
% 2. Matemáticas y símbolos
% ---------------------------------------------------------------------
\usepackage{amsmath,amssymb,amsthm,amsfonts}
\usepackage{bbm}
\usepackage{bm}

% ---------------------------------------------------------------------
% 3. Bibliografía y citas
% ---------------------------------------------------------------------
\usepackage[round]{natbib}
\bibliographystyle{erae}
\usepackage{bibentry}
\usepackage{url}
\let\htmlstyloaded=\relax

% Las citas se muestran en marrón para distinguirlas del contenido.
\let\oldcitet=\citet
\renewcommand{\citet}[1]{\textcolor{brown}{\oldcitet{#1}}}
\let\oldcitep=\citep
\renewcommand{\citep}[1]{\textcolor{brown}{\oldcitep{#1}}}
\let\oldcite=\cite
\renewcommand{\cite}[1]{\textcolor{brown}{\oldcite{#1}}}

% ---------------------------------------------------------------------
% 4. Figuras, tablas y colores
% ---------------------------------------------------------------------
\usepackage{float}
\usepackage{graphicx}
\usepackage{subcaption}
\usepackage[font={footnotesize}]{caption}
\usepackage{lscape}
\usepackage{longtable}
\usepackage{booktabs}
\usepackage{multirow}
\usepackage{array}
\usepackage[flushleft]{threeparttable}
\usepackage{color, colortbl}
\usepackage{eso-pic}
\usepackage{chngcntr}
\usepackage{adjustbox}
\usepackage{pdfpages}
\usepackage{media9}

% siunitx: usado para alinear/formatear números en tablas
\usepackage{siunitx}
\sisetup{
	detect-mode,
	tight-spacing		= true,
	group-digits		= false ,
	input-signs		= ,
	input-symbols		= ( ) [ ] - + *,
	input-open-uncertainty	= ,
	input-close-uncertainty	= ,
	table-align-text-post	= false
}

% ---------------------------------------------------------------------
% 5. TikZ / PGFPlots
% ---------------------------------------------------------------------
\usepackage{tikz}
\usepackage{pgfplots}
\usepgfplotslibrary{dateplot}
\usetikzlibrary{arrows,calc, patterns, positioning, shapes.geometric, decorations.pathreplacing,decorations.markings}

% ---------------------------------------------------------------------
% 6. Misceláneos de formato
% ---------------------------------------------------------------------
\usepackage{setspace}
\usepackage{lipsum}
\usepackage{framed}
\usepackage{mdframed}
\usepackage{fancyhdr}
\usepackage{beamerfoils}
\usepackage{pifont}
\usepackage{ragged2e}
\usepackage{xspace}
\usepackage{appendixnumberbeamer}
\usepackage[scale=2]{ccicons}
\usepackage[most]{tcolorbox}
\usepackage{enumitem}
\setbeamertemplate{note page}[plain]

% ---------------------------------------------------------------------
% 7. Hyperref (debe cargarse después de la mayoría de paquetes)
% ---------------------------------------------------------------------
\usepackage{hyperref}
\hypersetup {
	pdftitle={General},
	pdfauthor={Manuel Fern\'{a}ndez},
	colorlinks=true,
	citecolor=brown,
	linkcolor=black,
	urlcolor=blue
}
\makeatletter
\renewcommand{\hyper@natlinkstart}[1]{}
\renewcommand{\hyper@natlinkend}{}
\renewcommand{\hyper@natlinkbreak}[2]{#1}
\makeatother

% =====================================================================
% 8. Tema Beamer (metropolis)
% =====================================================================
\usetheme{metropolis}
\newcommand{\themename}{\textbf{\textsc{metropolis}}\xspace}
\metroset{sectionpage=none}

\setbeamerfont{citep}{family=\rm}
\setbeamerfont{caption}{size=\scriptsize}
\setbeamerfont{normal text}{size=7pt}

%\setcounter{tocdepth}{1}
\setbeamertemplate{section in toc}[sections numbered]
\AtBeginSection[]{
  \begin{frame}
  \vfill
  \centering
  \begin{beamercolorbox}[sep=8pt,center,shadow=true,rounded=true]{title}
    \usebeamerfont{title}\insertsectionhead\par%
  \end{beamercolorbox}
  \vfill
  \end{frame}
}
% NOTA: \setbeamercolor{button}, \setbeamercovered{invisible} y
% \AtBeginSubsection NO están aquí porque sólo algunas clases los usaban
% en el original (afectan botones de hipervínculo, el efecto de \pause,
% y los separadores de \subsection). Se preservan como "ajustes locales"
% en los .tex que ya los tenían, justo después de \input{preambulo}.

% =====================================================================
% 9. Macros de Estout (tablas de resultados)
% =====================================================================
\newcommand{\sym}[1]{\rlap{#1}}% Thanks to David Carlisle
\let\estinput=\input% define a new input command so that we can still flatten the document
\newcommand{\estwide}[3]{
	\vspace{.75ex}{
		\begin{tabular*}
			{\textwidth}{@{\hskip\tabcolsep\extracolsep\fill}l*{#2}{#3}}
			\toprule
			\estinput{#1}
			\bottomrule
			\addlinespace[.75ex]
		\end{tabular*}
	}
}
\newcommand{\estauto}[3]{
	\vspace{.75ex}{
		\begin{tabular}{l*{#2}{#3}}
			\toprule
			\estinput{#1}
			\bottomrule
			\addlinespace[.75ex]
		\end{tabular}
	}
}
% Allow line breaks with \\ in specialcells
\newcommand{\specialcell}[2][c]{%
	\begin{tabular}[#1]{@{}c@{}}#2\end{tabular}}

% =====================================================================
% 10. Notas y fuentes al pie de figuras/tablas
% =====================================================================
% Note/Source/Text after Tables
\newcommand{\figtext}[1]{
	\vspace{-1.9ex}
	\captionsetup{justification=justified,font=footnotesize}
	\caption*{\hspace{6pt}\hangindent=1.5em #1}
}
\newcommand{\fignote}[1]{\figtext{\emph{Note:~}~#1}}
\newcommand{\figsource}[1]{\figtext{\emph{Source:~}~#1}}
% Add significance note with \starnote
\newcommand{\starnote}{\figtext{* p $<$ 0.10 ** p $<$ 0.05 *** p $<$ 0.01. }}

% =====================================================================
% 11. Entornos tipo teorema (en español)
% =====================================================================
\newtheorem{Prop}{Proposition}[section]
\newtheorem{Def}{Definition}
%\newtheorem{Assum}{Assumption}
%\newtheorem{Question}{Question}
%\newtheorem{Comment}{Comment}
%\theoremstyle{plain}
%\numberwithin{Assum}{section}
%\numberwithin{figure}{section}
%\numberwithin{equation}{section}
%\numberwithin{Def}{section}

\newenvironment{Teorema}[1][Teorema.]{\begin{trivlist}
		\item[\hskip \labelsep {\bfseries #1}]}{\end{trivlist}}
\newenvironment{Corolario}[1][Corolario.]{\begin{trivlist}
		\item[\hskip \labelsep {\bfseries #1}]}{\end{trivlist}}
\newenvironment{Definicion}[1][Definición.]{\begin{trivlist}
		\item[\hskip \labelsep {\bfseries #1}]}{\end{trivlist}}
\newenvironment{Lema}[1][Lema.]{\begin{trivlist}
		\item[\hskip \labelsep {\bfseries #1}]}{\end{trivlist}}
\newenvironment{Ejemplo}[1][Ejemplo.]{\begin{trivlist}
		\item[\hskip \labelsep {\bfseries #1}]}{\end{trivlist}}

\newenvironment{myindentpar}[1]%
{\begin{list}{}%
		{\setlength{\leftmargin}{#1}}%
		\item[]%
	}
	{\end{list}}

% =====================================================================
% 12. Notación matemática y econométrica
% =====================================================================
\newcommand{\plim}{\operatornamewithlimits{plim}}
\newcommand{\Var}{\operatorname{Var}}
\newcommand{\pderiv}[2]{\frac{\partial #1}{\partial #2}}
\newcommand{\spderiv}[2]{\frac{\partial^2 #1}{\partial #2^2}}
\newcommand{\xderiv}[3]{\frac{\partial^2 #1}{\partial #2 \partial #3}}
\newcommand{\argmin}{\operatornamewithlimits{arg\,min}}
\newcommand{\argmax}{\operatornamewithlimits{arg\,max}}
\newcommand{\amax}{\operatornamewithlimits{max}}
\newcommand{\amin}{\operatornamewithlimits{min}}
\newcommand*\circled[1]{\tikz[baseline=(char.base)]{\node[shape=circle,draw,inner sep=1pt] (char) {#1};}}
\newcommand{\indep}{\perp \!\!\! \perp}
\makeatletter
\newcommand{\distas}[1]{\mathbin{\overset{#1}{\kern\z@\sim}}}%
\newsavebox{\mybox}\newsavebox{\mysim}
\newcommand{\distras}[1]{%
	\savebox{\mybox}{\hbox{\kern3pt$\scriptstyle#1$\kern3pt}}%
	\savebox{\mysim}{\hbox{$\sim$}}%
	\mathbin{\overset{#1}{\kern\z@\resizebox{\wd\mybox}{\ht\mysim}{$\sim$}}}%
}
\makeatother

\DeclareMathOperator{\EX}{\mathbb{E}}% expected value
\newcommand{\Prob}[1]{\text{Pr}\left[#1\right]}
\newcommand{\RK}[1]{\mathbb{R}^{#1}}
\newcommand{\HT}[1]{\mathbb{H}_{#1}}

% Vectores y matrices en negrilla (romana)
\newcommand{\xB}{\textbf{x}}
\newcommand{\yB}{\textbf{y}}
\newcommand{\zB}{\textbf{z}}
\newcommand{\wB}{\textbf{w}}
\newcommand{\eB}{\textbf{e}}
\newcommand{\vB}{\textbf{v}}
\newcommand{\AB}{\textbf{A}}
\newcommand{\XB}{\textbf{X}}
\newcommand{\YB}{\textbf{Y}}
\newcommand{\ZB}{\textbf{Z}}
\newcommand{\QB}{\textbf{Q}}
\newcommand{\WB}{\textbf{W}}

% Vectores griegos en negrilla
\newcommand{\betaB}{\boldsymbol{\beta}}
\newcommand{\gammaB}{\boldsymbol{\gamma}}
\newcommand{\alphaB}{\boldsymbol{\alpha}}
\newcommand{\thetaB}{\boldsymbol{\theta}}
\newcommand{\deltaB}{\boldsymbol{\delta}}
\newcommand{\phiB}{\boldsymbol{\phi}}
\newcommand{\epsilonB}{\boldsymbol{\epsilon}}
\newcommand{\etaB}{\boldsymbol{\eta}}
\newcommand{\betaBH}{\boldsymbol{\hat{\beta}}}
\newcommand{\betaH}{\hat{\beta}}

% Colores de énfasis
\newcommand{\brown}[1]{\textcolor{brown}{#1}}
\newcommand{\blue}[1]{\textcolor{blue}{#1}}
\newcommand{\red}[1]{\textcolor{red}{#1}}

% Espaciados usados en itemize/enumerate
\newcommand{\smallspace}{\vspace{0.1cm}}
\newcommand{\mediumspace}{\vspace{0.3cm}}
\newcommand{\largespace}{\vspace{6cm}}

% Tamaños de fuente auxiliares
\newcommand\fontvia{\fontsize{4}{6}\selectfont}
\newcommand\fontvib{\fontsize{5}{8}\selectfont}
\newcommand\fontvic{\fontsize{7}{10}\selectfont}
\newcommand\fontvid{\fontsize{8}{11}\selectfont}
\newcommand\fontvie{\fontsize{10}{11}\selectfont}

% Celdas de mapa de calor estilizado para la intuición de shift-share
% Monotone (brown) cell: value 0..100 -> brown intensity
\newcommand{\sharecell}[3]{%
  \pgfmathtruncatemacro{\vv}{#3}%
  \fill[brown!\vv!white] (#1,#2) rectangle ({#1+1},{#2+1});%
  \draw[gray!40] (#1,#2) rectangle ({#1+1},{#2+1});%
}
% Diverging (blue/red) cell: value -50..50 -> blue (neg) or red (pos)
\newcommand{\divcell}[3]{%
  \pgfmathtruncatemacro{\absvv}{abs(#3)*2}%
  \pgfmathtruncatemacro{\negflag}{(#3)<0 ? 1 : 0}%
  \ifnum\negflag=1
    \fill[blue!\absvv!white] (#1,#2) rectangle ({#1+1},{#2+1});%
  \else
    \fill[red!\absvv!white] (#1,#2) rectangle ({#1+1},{#2+1});%
  \fi
  \draw[gray!40] (#1,#2) rectangle ({#1+1},{#2+1});%
}
% Uniform cell: same color regardless of (col,row)
\newcommand{\unifcell}[3]{%
  \pgfmathtruncatemacro{\vv}{#3}%
  \fill[red!\vv!white] (#1,#2) rectangle ({#1+1},{#2+1});%
  \draw[gray!40] (#1,#2) rectangle ({#1+1},{#2+1});%
}

% =====================================================================
% 13. Datos de portada (comunes a todas las clases)
% =====================================================================
\title{Econometría Avanzada}
\author{Manuel Fernández}
\date{\today}
%\titlegraphic{\hfill\includegraphics[height=1.5cm]{logo.pdf}}

%   3) \subtitle{...}                   (título de la clase puntual)
% Todo lo demás (paquetes, macros, tema, título, autor, hypersetup, etc.)
% vive aquí para poder hacer cambios globales en un solo lugar.
% =====================================================================

% ---------------------------------------------------------------------
% 1. Codificación, idioma y fuentes
% ---------------------------------------------------------------------
\usepackage{etex}
\usepackage[utf8]{inputenc}
\usepackage[T1]{fontenc}
\usepackage[spanish,es-nodecimaldot]{babel}
\usepackage{lmodern}
\usepackage{type1cm}
\usepackage{textcomp}

% ---------------------------------------------------------------------
% 2. Matemáticas y símbolos
% ---------------------------------------------------------------------
\usepackage{amsmath,amssymb,amsthm,amsfonts}
\usepackage{bbm}
\usepackage{bm}

% ---------------------------------------------------------------------
% 3. Bibliografía y citas
% ---------------------------------------------------------------------
\usepackage[round]{natbib}
\bibliographystyle{erae}
\usepackage{bibentry}
\usepackage{url}
\let\htmlstyloaded=\relax

% Las citas se muestran en marrón para distinguirlas del contenido.
\let\oldcitet=\citet
\renewcommand{\citet}[1]{\textcolor{brown}{\oldcitet{#1}}}
\let\oldcitep=\citep
\renewcommand{\citep}[1]{\textcolor{brown}{\oldcitep{#1}}}
\let\oldcite=\cite
\renewcommand{\cite}[1]{\textcolor{brown}{\oldcite{#1}}}

% ---------------------------------------------------------------------
% 4. Figuras, tablas y colores
% ---------------------------------------------------------------------
\usepackage{float}
\usepackage{graphicx}
\usepackage{subcaption}
\usepackage[font={footnotesize}]{caption}
\usepackage{lscape}
\usepackage{longtable}
\usepackage{booktabs}
\usepackage{multirow}
\usepackage{array}
\usepackage[flushleft]{threeparttable}
\usepackage{color, colortbl}
\usepackage{eso-pic}
\usepackage{chngcntr}
\usepackage{adjustbox}
\usepackage{pdfpages}
\usepackage{media9}

% siunitx: usado para alinear/formatear números en tablas
\usepackage{siunitx}
\sisetup{
	detect-mode,
	tight-spacing		= true,
	group-digits		= false ,
	input-signs		= ,
	input-symbols		= ( ) [ ] - + *,
	input-open-uncertainty	= ,
	input-close-uncertainty	= ,
	table-align-text-post	= false
}

% ---------------------------------------------------------------------
% 5. TikZ / PGFPlots
% ---------------------------------------------------------------------
\usepackage{tikz}
\usepackage{pgfplots}
\usepgfplotslibrary{dateplot}
\usetikzlibrary{arrows,calc, patterns, positioning, shapes.geometric, decorations.pathreplacing,decorations.markings}

% ---------------------------------------------------------------------
% 6. Misceláneos de formato
% ---------------------------------------------------------------------
\usepackage{setspace}
\usepackage{lipsum}
\usepackage{framed}
\usepackage{mdframed}
\usepackage{fancyhdr}
\usepackage{beamerfoils}
\usepackage{pifont}
\usepackage{ragged2e}
\usepackage{xspace}
\usepackage{appendixnumberbeamer}
\usepackage[scale=2]{ccicons}
\usepackage[most]{tcolorbox}
\usepackage{enumitem}
\setbeamertemplate{note page}[plain]

% ---------------------------------------------------------------------
% 7. Hyperref (debe cargarse después de la mayoría de paquetes)
% ---------------------------------------------------------------------
\usepackage{hyperref}
\hypersetup {
	pdftitle={General},
	pdfauthor={Manuel Fern\'{a}ndez},
	colorlinks=true,
	citecolor=brown,
	linkcolor=black,
	urlcolor=blue
}
\makeatletter
\renewcommand{\hyper@natlinkstart}[1]{}
\renewcommand{\hyper@natlinkend}{}
\renewcommand{\hyper@natlinkbreak}[2]{#1}
\makeatother

% =====================================================================
% 8. Tema Beamer (metropolis)
% =====================================================================
\usetheme{metropolis}
\newcommand{\themename}{\textbf{\textsc{metropolis}}\xspace}
\metroset{sectionpage=none}

\setbeamerfont{citep}{family=\rm}
\setbeamerfont{caption}{size=\scriptsize}
\setbeamerfont{normal text}{size=7pt}

%\setcounter{tocdepth}{1}
\setbeamertemplate{section in toc}[sections numbered]
\AtBeginSection[]{
  \begin{frame}
  \vfill
  \centering
  \begin{beamercolorbox}[sep=8pt,center,shadow=true,rounded=true]{title}
    \usebeamerfont{title}\insertsectionhead\par%
  \end{beamercolorbox}
  \vfill
  \end{frame}
}
% NOTA: \setbeamercolor{button}, \setbeamercovered{invisible} y
% \AtBeginSubsection NO están aquí porque sólo algunas clases los usaban
% en el original (afectan botones de hipervínculo, el efecto de \pause,
% y los separadores de \subsection). Se preservan como "ajustes locales"
% en los .tex que ya los tenían, justo después de \input{preambulo}.

% =====================================================================
% 9. Macros de Estout (tablas de resultados)
% =====================================================================
\newcommand{\sym}[1]{\rlap{#1}}% Thanks to David Carlisle
\let\estinput=\input% define a new input command so that we can still flatten the document
\newcommand{\estwide}[3]{
	\vspace{.75ex}{
		\begin{tabular*}
			{\textwidth}{@{\hskip\tabcolsep\extracolsep\fill}l*{#2}{#3}}
			\toprule
			\estinput{#1}
			\bottomrule
			\addlinespace[.75ex]
		\end{tabular*}
	}
}
\newcommand{\estauto}[3]{
	\vspace{.75ex}{
		\begin{tabular}{l*{#2}{#3}}
			\toprule
			\estinput{#1}
			\bottomrule
			\addlinespace[.75ex]
		\end{tabular}
	}
}
% Allow line breaks with \\ in specialcells
\newcommand{\specialcell}[2][c]{%
	\begin{tabular}[#1]{@{}c@{}}#2\end{tabular}}

% =====================================================================
% 10. Notas y fuentes al pie de figuras/tablas
% =====================================================================
% Note/Source/Text after Tables
\newcommand{\figtext}[1]{
	\vspace{-1.9ex}
	\captionsetup{justification=justified,font=footnotesize}
	\caption*{\hspace{6pt}\hangindent=1.5em #1}
}
\newcommand{\fignote}[1]{\figtext{\emph{Note:~}~#1}}
\newcommand{\figsource}[1]{\figtext{\emph{Source:~}~#1}}
% Add significance note with \starnote
\newcommand{\starnote}{\figtext{* p $<$ 0.10 ** p $<$ 0.05 *** p $<$ 0.01. }}

% =====================================================================
% 11. Entornos tipo teorema (en español)
% =====================================================================
\newtheorem{Prop}{Proposition}[section]
\newtheorem{Def}{Definition}
%\newtheorem{Assum}{Assumption}
%\newtheorem{Question}{Question}
%\newtheorem{Comment}{Comment}
%\theoremstyle{plain}
%\numberwithin{Assum}{section}
%\numberwithin{figure}{section}
%\numberwithin{equation}{section}
%\numberwithin{Def}{section}

\newenvironment{Teorema}[1][Teorema.]{\begin{trivlist}
		\item[\hskip \labelsep {\bfseries #1}]}{\end{trivlist}}
\newenvironment{Corolario}[1][Corolario.]{\begin{trivlist}
		\item[\hskip \labelsep {\bfseries #1}]}{\end{trivlist}}
\newenvironment{Definicion}[1][Definición.]{\begin{trivlist}
		\item[\hskip \labelsep {\bfseries #1}]}{\end{trivlist}}
\newenvironment{Lema}[1][Lema.]{\begin{trivlist}
		\item[\hskip \labelsep {\bfseries #1}]}{\end{trivlist}}
\newenvironment{Ejemplo}[1][Ejemplo.]{\begin{trivlist}
		\item[\hskip \labelsep {\bfseries #1}]}{\end{trivlist}}

\newenvironment{myindentpar}[1]%
{\begin{list}{}%
		{\setlength{\leftmargin}{#1}}%
		\item[]%
	}
	{\end{list}}

% =====================================================================
% 12. Notación matemática y econométrica
% =====================================================================
\newcommand{\plim}{\operatornamewithlimits{plim}}
\newcommand{\Var}{\operatorname{Var}}
\newcommand{\pderiv}[2]{\frac{\partial #1}{\partial #2}}
\newcommand{\spderiv}[2]{\frac{\partial^2 #1}{\partial #2^2}}
\newcommand{\xderiv}[3]{\frac{\partial^2 #1}{\partial #2 \partial #3}}
\newcommand{\argmin}{\operatornamewithlimits{arg\,min}}
\newcommand{\argmax}{\operatornamewithlimits{arg\,max}}
\newcommand{\amax}{\operatornamewithlimits{max}}
\newcommand{\amin}{\operatornamewithlimits{min}}
\newcommand*\circled[1]{\tikz[baseline=(char.base)]{\node[shape=circle,draw,inner sep=1pt] (char) {#1};}}
\newcommand{\indep}{\perp \!\!\! \perp}
\makeatletter
\newcommand{\distas}[1]{\mathbin{\overset{#1}{\kern\z@\sim}}}%
\newsavebox{\mybox}\newsavebox{\mysim}
\newcommand{\distras}[1]{%
	\savebox{\mybox}{\hbox{\kern3pt$\scriptstyle#1$\kern3pt}}%
	\savebox{\mysim}{\hbox{$\sim$}}%
	\mathbin{\overset{#1}{\kern\z@\resizebox{\wd\mybox}{\ht\mysim}{$\sim$}}}%
}
\makeatother

\DeclareMathOperator{\EX}{\mathbb{E}}% expected value
\newcommand{\Prob}[1]{\text{Pr}\left[#1\right]}
\newcommand{\RK}[1]{\mathbb{R}^{#1}}
\newcommand{\HT}[1]{\mathbb{H}_{#1}}

% Vectores y matrices en negrilla (romana)
\newcommand{\xB}{\textbf{x}}
\newcommand{\yB}{\textbf{y}}
\newcommand{\zB}{\textbf{z}}
\newcommand{\wB}{\textbf{w}}
\newcommand{\eB}{\textbf{e}}
\newcommand{\vB}{\textbf{v}}
\newcommand{\AB}{\textbf{A}}
\newcommand{\XB}{\textbf{X}}
\newcommand{\YB}{\textbf{Y}}
\newcommand{\ZB}{\textbf{Z}}
\newcommand{\QB}{\textbf{Q}}
\newcommand{\WB}{\textbf{W}}

% Vectores griegos en negrilla
\newcommand{\betaB}{\boldsymbol{\beta}}
\newcommand{\gammaB}{\boldsymbol{\gamma}}
\newcommand{\alphaB}{\boldsymbol{\alpha}}
\newcommand{\thetaB}{\boldsymbol{\theta}}
\newcommand{\deltaB}{\boldsymbol{\delta}}
\newcommand{\phiB}{\boldsymbol{\phi}}
\newcommand{\epsilonB}{\boldsymbol{\epsilon}}
\newcommand{\etaB}{\boldsymbol{\eta}}
\newcommand{\betaBH}{\boldsymbol{\hat{\beta}}}
\newcommand{\betaH}{\hat{\beta}}

% Colores de énfasis
\newcommand{\brown}[1]{\textcolor{brown}{#1}}
\newcommand{\blue}[1]{\textcolor{blue}{#1}}
\newcommand{\red}[1]{\textcolor{red}{#1}}

% Espaciados usados en itemize/enumerate
\newcommand{\smallspace}{\vspace{0.1cm}}
\newcommand{\mediumspace}{\vspace{0.3cm}}
\newcommand{\largespace}{\vspace{6cm}}

% Tamaños de fuente auxiliares
\newcommand\fontvia{\fontsize{4}{6}\selectfont}
\newcommand\fontvib{\fontsize{5}{8}\selectfont}
\newcommand\fontvic{\fontsize{7}{10}\selectfont}
\newcommand\fontvid{\fontsize{8}{11}\selectfont}
\newcommand\fontvie{\fontsize{10}{11}\selectfont}

% Celdas de mapa de calor estilizado para la intuición de shift-share
% Monotone (brown) cell: value 0..100 -> brown intensity
\newcommand{\sharecell}[3]{%
  \pgfmathtruncatemacro{\vv}{#3}%
  \fill[brown!\vv!white] (#1,#2) rectangle ({#1+1},{#2+1});%
  \draw[gray!40] (#1,#2) rectangle ({#1+1},{#2+1});%
}
% Diverging (blue/red) cell: value -50..50 -> blue (neg) or red (pos)
\newcommand{\divcell}[3]{%
  \pgfmathtruncatemacro{\absvv}{abs(#3)*2}%
  \pgfmathtruncatemacro{\negflag}{(#3)<0 ? 1 : 0}%
  \ifnum\negflag=1
    \fill[blue!\absvv!white] (#1,#2) rectangle ({#1+1},{#2+1});%
  \else
    \fill[red!\absvv!white] (#1,#2) rectangle ({#1+1},{#2+1});%
  \fi
  \draw[gray!40] (#1,#2) rectangle ({#1+1},{#2+1});%
}
% Uniform cell: same color regardless of (col,row)
\newcommand{\unifcell}[3]{%
  \pgfmathtruncatemacro{\vv}{#3}%
  \fill[red!\vv!white] (#1,#2) rectangle ({#1+1},{#2+1});%
  \draw[gray!40] (#1,#2) rectangle ({#1+1},{#2+1});%
}

% =====================================================================
% 13. Datos de portada (comunes a todas las clases)
% =====================================================================
\title{Econometría Avanzada}
\author{Manuel Fernández}
\date{\today}
%\titlegraphic{\hfill\includegraphics[height=1.5cm]{logo.pdf}}
 DESPUÉS de
% \documentclass[<tam>pt, xcolor=dvipsnames]{beamer} en cada archivo.
%
% Lo único que debe quedar en cada .tex individual es:
%   1) \documentclass[...]{beamer}      (el tamaño de fuente puede variar)
%   2) % =====================================================================
% Preámbulo común — Diapositivas Magistrales, Econometría Avanzada
% =====================================================================
% Se incluye con % =====================================================================
% Preámbulo común — Diapositivas Magistrales, Econometría Avanzada
% =====================================================================
% Se incluye con \input{preambulo/preambulo-magistral} DESPUÉS de
% \documentclass[<tam>pt, xcolor=dvipsnames]{beamer} en cada archivo.
%
% Lo único que debe quedar en cada .tex individual es:
%   1) \documentclass[...]{beamer}      (el tamaño de fuente puede variar)
%   2) \input{preambulo/preambulo-magistral}
%   3) \subtitle{...}                   (título de la clase puntual)
% Todo lo demás (paquetes, macros, tema, título, autor, hypersetup, etc.)
% vive aquí para poder hacer cambios globales en un solo lugar.
% =====================================================================

% ---------------------------------------------------------------------
% 1. Codificación, idioma y fuentes
% ---------------------------------------------------------------------
\usepackage{etex}
\usepackage[utf8]{inputenc}
\usepackage[T1]{fontenc}
\usepackage[spanish,es-nodecimaldot]{babel}
\usepackage{lmodern}
\usepackage{type1cm}
\usepackage{textcomp}

% ---------------------------------------------------------------------
% 2. Matemáticas y símbolos
% ---------------------------------------------------------------------
\usepackage{amsmath,amssymb,amsthm,amsfonts}
\usepackage{bbm}
\usepackage{bm}

% ---------------------------------------------------------------------
% 3. Bibliografía y citas
% ---------------------------------------------------------------------
\usepackage[round]{natbib}
\bibliographystyle{erae}
\usepackage{bibentry}
\usepackage{url}
\let\htmlstyloaded=\relax

% Las citas se muestran en marrón para distinguirlas del contenido.
\let\oldcitet=\citet
\renewcommand{\citet}[1]{\textcolor{brown}{\oldcitet{#1}}}
\let\oldcitep=\citep
\renewcommand{\citep}[1]{\textcolor{brown}{\oldcitep{#1}}}
\let\oldcite=\cite
\renewcommand{\cite}[1]{\textcolor{brown}{\oldcite{#1}}}

% ---------------------------------------------------------------------
% 4. Figuras, tablas y colores
% ---------------------------------------------------------------------
\usepackage{float}
\usepackage{graphicx}
\usepackage{subcaption}
\usepackage[font={footnotesize}]{caption}
\usepackage{lscape}
\usepackage{longtable}
\usepackage{booktabs}
\usepackage{multirow}
\usepackage{array}
\usepackage[flushleft]{threeparttable}
\usepackage{color, colortbl}
\usepackage{eso-pic}
\usepackage{chngcntr}
\usepackage{adjustbox}
\usepackage{pdfpages}
\usepackage{media9}

% siunitx: usado para alinear/formatear números en tablas
\usepackage{siunitx}
\sisetup{
	detect-mode,
	tight-spacing		= true,
	group-digits		= false ,
	input-signs		= ,
	input-symbols		= ( ) [ ] - + *,
	input-open-uncertainty	= ,
	input-close-uncertainty	= ,
	table-align-text-post	= false
}

% ---------------------------------------------------------------------
% 5. TikZ / PGFPlots
% ---------------------------------------------------------------------
\usepackage{tikz}
\usepackage{pgfplots}
\usepgfplotslibrary{dateplot}
\usetikzlibrary{arrows,calc, patterns, positioning, shapes.geometric, decorations.pathreplacing,decorations.markings}

% ---------------------------------------------------------------------
% 6. Misceláneos de formato
% ---------------------------------------------------------------------
\usepackage{setspace}
\usepackage{lipsum}
\usepackage{framed}
\usepackage{mdframed}
\usepackage{fancyhdr}
\usepackage{beamerfoils}
\usepackage{pifont}
\usepackage{ragged2e}
\usepackage{xspace}
\usepackage{appendixnumberbeamer}
\usepackage[scale=2]{ccicons}
\usepackage[most]{tcolorbox}
\usepackage{enumitem}
\setbeamertemplate{note page}[plain]

% ---------------------------------------------------------------------
% 7. Hyperref (debe cargarse después de la mayoría de paquetes)
% ---------------------------------------------------------------------
\usepackage{hyperref}
\hypersetup {
	pdftitle={General},
	pdfauthor={Manuel Fern\'{a}ndez},
	colorlinks=true,
	citecolor=brown,
	linkcolor=black,
	urlcolor=blue
}
\makeatletter
\renewcommand{\hyper@natlinkstart}[1]{}
\renewcommand{\hyper@natlinkend}{}
\renewcommand{\hyper@natlinkbreak}[2]{#1}
\makeatother

% =====================================================================
% 8. Tema Beamer (metropolis)
% =====================================================================
\usetheme{metropolis}
\newcommand{\themename}{\textbf{\textsc{metropolis}}\xspace}
\metroset{sectionpage=none}

\setbeamerfont{citep}{family=\rm}
\setbeamerfont{caption}{size=\scriptsize}
\setbeamerfont{normal text}{size=7pt}

%\setcounter{tocdepth}{1}
\setbeamertemplate{section in toc}[sections numbered]
\AtBeginSection[]{
  \begin{frame}
  \vfill
  \centering
  \begin{beamercolorbox}[sep=8pt,center,shadow=true,rounded=true]{title}
    \usebeamerfont{title}\insertsectionhead\par%
  \end{beamercolorbox}
  \vfill
  \end{frame}
}
% NOTA: \setbeamercolor{button}, \setbeamercovered{invisible} y
% \AtBeginSubsection NO están aquí porque sólo algunas clases los usaban
% en el original (afectan botones de hipervínculo, el efecto de \pause,
% y los separadores de \subsection). Se preservan como "ajustes locales"
% en los .tex que ya los tenían, justo después de \input{preambulo}.

% =====================================================================
% 9. Macros de Estout (tablas de resultados)
% =====================================================================
\newcommand{\sym}[1]{\rlap{#1}}% Thanks to David Carlisle
\let\estinput=\input% define a new input command so that we can still flatten the document
\newcommand{\estwide}[3]{
	\vspace{.75ex}{
		\begin{tabular*}
			{\textwidth}{@{\hskip\tabcolsep\extracolsep\fill}l*{#2}{#3}}
			\toprule
			\estinput{#1}
			\bottomrule
			\addlinespace[.75ex]
		\end{tabular*}
	}
}
\newcommand{\estauto}[3]{
	\vspace{.75ex}{
		\begin{tabular}{l*{#2}{#3}}
			\toprule
			\estinput{#1}
			\bottomrule
			\addlinespace[.75ex]
		\end{tabular}
	}
}
% Allow line breaks with \\ in specialcells
\newcommand{\specialcell}[2][c]{%
	\begin{tabular}[#1]{@{}c@{}}#2\end{tabular}}

% =====================================================================
% 10. Notas y fuentes al pie de figuras/tablas
% =====================================================================
% Note/Source/Text after Tables
\newcommand{\figtext}[1]{
	\vspace{-1.9ex}
	\captionsetup{justification=justified,font=footnotesize}
	\caption*{\hspace{6pt}\hangindent=1.5em #1}
}
\newcommand{\fignote}[1]{\figtext{\emph{Note:~}~#1}}
\newcommand{\figsource}[1]{\figtext{\emph{Source:~}~#1}}
% Add significance note with \starnote
\newcommand{\starnote}{\figtext{* p $<$ 0.10 ** p $<$ 0.05 *** p $<$ 0.01. }}

% =====================================================================
% 11. Entornos tipo teorema (en español)
% =====================================================================
\newtheorem{Prop}{Proposition}[section]
\newtheorem{Def}{Definition}
%\newtheorem{Assum}{Assumption}
%\newtheorem{Question}{Question}
%\newtheorem{Comment}{Comment}
%\theoremstyle{plain}
%\numberwithin{Assum}{section}
%\numberwithin{figure}{section}
%\numberwithin{equation}{section}
%\numberwithin{Def}{section}

\newenvironment{Teorema}[1][Teorema.]{\begin{trivlist}
		\item[\hskip \labelsep {\bfseries #1}]}{\end{trivlist}}
\newenvironment{Corolario}[1][Corolario.]{\begin{trivlist}
		\item[\hskip \labelsep {\bfseries #1}]}{\end{trivlist}}
\newenvironment{Definicion}[1][Definición.]{\begin{trivlist}
		\item[\hskip \labelsep {\bfseries #1}]}{\end{trivlist}}
\newenvironment{Lema}[1][Lema.]{\begin{trivlist}
		\item[\hskip \labelsep {\bfseries #1}]}{\end{trivlist}}
\newenvironment{Ejemplo}[1][Ejemplo.]{\begin{trivlist}
		\item[\hskip \labelsep {\bfseries #1}]}{\end{trivlist}}

\newenvironment{myindentpar}[1]%
{\begin{list}{}%
		{\setlength{\leftmargin}{#1}}%
		\item[]%
	}
	{\end{list}}

% =====================================================================
% 12. Notación matemática y econométrica
% =====================================================================
\newcommand{\plim}{\operatornamewithlimits{plim}}
\newcommand{\Var}{\operatorname{Var}}
\newcommand{\pderiv}[2]{\frac{\partial #1}{\partial #2}}
\newcommand{\spderiv}[2]{\frac{\partial^2 #1}{\partial #2^2}}
\newcommand{\xderiv}[3]{\frac{\partial^2 #1}{\partial #2 \partial #3}}
\newcommand{\argmin}{\operatornamewithlimits{arg\,min}}
\newcommand{\argmax}{\operatornamewithlimits{arg\,max}}
\newcommand{\amax}{\operatornamewithlimits{max}}
\newcommand{\amin}{\operatornamewithlimits{min}}
\newcommand*\circled[1]{\tikz[baseline=(char.base)]{\node[shape=circle,draw,inner sep=1pt] (char) {#1};}}
\newcommand{\indep}{\perp \!\!\! \perp}
\makeatletter
\newcommand{\distas}[1]{\mathbin{\overset{#1}{\kern\z@\sim}}}%
\newsavebox{\mybox}\newsavebox{\mysim}
\newcommand{\distras}[1]{%
	\savebox{\mybox}{\hbox{\kern3pt$\scriptstyle#1$\kern3pt}}%
	\savebox{\mysim}{\hbox{$\sim$}}%
	\mathbin{\overset{#1}{\kern\z@\resizebox{\wd\mybox}{\ht\mysim}{$\sim$}}}%
}
\makeatother

\DeclareMathOperator{\EX}{\mathbb{E}}% expected value
\newcommand{\Prob}[1]{\text{Pr}\left[#1\right]}
\newcommand{\RK}[1]{\mathbb{R}^{#1}}
\newcommand{\HT}[1]{\mathbb{H}_{#1}}

% Vectores y matrices en negrilla (romana)
\newcommand{\xB}{\textbf{x}}
\newcommand{\yB}{\textbf{y}}
\newcommand{\zB}{\textbf{z}}
\newcommand{\wB}{\textbf{w}}
\newcommand{\eB}{\textbf{e}}
\newcommand{\vB}{\textbf{v}}
\newcommand{\AB}{\textbf{A}}
\newcommand{\XB}{\textbf{X}}
\newcommand{\YB}{\textbf{Y}}
\newcommand{\ZB}{\textbf{Z}}
\newcommand{\QB}{\textbf{Q}}
\newcommand{\WB}{\textbf{W}}

% Vectores griegos en negrilla
\newcommand{\betaB}{\boldsymbol{\beta}}
\newcommand{\gammaB}{\boldsymbol{\gamma}}
\newcommand{\alphaB}{\boldsymbol{\alpha}}
\newcommand{\thetaB}{\boldsymbol{\theta}}
\newcommand{\deltaB}{\boldsymbol{\delta}}
\newcommand{\phiB}{\boldsymbol{\phi}}
\newcommand{\epsilonB}{\boldsymbol{\epsilon}}
\newcommand{\etaB}{\boldsymbol{\eta}}
\newcommand{\betaBH}{\boldsymbol{\hat{\beta}}}
\newcommand{\betaH}{\hat{\beta}}

% Colores de énfasis
\newcommand{\brown}[1]{\textcolor{brown}{#1}}
\newcommand{\blue}[1]{\textcolor{blue}{#1}}
\newcommand{\red}[1]{\textcolor{red}{#1}}

% Espaciados usados en itemize/enumerate
\newcommand{\smallspace}{\vspace{0.1cm}}
\newcommand{\mediumspace}{\vspace{0.3cm}}
\newcommand{\largespace}{\vspace{6cm}}

% Tamaños de fuente auxiliares
\newcommand\fontvia{\fontsize{4}{6}\selectfont}
\newcommand\fontvib{\fontsize{5}{8}\selectfont}
\newcommand\fontvic{\fontsize{7}{10}\selectfont}
\newcommand\fontvid{\fontsize{8}{11}\selectfont}
\newcommand\fontvie{\fontsize{10}{11}\selectfont}

% Celdas de mapa de calor estilizado para la intuición de shift-share
% Monotone (brown) cell: value 0..100 -> brown intensity
\newcommand{\sharecell}[3]{%
  \pgfmathtruncatemacro{\vv}{#3}%
  \fill[brown!\vv!white] (#1,#2) rectangle ({#1+1},{#2+1});%
  \draw[gray!40] (#1,#2) rectangle ({#1+1},{#2+1});%
}
% Diverging (blue/red) cell: value -50..50 -> blue (neg) or red (pos)
\newcommand{\divcell}[3]{%
  \pgfmathtruncatemacro{\absvv}{abs(#3)*2}%
  \pgfmathtruncatemacro{\negflag}{(#3)<0 ? 1 : 0}%
  \ifnum\negflag=1
    \fill[blue!\absvv!white] (#1,#2) rectangle ({#1+1},{#2+1});%
  \else
    \fill[red!\absvv!white] (#1,#2) rectangle ({#1+1},{#2+1});%
  \fi
  \draw[gray!40] (#1,#2) rectangle ({#1+1},{#2+1});%
}
% Uniform cell: same color regardless of (col,row)
\newcommand{\unifcell}[3]{%
  \pgfmathtruncatemacro{\vv}{#3}%
  \fill[red!\vv!white] (#1,#2) rectangle ({#1+1},{#2+1});%
  \draw[gray!40] (#1,#2) rectangle ({#1+1},{#2+1});%
}

% =====================================================================
% 13. Datos de portada (comunes a todas las clases)
% =====================================================================
\title{Econometría Avanzada}
\author{Manuel Fernández}
\date{\today}
%\titlegraphic{\hfill\includegraphics[height=1.5cm]{logo.pdf}}
 DESPUÉS de
% \documentclass[<tam>pt, xcolor=dvipsnames]{beamer} en cada archivo.
%
% Lo único que debe quedar en cada .tex individual es:
%   1) \documentclass[...]{beamer}      (el tamaño de fuente puede variar)
%   2) % =====================================================================
% Preámbulo común — Diapositivas Magistrales, Econometría Avanzada
% =====================================================================
% Se incluye con \input{preambulo/preambulo-magistral} DESPUÉS de
% \documentclass[<tam>pt, xcolor=dvipsnames]{beamer} en cada archivo.
%
% Lo único que debe quedar en cada .tex individual es:
%   1) \documentclass[...]{beamer}      (el tamaño de fuente puede variar)
%   2) \input{preambulo/preambulo-magistral}
%   3) \subtitle{...}                   (título de la clase puntual)
% Todo lo demás (paquetes, macros, tema, título, autor, hypersetup, etc.)
% vive aquí para poder hacer cambios globales en un solo lugar.
% =====================================================================

% ---------------------------------------------------------------------
% 1. Codificación, idioma y fuentes
% ---------------------------------------------------------------------
\usepackage{etex}
\usepackage[utf8]{inputenc}
\usepackage[T1]{fontenc}
\usepackage[spanish,es-nodecimaldot]{babel}
\usepackage{lmodern}
\usepackage{type1cm}
\usepackage{textcomp}

% ---------------------------------------------------------------------
% 2. Matemáticas y símbolos
% ---------------------------------------------------------------------
\usepackage{amsmath,amssymb,amsthm,amsfonts}
\usepackage{bbm}
\usepackage{bm}

% ---------------------------------------------------------------------
% 3. Bibliografía y citas
% ---------------------------------------------------------------------
\usepackage[round]{natbib}
\bibliographystyle{erae}
\usepackage{bibentry}
\usepackage{url}
\let\htmlstyloaded=\relax

% Las citas se muestran en marrón para distinguirlas del contenido.
\let\oldcitet=\citet
\renewcommand{\citet}[1]{\textcolor{brown}{\oldcitet{#1}}}
\let\oldcitep=\citep
\renewcommand{\citep}[1]{\textcolor{brown}{\oldcitep{#1}}}
\let\oldcite=\cite
\renewcommand{\cite}[1]{\textcolor{brown}{\oldcite{#1}}}

% ---------------------------------------------------------------------
% 4. Figuras, tablas y colores
% ---------------------------------------------------------------------
\usepackage{float}
\usepackage{graphicx}
\usepackage{subcaption}
\usepackage[font={footnotesize}]{caption}
\usepackage{lscape}
\usepackage{longtable}
\usepackage{booktabs}
\usepackage{multirow}
\usepackage{array}
\usepackage[flushleft]{threeparttable}
\usepackage{color, colortbl}
\usepackage{eso-pic}
\usepackage{chngcntr}
\usepackage{adjustbox}
\usepackage{pdfpages}
\usepackage{media9}

% siunitx: usado para alinear/formatear números en tablas
\usepackage{siunitx}
\sisetup{
	detect-mode,
	tight-spacing		= true,
	group-digits		= false ,
	input-signs		= ,
	input-symbols		= ( ) [ ] - + *,
	input-open-uncertainty	= ,
	input-close-uncertainty	= ,
	table-align-text-post	= false
}

% ---------------------------------------------------------------------
% 5. TikZ / PGFPlots
% ---------------------------------------------------------------------
\usepackage{tikz}
\usepackage{pgfplots}
\usepgfplotslibrary{dateplot}
\usetikzlibrary{arrows,calc, patterns, positioning, shapes.geometric, decorations.pathreplacing,decorations.markings}

% ---------------------------------------------------------------------
% 6. Misceláneos de formato
% ---------------------------------------------------------------------
\usepackage{setspace}
\usepackage{lipsum}
\usepackage{framed}
\usepackage{mdframed}
\usepackage{fancyhdr}
\usepackage{beamerfoils}
\usepackage{pifont}
\usepackage{ragged2e}
\usepackage{xspace}
\usepackage{appendixnumberbeamer}
\usepackage[scale=2]{ccicons}
\usepackage[most]{tcolorbox}
\usepackage{enumitem}
\setbeamertemplate{note page}[plain]

% ---------------------------------------------------------------------
% 7. Hyperref (debe cargarse después de la mayoría de paquetes)
% ---------------------------------------------------------------------
\usepackage{hyperref}
\hypersetup {
	pdftitle={General},
	pdfauthor={Manuel Fern\'{a}ndez},
	colorlinks=true,
	citecolor=brown,
	linkcolor=black,
	urlcolor=blue
}
\makeatletter
\renewcommand{\hyper@natlinkstart}[1]{}
\renewcommand{\hyper@natlinkend}{}
\renewcommand{\hyper@natlinkbreak}[2]{#1}
\makeatother

% =====================================================================
% 8. Tema Beamer (metropolis)
% =====================================================================
\usetheme{metropolis}
\newcommand{\themename}{\textbf{\textsc{metropolis}}\xspace}
\metroset{sectionpage=none}

\setbeamerfont{citep}{family=\rm}
\setbeamerfont{caption}{size=\scriptsize}
\setbeamerfont{normal text}{size=7pt}

%\setcounter{tocdepth}{1}
\setbeamertemplate{section in toc}[sections numbered]
\AtBeginSection[]{
  \begin{frame}
  \vfill
  \centering
  \begin{beamercolorbox}[sep=8pt,center,shadow=true,rounded=true]{title}
    \usebeamerfont{title}\insertsectionhead\par%
  \end{beamercolorbox}
  \vfill
  \end{frame}
}
% NOTA: \setbeamercolor{button}, \setbeamercovered{invisible} y
% \AtBeginSubsection NO están aquí porque sólo algunas clases los usaban
% en el original (afectan botones de hipervínculo, el efecto de \pause,
% y los separadores de \subsection). Se preservan como "ajustes locales"
% en los .tex que ya los tenían, justo después de \input{preambulo}.

% =====================================================================
% 9. Macros de Estout (tablas de resultados)
% =====================================================================
\newcommand{\sym}[1]{\rlap{#1}}% Thanks to David Carlisle
\let\estinput=\input% define a new input command so that we can still flatten the document
\newcommand{\estwide}[3]{
	\vspace{.75ex}{
		\begin{tabular*}
			{\textwidth}{@{\hskip\tabcolsep\extracolsep\fill}l*{#2}{#3}}
			\toprule
			\estinput{#1}
			\bottomrule
			\addlinespace[.75ex]
		\end{tabular*}
	}
}
\newcommand{\estauto}[3]{
	\vspace{.75ex}{
		\begin{tabular}{l*{#2}{#3}}
			\toprule
			\estinput{#1}
			\bottomrule
			\addlinespace[.75ex]
		\end{tabular}
	}
}
% Allow line breaks with \\ in specialcells
\newcommand{\specialcell}[2][c]{%
	\begin{tabular}[#1]{@{}c@{}}#2\end{tabular}}

% =====================================================================
% 10. Notas y fuentes al pie de figuras/tablas
% =====================================================================
% Note/Source/Text after Tables
\newcommand{\figtext}[1]{
	\vspace{-1.9ex}
	\captionsetup{justification=justified,font=footnotesize}
	\caption*{\hspace{6pt}\hangindent=1.5em #1}
}
\newcommand{\fignote}[1]{\figtext{\emph{Note:~}~#1}}
\newcommand{\figsource}[1]{\figtext{\emph{Source:~}~#1}}
% Add significance note with \starnote
\newcommand{\starnote}{\figtext{* p $<$ 0.10 ** p $<$ 0.05 *** p $<$ 0.01. }}

% =====================================================================
% 11. Entornos tipo teorema (en español)
% =====================================================================
\newtheorem{Prop}{Proposition}[section]
\newtheorem{Def}{Definition}
%\newtheorem{Assum}{Assumption}
%\newtheorem{Question}{Question}
%\newtheorem{Comment}{Comment}
%\theoremstyle{plain}
%\numberwithin{Assum}{section}
%\numberwithin{figure}{section}
%\numberwithin{equation}{section}
%\numberwithin{Def}{section}

\newenvironment{Teorema}[1][Teorema.]{\begin{trivlist}
		\item[\hskip \labelsep {\bfseries #1}]}{\end{trivlist}}
\newenvironment{Corolario}[1][Corolario.]{\begin{trivlist}
		\item[\hskip \labelsep {\bfseries #1}]}{\end{trivlist}}
\newenvironment{Definicion}[1][Definición.]{\begin{trivlist}
		\item[\hskip \labelsep {\bfseries #1}]}{\end{trivlist}}
\newenvironment{Lema}[1][Lema.]{\begin{trivlist}
		\item[\hskip \labelsep {\bfseries #1}]}{\end{trivlist}}
\newenvironment{Ejemplo}[1][Ejemplo.]{\begin{trivlist}
		\item[\hskip \labelsep {\bfseries #1}]}{\end{trivlist}}

\newenvironment{myindentpar}[1]%
{\begin{list}{}%
		{\setlength{\leftmargin}{#1}}%
		\item[]%
	}
	{\end{list}}

% =====================================================================
% 12. Notación matemática y econométrica
% =====================================================================
\newcommand{\plim}{\operatornamewithlimits{plim}}
\newcommand{\Var}{\operatorname{Var}}
\newcommand{\pderiv}[2]{\frac{\partial #1}{\partial #2}}
\newcommand{\spderiv}[2]{\frac{\partial^2 #1}{\partial #2^2}}
\newcommand{\xderiv}[3]{\frac{\partial^2 #1}{\partial #2 \partial #3}}
\newcommand{\argmin}{\operatornamewithlimits{arg\,min}}
\newcommand{\argmax}{\operatornamewithlimits{arg\,max}}
\newcommand{\amax}{\operatornamewithlimits{max}}
\newcommand{\amin}{\operatornamewithlimits{min}}
\newcommand*\circled[1]{\tikz[baseline=(char.base)]{\node[shape=circle,draw,inner sep=1pt] (char) {#1};}}
\newcommand{\indep}{\perp \!\!\! \perp}
\makeatletter
\newcommand{\distas}[1]{\mathbin{\overset{#1}{\kern\z@\sim}}}%
\newsavebox{\mybox}\newsavebox{\mysim}
\newcommand{\distras}[1]{%
	\savebox{\mybox}{\hbox{\kern3pt$\scriptstyle#1$\kern3pt}}%
	\savebox{\mysim}{\hbox{$\sim$}}%
	\mathbin{\overset{#1}{\kern\z@\resizebox{\wd\mybox}{\ht\mysim}{$\sim$}}}%
}
\makeatother

\DeclareMathOperator{\EX}{\mathbb{E}}% expected value
\newcommand{\Prob}[1]{\text{Pr}\left[#1\right]}
\newcommand{\RK}[1]{\mathbb{R}^{#1}}
\newcommand{\HT}[1]{\mathbb{H}_{#1}}

% Vectores y matrices en negrilla (romana)
\newcommand{\xB}{\textbf{x}}
\newcommand{\yB}{\textbf{y}}
\newcommand{\zB}{\textbf{z}}
\newcommand{\wB}{\textbf{w}}
\newcommand{\eB}{\textbf{e}}
\newcommand{\vB}{\textbf{v}}
\newcommand{\AB}{\textbf{A}}
\newcommand{\XB}{\textbf{X}}
\newcommand{\YB}{\textbf{Y}}
\newcommand{\ZB}{\textbf{Z}}
\newcommand{\QB}{\textbf{Q}}
\newcommand{\WB}{\textbf{W}}

% Vectores griegos en negrilla
\newcommand{\betaB}{\boldsymbol{\beta}}
\newcommand{\gammaB}{\boldsymbol{\gamma}}
\newcommand{\alphaB}{\boldsymbol{\alpha}}
\newcommand{\thetaB}{\boldsymbol{\theta}}
\newcommand{\deltaB}{\boldsymbol{\delta}}
\newcommand{\phiB}{\boldsymbol{\phi}}
\newcommand{\epsilonB}{\boldsymbol{\epsilon}}
\newcommand{\etaB}{\boldsymbol{\eta}}
\newcommand{\betaBH}{\boldsymbol{\hat{\beta}}}
\newcommand{\betaH}{\hat{\beta}}

% Colores de énfasis
\newcommand{\brown}[1]{\textcolor{brown}{#1}}
\newcommand{\blue}[1]{\textcolor{blue}{#1}}
\newcommand{\red}[1]{\textcolor{red}{#1}}

% Espaciados usados en itemize/enumerate
\newcommand{\smallspace}{\vspace{0.1cm}}
\newcommand{\mediumspace}{\vspace{0.3cm}}
\newcommand{\largespace}{\vspace{6cm}}

% Tamaños de fuente auxiliares
\newcommand\fontvia{\fontsize{4}{6}\selectfont}
\newcommand\fontvib{\fontsize{5}{8}\selectfont}
\newcommand\fontvic{\fontsize{7}{10}\selectfont}
\newcommand\fontvid{\fontsize{8}{11}\selectfont}
\newcommand\fontvie{\fontsize{10}{11}\selectfont}

% Celdas de mapa de calor estilizado para la intuición de shift-share
% Monotone (brown) cell: value 0..100 -> brown intensity
\newcommand{\sharecell}[3]{%
  \pgfmathtruncatemacro{\vv}{#3}%
  \fill[brown!\vv!white] (#1,#2) rectangle ({#1+1},{#2+1});%
  \draw[gray!40] (#1,#2) rectangle ({#1+1},{#2+1});%
}
% Diverging (blue/red) cell: value -50..50 -> blue (neg) or red (pos)
\newcommand{\divcell}[3]{%
  \pgfmathtruncatemacro{\absvv}{abs(#3)*2}%
  \pgfmathtruncatemacro{\negflag}{(#3)<0 ? 1 : 0}%
  \ifnum\negflag=1
    \fill[blue!\absvv!white] (#1,#2) rectangle ({#1+1},{#2+1});%
  \else
    \fill[red!\absvv!white] (#1,#2) rectangle ({#1+1},{#2+1});%
  \fi
  \draw[gray!40] (#1,#2) rectangle ({#1+1},{#2+1});%
}
% Uniform cell: same color regardless of (col,row)
\newcommand{\unifcell}[3]{%
  \pgfmathtruncatemacro{\vv}{#3}%
  \fill[red!\vv!white] (#1,#2) rectangle ({#1+1},{#2+1});%
  \draw[gray!40] (#1,#2) rectangle ({#1+1},{#2+1});%
}

% =====================================================================
% 13. Datos de portada (comunes a todas las clases)
% =====================================================================
\title{Econometría Avanzada}
\author{Manuel Fernández}
\date{\today}
%\titlegraphic{\hfill\includegraphics[height=1.5cm]{logo.pdf}}

%   3) \subtitle{...}                   (título de la clase puntual)
% Todo lo demás (paquetes, macros, tema, título, autor, hypersetup, etc.)
% vive aquí para poder hacer cambios globales en un solo lugar.
% =====================================================================

% ---------------------------------------------------------------------
% 1. Codificación, idioma y fuentes
% ---------------------------------------------------------------------
\usepackage{etex}
\usepackage[utf8]{inputenc}
\usepackage[T1]{fontenc}
\usepackage[spanish,es-nodecimaldot]{babel}
\usepackage{lmodern}
\usepackage{type1cm}
\usepackage{textcomp}

% ---------------------------------------------------------------------
% 2. Matemáticas y símbolos
% ---------------------------------------------------------------------
\usepackage{amsmath,amssymb,amsthm,amsfonts}
\usepackage{bbm}
\usepackage{bm}

% ---------------------------------------------------------------------
% 3. Bibliografía y citas
% ---------------------------------------------------------------------
\usepackage[round]{natbib}
\bibliographystyle{erae}
\usepackage{bibentry}
\usepackage{url}
\let\htmlstyloaded=\relax

% Las citas se muestran en marrón para distinguirlas del contenido.
\let\oldcitet=\citet
\renewcommand{\citet}[1]{\textcolor{brown}{\oldcitet{#1}}}
\let\oldcitep=\citep
\renewcommand{\citep}[1]{\textcolor{brown}{\oldcitep{#1}}}
\let\oldcite=\cite
\renewcommand{\cite}[1]{\textcolor{brown}{\oldcite{#1}}}

% ---------------------------------------------------------------------
% 4. Figuras, tablas y colores
% ---------------------------------------------------------------------
\usepackage{float}
\usepackage{graphicx}
\usepackage{subcaption}
\usepackage[font={footnotesize}]{caption}
\usepackage{lscape}
\usepackage{longtable}
\usepackage{booktabs}
\usepackage{multirow}
\usepackage{array}
\usepackage[flushleft]{threeparttable}
\usepackage{color, colortbl}
\usepackage{eso-pic}
\usepackage{chngcntr}
\usepackage{adjustbox}
\usepackage{pdfpages}
\usepackage{media9}

% siunitx: usado para alinear/formatear números en tablas
\usepackage{siunitx}
\sisetup{
	detect-mode,
	tight-spacing		= true,
	group-digits		= false ,
	input-signs		= ,
	input-symbols		= ( ) [ ] - + *,
	input-open-uncertainty	= ,
	input-close-uncertainty	= ,
	table-align-text-post	= false
}

% ---------------------------------------------------------------------
% 5. TikZ / PGFPlots
% ---------------------------------------------------------------------
\usepackage{tikz}
\usepackage{pgfplots}
\usepgfplotslibrary{dateplot}
\usetikzlibrary{arrows,calc, patterns, positioning, shapes.geometric, decorations.pathreplacing,decorations.markings}

% ---------------------------------------------------------------------
% 6. Misceláneos de formato
% ---------------------------------------------------------------------
\usepackage{setspace}
\usepackage{lipsum}
\usepackage{framed}
\usepackage{mdframed}
\usepackage{fancyhdr}
\usepackage{beamerfoils}
\usepackage{pifont}
\usepackage{ragged2e}
\usepackage{xspace}
\usepackage{appendixnumberbeamer}
\usepackage[scale=2]{ccicons}
\usepackage[most]{tcolorbox}
\usepackage{enumitem}
\setbeamertemplate{note page}[plain]

% ---------------------------------------------------------------------
% 7. Hyperref (debe cargarse después de la mayoría de paquetes)
% ---------------------------------------------------------------------
\usepackage{hyperref}
\hypersetup {
	pdftitle={General},
	pdfauthor={Manuel Fern\'{a}ndez},
	colorlinks=true,
	citecolor=brown,
	linkcolor=black,
	urlcolor=blue
}
\makeatletter
\renewcommand{\hyper@natlinkstart}[1]{}
\renewcommand{\hyper@natlinkend}{}
\renewcommand{\hyper@natlinkbreak}[2]{#1}
\makeatother

% =====================================================================
% 8. Tema Beamer (metropolis)
% =====================================================================
\usetheme{metropolis}
\newcommand{\themename}{\textbf{\textsc{metropolis}}\xspace}
\metroset{sectionpage=none}

\setbeamerfont{citep}{family=\rm}
\setbeamerfont{caption}{size=\scriptsize}
\setbeamerfont{normal text}{size=7pt}

%\setcounter{tocdepth}{1}
\setbeamertemplate{section in toc}[sections numbered]
\AtBeginSection[]{
  \begin{frame}
  \vfill
  \centering
  \begin{beamercolorbox}[sep=8pt,center,shadow=true,rounded=true]{title}
    \usebeamerfont{title}\insertsectionhead\par%
  \end{beamercolorbox}
  \vfill
  \end{frame}
}
% NOTA: \setbeamercolor{button}, \setbeamercovered{invisible} y
% \AtBeginSubsection NO están aquí porque sólo algunas clases los usaban
% en el original (afectan botones de hipervínculo, el efecto de \pause,
% y los separadores de \subsection). Se preservan como "ajustes locales"
% en los .tex que ya los tenían, justo después de \input{preambulo}.

% =====================================================================
% 9. Macros de Estout (tablas de resultados)
% =====================================================================
\newcommand{\sym}[1]{\rlap{#1}}% Thanks to David Carlisle
\let\estinput=\input% define a new input command so that we can still flatten the document
\newcommand{\estwide}[3]{
	\vspace{.75ex}{
		\begin{tabular*}
			{\textwidth}{@{\hskip\tabcolsep\extracolsep\fill}l*{#2}{#3}}
			\toprule
			\estinput{#1}
			\bottomrule
			\addlinespace[.75ex]
		\end{tabular*}
	}
}
\newcommand{\estauto}[3]{
	\vspace{.75ex}{
		\begin{tabular}{l*{#2}{#3}}
			\toprule
			\estinput{#1}
			\bottomrule
			\addlinespace[.75ex]
		\end{tabular}
	}
}
% Allow line breaks with \\ in specialcells
\newcommand{\specialcell}[2][c]{%
	\begin{tabular}[#1]{@{}c@{}}#2\end{tabular}}

% =====================================================================
% 10. Notas y fuentes al pie de figuras/tablas
% =====================================================================
% Note/Source/Text after Tables
\newcommand{\figtext}[1]{
	\vspace{-1.9ex}
	\captionsetup{justification=justified,font=footnotesize}
	\caption*{\hspace{6pt}\hangindent=1.5em #1}
}
\newcommand{\fignote}[1]{\figtext{\emph{Note:~}~#1}}
\newcommand{\figsource}[1]{\figtext{\emph{Source:~}~#1}}
% Add significance note with \starnote
\newcommand{\starnote}{\figtext{* p $<$ 0.10 ** p $<$ 0.05 *** p $<$ 0.01. }}

% =====================================================================
% 11. Entornos tipo teorema (en español)
% =====================================================================
\newtheorem{Prop}{Proposition}[section]
\newtheorem{Def}{Definition}
%\newtheorem{Assum}{Assumption}
%\newtheorem{Question}{Question}
%\newtheorem{Comment}{Comment}
%\theoremstyle{plain}
%\numberwithin{Assum}{section}
%\numberwithin{figure}{section}
%\numberwithin{equation}{section}
%\numberwithin{Def}{section}

\newenvironment{Teorema}[1][Teorema.]{\begin{trivlist}
		\item[\hskip \labelsep {\bfseries #1}]}{\end{trivlist}}
\newenvironment{Corolario}[1][Corolario.]{\begin{trivlist}
		\item[\hskip \labelsep {\bfseries #1}]}{\end{trivlist}}
\newenvironment{Definicion}[1][Definición.]{\begin{trivlist}
		\item[\hskip \labelsep {\bfseries #1}]}{\end{trivlist}}
\newenvironment{Lema}[1][Lema.]{\begin{trivlist}
		\item[\hskip \labelsep {\bfseries #1}]}{\end{trivlist}}
\newenvironment{Ejemplo}[1][Ejemplo.]{\begin{trivlist}
		\item[\hskip \labelsep {\bfseries #1}]}{\end{trivlist}}

\newenvironment{myindentpar}[1]%
{\begin{list}{}%
		{\setlength{\leftmargin}{#1}}%
		\item[]%
	}
	{\end{list}}

% =====================================================================
% 12. Notación matemática y econométrica
% =====================================================================
\newcommand{\plim}{\operatornamewithlimits{plim}}
\newcommand{\Var}{\operatorname{Var}}
\newcommand{\pderiv}[2]{\frac{\partial #1}{\partial #2}}
\newcommand{\spderiv}[2]{\frac{\partial^2 #1}{\partial #2^2}}
\newcommand{\xderiv}[3]{\frac{\partial^2 #1}{\partial #2 \partial #3}}
\newcommand{\argmin}{\operatornamewithlimits{arg\,min}}
\newcommand{\argmax}{\operatornamewithlimits{arg\,max}}
\newcommand{\amax}{\operatornamewithlimits{max}}
\newcommand{\amin}{\operatornamewithlimits{min}}
\newcommand*\circled[1]{\tikz[baseline=(char.base)]{\node[shape=circle,draw,inner sep=1pt] (char) {#1};}}
\newcommand{\indep}{\perp \!\!\! \perp}
\makeatletter
\newcommand{\distas}[1]{\mathbin{\overset{#1}{\kern\z@\sim}}}%
\newsavebox{\mybox}\newsavebox{\mysim}
\newcommand{\distras}[1]{%
	\savebox{\mybox}{\hbox{\kern3pt$\scriptstyle#1$\kern3pt}}%
	\savebox{\mysim}{\hbox{$\sim$}}%
	\mathbin{\overset{#1}{\kern\z@\resizebox{\wd\mybox}{\ht\mysim}{$\sim$}}}%
}
\makeatother

\DeclareMathOperator{\EX}{\mathbb{E}}% expected value
\newcommand{\Prob}[1]{\text{Pr}\left[#1\right]}
\newcommand{\RK}[1]{\mathbb{R}^{#1}}
\newcommand{\HT}[1]{\mathbb{H}_{#1}}

% Vectores y matrices en negrilla (romana)
\newcommand{\xB}{\textbf{x}}
\newcommand{\yB}{\textbf{y}}
\newcommand{\zB}{\textbf{z}}
\newcommand{\wB}{\textbf{w}}
\newcommand{\eB}{\textbf{e}}
\newcommand{\vB}{\textbf{v}}
\newcommand{\AB}{\textbf{A}}
\newcommand{\XB}{\textbf{X}}
\newcommand{\YB}{\textbf{Y}}
\newcommand{\ZB}{\textbf{Z}}
\newcommand{\QB}{\textbf{Q}}
\newcommand{\WB}{\textbf{W}}

% Vectores griegos en negrilla
\newcommand{\betaB}{\boldsymbol{\beta}}
\newcommand{\gammaB}{\boldsymbol{\gamma}}
\newcommand{\alphaB}{\boldsymbol{\alpha}}
\newcommand{\thetaB}{\boldsymbol{\theta}}
\newcommand{\deltaB}{\boldsymbol{\delta}}
\newcommand{\phiB}{\boldsymbol{\phi}}
\newcommand{\epsilonB}{\boldsymbol{\epsilon}}
\newcommand{\etaB}{\boldsymbol{\eta}}
\newcommand{\betaBH}{\boldsymbol{\hat{\beta}}}
\newcommand{\betaH}{\hat{\beta}}

% Colores de énfasis
\newcommand{\brown}[1]{\textcolor{brown}{#1}}
\newcommand{\blue}[1]{\textcolor{blue}{#1}}
\newcommand{\red}[1]{\textcolor{red}{#1}}

% Espaciados usados en itemize/enumerate
\newcommand{\smallspace}{\vspace{0.1cm}}
\newcommand{\mediumspace}{\vspace{0.3cm}}
\newcommand{\largespace}{\vspace{6cm}}

% Tamaños de fuente auxiliares
\newcommand\fontvia{\fontsize{4}{6}\selectfont}
\newcommand\fontvib{\fontsize{5}{8}\selectfont}
\newcommand\fontvic{\fontsize{7}{10}\selectfont}
\newcommand\fontvid{\fontsize{8}{11}\selectfont}
\newcommand\fontvie{\fontsize{10}{11}\selectfont}

% Celdas de mapa de calor estilizado para la intuición de shift-share
% Monotone (brown) cell: value 0..100 -> brown intensity
\newcommand{\sharecell}[3]{%
  \pgfmathtruncatemacro{\vv}{#3}%
  \fill[brown!\vv!white] (#1,#2) rectangle ({#1+1},{#2+1});%
  \draw[gray!40] (#1,#2) rectangle ({#1+1},{#2+1});%
}
% Diverging (blue/red) cell: value -50..50 -> blue (neg) or red (pos)
\newcommand{\divcell}[3]{%
  \pgfmathtruncatemacro{\absvv}{abs(#3)*2}%
  \pgfmathtruncatemacro{\negflag}{(#3)<0 ? 1 : 0}%
  \ifnum\negflag=1
    \fill[blue!\absvv!white] (#1,#2) rectangle ({#1+1},{#2+1});%
  \else
    \fill[red!\absvv!white] (#1,#2) rectangle ({#1+1},{#2+1});%
  \fi
  \draw[gray!40] (#1,#2) rectangle ({#1+1},{#2+1});%
}
% Uniform cell: same color regardless of (col,row)
\newcommand{\unifcell}[3]{%
  \pgfmathtruncatemacro{\vv}{#3}%
  \fill[red!\vv!white] (#1,#2) rectangle ({#1+1},{#2+1});%
  \draw[gray!40] (#1,#2) rectangle ({#1+1},{#2+1});%
}

% =====================================================================
% 13. Datos de portada (comunes a todas las clases)
% =====================================================================
\title{Econometría Avanzada}
\author{Manuel Fernández}
\date{\today}
%\titlegraphic{\hfill\includegraphics[height=1.5cm]{logo.pdf}}

%   3) \subtitle{...}                   (título de la clase puntual)
% Todo lo demás (paquetes, macros, tema, título, autor, hypersetup, etc.)
% vive aquí para poder hacer cambios globales en un solo lugar.
% =====================================================================

% ---------------------------------------------------------------------
% 1. Codificación, idioma y fuentes
% ---------------------------------------------------------------------
\usepackage{etex}
\usepackage[utf8]{inputenc}
\usepackage[T1]{fontenc}
\usepackage[spanish,es-nodecimaldot]{babel}
\usepackage{lmodern}
\usepackage{type1cm}
\usepackage{textcomp}

% ---------------------------------------------------------------------
% 2. Matemáticas y símbolos
% ---------------------------------------------------------------------
\usepackage{amsmath,amssymb,amsthm,amsfonts}
\usepackage{bbm}
\usepackage{bm}

% ---------------------------------------------------------------------
% 3. Bibliografía y citas
% ---------------------------------------------------------------------
\usepackage[round]{natbib}
\bibliographystyle{erae}
\usepackage{bibentry}
\usepackage{url}
\let\htmlstyloaded=\relax

% Las citas se muestran en marrón para distinguirlas del contenido.
\let\oldcitet=\citet
\renewcommand{\citet}[1]{\textcolor{brown}{\oldcitet{#1}}}
\let\oldcitep=\citep
\renewcommand{\citep}[1]{\textcolor{brown}{\oldcitep{#1}}}
\let\oldcite=\cite
\renewcommand{\cite}[1]{\textcolor{brown}{\oldcite{#1}}}

% ---------------------------------------------------------------------
% 4. Figuras, tablas y colores
% ---------------------------------------------------------------------
\usepackage{float}
\usepackage{graphicx}
\usepackage{subcaption}
\usepackage[font={footnotesize}]{caption}
\usepackage{lscape}
\usepackage{longtable}
\usepackage{booktabs}
\usepackage{multirow}
\usepackage{array}
\usepackage[flushleft]{threeparttable}
\usepackage{color, colortbl}
\usepackage{eso-pic}
\usepackage{chngcntr}
\usepackage{adjustbox}
\usepackage{pdfpages}
\usepackage{media9}

% siunitx: usado para alinear/formatear números en tablas
\usepackage{siunitx}
\sisetup{
	detect-mode,
	tight-spacing		= true,
	group-digits		= false ,
	input-signs		= ,
	input-symbols		= ( ) [ ] - + *,
	input-open-uncertainty	= ,
	input-close-uncertainty	= ,
	table-align-text-post	= false
}

% ---------------------------------------------------------------------
% 5. TikZ / PGFPlots
% ---------------------------------------------------------------------
\usepackage{tikz}
\usepackage{pgfplots}
\usepgfplotslibrary{dateplot}
\usetikzlibrary{arrows,calc, patterns, positioning, shapes.geometric, decorations.pathreplacing,decorations.markings}

% ---------------------------------------------------------------------
% 6. Misceláneos de formato
% ---------------------------------------------------------------------
\usepackage{setspace}
\usepackage{lipsum}
\usepackage{framed}
\usepackage{mdframed}
\usepackage{fancyhdr}
\usepackage{beamerfoils}
\usepackage{pifont}
\usepackage{ragged2e}
\usepackage{xspace}
\usepackage{appendixnumberbeamer}
\usepackage[scale=2]{ccicons}
\usepackage[most]{tcolorbox}
\usepackage{enumitem}
\setbeamertemplate{note page}[plain]

% ---------------------------------------------------------------------
% 7. Hyperref (debe cargarse después de la mayoría de paquetes)
% ---------------------------------------------------------------------
\usepackage{hyperref}
\hypersetup {
	pdftitle={General},
	pdfauthor={Manuel Fern\'{a}ndez},
	colorlinks=true,
	citecolor=brown,
	linkcolor=black,
	urlcolor=blue
}
\makeatletter
\renewcommand{\hyper@natlinkstart}[1]{}
\renewcommand{\hyper@natlinkend}{}
\renewcommand{\hyper@natlinkbreak}[2]{#1}
\makeatother

% =====================================================================
% 8. Tema Beamer (metropolis)
% =====================================================================
\usetheme{metropolis}
\newcommand{\themename}{\textbf{\textsc{metropolis}}\xspace}
\metroset{sectionpage=none}

\setbeamerfont{citep}{family=\rm}
\setbeamerfont{caption}{size=\scriptsize}
\setbeamerfont{normal text}{size=7pt}

%\setcounter{tocdepth}{1}
\setbeamertemplate{section in toc}[sections numbered]
\AtBeginSection[]{
  \begin{frame}
  \vfill
  \centering
  \begin{beamercolorbox}[sep=8pt,center,shadow=true,rounded=true]{title}
    \usebeamerfont{title}\insertsectionhead\par%
  \end{beamercolorbox}
  \vfill
  \end{frame}
}
% NOTA: \setbeamercolor{button}, \setbeamercovered{invisible} y
% \AtBeginSubsection NO están aquí porque sólo algunas clases los usaban
% en el original (afectan botones de hipervínculo, el efecto de \pause,
% y los separadores de \subsection). Se preservan como "ajustes locales"
% en los .tex que ya los tenían, justo después de \input{preambulo}.

% =====================================================================
% 9. Macros de Estout (tablas de resultados)
% =====================================================================
\newcommand{\sym}[1]{\rlap{#1}}% Thanks to David Carlisle
\let\estinput=\input% define a new input command so that we can still flatten the document
\newcommand{\estwide}[3]{
	\vspace{.75ex}{
		\begin{tabular*}
			{\textwidth}{@{\hskip\tabcolsep\extracolsep\fill}l*{#2}{#3}}
			\toprule
			\estinput{#1}
			\bottomrule
			\addlinespace[.75ex]
		\end{tabular*}
	}
}
\newcommand{\estauto}[3]{
	\vspace{.75ex}{
		\begin{tabular}{l*{#2}{#3}}
			\toprule
			\estinput{#1}
			\bottomrule
			\addlinespace[.75ex]
		\end{tabular}
	}
}
% Allow line breaks with \\ in specialcells
\newcommand{\specialcell}[2][c]{%
	\begin{tabular}[#1]{@{}c@{}}#2\end{tabular}}

% =====================================================================
% 10. Notas y fuentes al pie de figuras/tablas
% =====================================================================
% Note/Source/Text after Tables
\newcommand{\figtext}[1]{
	\vspace{-1.9ex}
	\captionsetup{justification=justified,font=footnotesize}
	\caption*{\hspace{6pt}\hangindent=1.5em #1}
}
\newcommand{\fignote}[1]{\figtext{\emph{Note:~}~#1}}
\newcommand{\figsource}[1]{\figtext{\emph{Source:~}~#1}}
% Add significance note with \starnote
\newcommand{\starnote}{\figtext{* p $<$ 0.10 ** p $<$ 0.05 *** p $<$ 0.01. }}

% =====================================================================
% 11. Entornos tipo teorema (en español)
% =====================================================================
\newtheorem{Prop}{Proposition}[section]
\newtheorem{Def}{Definition}
%\newtheorem{Assum}{Assumption}
%\newtheorem{Question}{Question}
%\newtheorem{Comment}{Comment}
%\theoremstyle{plain}
%\numberwithin{Assum}{section}
%\numberwithin{figure}{section}
%\numberwithin{equation}{section}
%\numberwithin{Def}{section}

\newenvironment{Teorema}[1][Teorema.]{\begin{trivlist}
		\item[\hskip \labelsep {\bfseries #1}]}{\end{trivlist}}
\newenvironment{Corolario}[1][Corolario.]{\begin{trivlist}
		\item[\hskip \labelsep {\bfseries #1}]}{\end{trivlist}}
\newenvironment{Definicion}[1][Definición.]{\begin{trivlist}
		\item[\hskip \labelsep {\bfseries #1}]}{\end{trivlist}}
\newenvironment{Lema}[1][Lema.]{\begin{trivlist}
		\item[\hskip \labelsep {\bfseries #1}]}{\end{trivlist}}
\newenvironment{Ejemplo}[1][Ejemplo.]{\begin{trivlist}
		\item[\hskip \labelsep {\bfseries #1}]}{\end{trivlist}}

\newenvironment{myindentpar}[1]%
{\begin{list}{}%
		{\setlength{\leftmargin}{#1}}%
		\item[]%
	}
	{\end{list}}

% =====================================================================
% 12. Notación matemática y econométrica
% =====================================================================
\newcommand{\plim}{\operatornamewithlimits{plim}}
\newcommand{\Var}{\operatorname{Var}}
\newcommand{\pderiv}[2]{\frac{\partial #1}{\partial #2}}
\newcommand{\spderiv}[2]{\frac{\partial^2 #1}{\partial #2^2}}
\newcommand{\xderiv}[3]{\frac{\partial^2 #1}{\partial #2 \partial #3}}
\newcommand{\argmin}{\operatornamewithlimits{arg\,min}}
\newcommand{\argmax}{\operatornamewithlimits{arg\,max}}
\newcommand{\amax}{\operatornamewithlimits{max}}
\newcommand{\amin}{\operatornamewithlimits{min}}
\newcommand*\circled[1]{\tikz[baseline=(char.base)]{\node[shape=circle,draw,inner sep=1pt] (char) {#1};}}
\newcommand{\indep}{\perp \!\!\! \perp}
\makeatletter
\newcommand{\distas}[1]{\mathbin{\overset{#1}{\kern\z@\sim}}}%
\newsavebox{\mybox}\newsavebox{\mysim}
\newcommand{\distras}[1]{%
	\savebox{\mybox}{\hbox{\kern3pt$\scriptstyle#1$\kern3pt}}%
	\savebox{\mysim}{\hbox{$\sim$}}%
	\mathbin{\overset{#1}{\kern\z@\resizebox{\wd\mybox}{\ht\mysim}{$\sim$}}}%
}
\makeatother

\DeclareMathOperator{\EX}{\mathbb{E}}% expected value
\newcommand{\Prob}[1]{\text{Pr}\left[#1\right]}
\newcommand{\RK}[1]{\mathbb{R}^{#1}}
\newcommand{\HT}[1]{\mathbb{H}_{#1}}

% Vectores y matrices en negrilla (romana)
\newcommand{\xB}{\textbf{x}}
\newcommand{\yB}{\textbf{y}}
\newcommand{\zB}{\textbf{z}}
\newcommand{\wB}{\textbf{w}}
\newcommand{\eB}{\textbf{e}}
\newcommand{\vB}{\textbf{v}}
\newcommand{\AB}{\textbf{A}}
\newcommand{\XB}{\textbf{X}}
\newcommand{\YB}{\textbf{Y}}
\newcommand{\ZB}{\textbf{Z}}
\newcommand{\QB}{\textbf{Q}}
\newcommand{\WB}{\textbf{W}}

% Vectores griegos en negrilla
\newcommand{\betaB}{\boldsymbol{\beta}}
\newcommand{\gammaB}{\boldsymbol{\gamma}}
\newcommand{\alphaB}{\boldsymbol{\alpha}}
\newcommand{\thetaB}{\boldsymbol{\theta}}
\newcommand{\deltaB}{\boldsymbol{\delta}}
\newcommand{\phiB}{\boldsymbol{\phi}}
\newcommand{\epsilonB}{\boldsymbol{\epsilon}}
\newcommand{\etaB}{\boldsymbol{\eta}}
\newcommand{\betaBH}{\boldsymbol{\hat{\beta}}}
\newcommand{\betaH}{\hat{\beta}}

% Colores de énfasis
\newcommand{\brown}[1]{\textcolor{brown}{#1}}
\newcommand{\blue}[1]{\textcolor{blue}{#1}}
\newcommand{\red}[1]{\textcolor{red}{#1}}

% Espaciados usados en itemize/enumerate
\newcommand{\smallspace}{\vspace{0.1cm}}
\newcommand{\mediumspace}{\vspace{0.3cm}}
\newcommand{\largespace}{\vspace{6cm}}

% Tamaños de fuente auxiliares
\newcommand\fontvia{\fontsize{4}{6}\selectfont}
\newcommand\fontvib{\fontsize{5}{8}\selectfont}
\newcommand\fontvic{\fontsize{7}{10}\selectfont}
\newcommand\fontvid{\fontsize{8}{11}\selectfont}
\newcommand\fontvie{\fontsize{10}{11}\selectfont}

% Celdas de mapa de calor estilizado para la intuición de shift-share
% Monotone (brown) cell: value 0..100 -> brown intensity
\newcommand{\sharecell}[3]{%
  \pgfmathtruncatemacro{\vv}{#3}%
  \fill[brown!\vv!white] (#1,#2) rectangle ({#1+1},{#2+1});%
  \draw[gray!40] (#1,#2) rectangle ({#1+1},{#2+1});%
}
% Diverging (blue/red) cell: value -50..50 -> blue (neg) or red (pos)
\newcommand{\divcell}[3]{%
  \pgfmathtruncatemacro{\absvv}{abs(#3)*2}%
  \pgfmathtruncatemacro{\negflag}{(#3)<0 ? 1 : 0}%
  \ifnum\negflag=1
    \fill[blue!\absvv!white] (#1,#2) rectangle ({#1+1},{#2+1});%
  \else
    \fill[red!\absvv!white] (#1,#2) rectangle ({#1+1},{#2+1});%
  \fi
  \draw[gray!40] (#1,#2) rectangle ({#1+1},{#2+1});%
}
% Uniform cell: same color regardless of (col,row)
\newcommand{\unifcell}[3]{%
  \pgfmathtruncatemacro{\vv}{#3}%
  \fill[red!\vv!white] (#1,#2) rectangle ({#1+1},{#2+1});%
  \draw[gray!40] (#1,#2) rectangle ({#1+1},{#2+1});%
}

% =====================================================================
% 13. Datos de portada (comunes a todas las clases)
% =====================================================================
\title{Econometría Avanzada}
\author{Manuel Fernández}
\date{\today}
%\titlegraphic{\hfill\includegraphics[height=1.5cm]{logo.pdf}}

%   3) \subtitle{...}                   (título de la clase puntual)
% Todo lo demás (paquetes, macros, tema, título, autor, hypersetup, etc.)
% vive aquí para poder hacer cambios globales en un solo lugar.
% =====================================================================

% ---------------------------------------------------------------------
% 1. Codificación, idioma y fuentes
% ---------------------------------------------------------------------
\usepackage{etex}
\usepackage[utf8]{inputenc}
\usepackage[T1]{fontenc}
\usepackage[spanish,es-nodecimaldot]{babel}
\usepackage{lmodern}
\usepackage{type1cm}
\usepackage{textcomp}

% ---------------------------------------------------------------------
% 2. Matemáticas y símbolos
% ---------------------------------------------------------------------
\usepackage{amsmath,amssymb,amsthm,amsfonts}
\usepackage{bbm}
\usepackage{bm}

% ---------------------------------------------------------------------
% 3. Bibliografía y citas
% ---------------------------------------------------------------------
\usepackage[round]{natbib}
\bibliographystyle{erae}
\usepackage{bibentry}
\usepackage{url}
\let\htmlstyloaded=\relax

% Las citas se muestran en marrón para distinguirlas del contenido.
\let\oldcitet=\citet
\renewcommand{\citet}[1]{\textcolor{brown}{\oldcitet{#1}}}
\let\oldcitep=\citep
\renewcommand{\citep}[1]{\textcolor{brown}{\oldcitep{#1}}}
\let\oldcite=\cite
\renewcommand{\cite}[1]{\textcolor{brown}{\oldcite{#1}}}

% ---------------------------------------------------------------------
% 4. Figuras, tablas y colores
% ---------------------------------------------------------------------
\usepackage{float}
\usepackage{graphicx}
\usepackage{subcaption}
\usepackage[font={footnotesize}]{caption}
\usepackage{lscape}
\usepackage{longtable}
\usepackage{booktabs}
\usepackage{multirow}
\usepackage{array}
\usepackage[flushleft]{threeparttable}
\usepackage{color, colortbl}
\usepackage{eso-pic}
\usepackage{chngcntr}
\usepackage{adjustbox}
\usepackage{pdfpages}
\usepackage{media9}

% siunitx: usado para alinear/formatear números en tablas
\usepackage{siunitx}
\sisetup{
	detect-mode,
	tight-spacing		= true,
	group-digits		= false ,
	input-signs		= ,
	input-symbols		= ( ) [ ] - + *,
	input-open-uncertainty	= ,
	input-close-uncertainty	= ,
	table-align-text-post	= false
}

% ---------------------------------------------------------------------
% 5. TikZ / PGFPlots
% ---------------------------------------------------------------------
\usepackage{tikz}
\usepackage{pgfplots}
\usepgfplotslibrary{dateplot}
\usetikzlibrary{arrows,calc, patterns, positioning, shapes.geometric, decorations.pathreplacing,decorations.markings}

% ---------------------------------------------------------------------
% 6. Misceláneos de formato
% ---------------------------------------------------------------------
\usepackage{setspace}
\usepackage{lipsum}
\usepackage{framed}
\usepackage{mdframed}
\usepackage{fancyhdr}
\usepackage{beamerfoils}
\usepackage{pifont}
\usepackage{ragged2e}
\usepackage{xspace}
\usepackage{appendixnumberbeamer}
\usepackage[scale=2]{ccicons}
\usepackage[most]{tcolorbox}
\usepackage{enumitem}
\setbeamertemplate{note page}[plain]

% ---------------------------------------------------------------------
% 7. Hyperref (debe cargarse después de la mayoría de paquetes)
% ---------------------------------------------------------------------
\usepackage{hyperref}
\hypersetup {
	pdftitle={General},
	pdfauthor={Manuel Fern\'{a}ndez},
	colorlinks=true,
	citecolor=brown,
	linkcolor=black,
	urlcolor=blue
}
\makeatletter
\renewcommand{\hyper@natlinkstart}[1]{}
\renewcommand{\hyper@natlinkend}{}
\renewcommand{\hyper@natlinkbreak}[2]{#1}
\makeatother

% =====================================================================
% 8. Tema Beamer (metropolis)
% =====================================================================
\usetheme{metropolis}
\newcommand{\themename}{\textbf{\textsc{metropolis}}\xspace}
\metroset{sectionpage=none}

\setbeamerfont{citep}{family=\rm}
\setbeamerfont{caption}{size=\scriptsize}
\setbeamerfont{normal text}{size=7pt}

%\setcounter{tocdepth}{1}
\setbeamertemplate{section in toc}[sections numbered]
\AtBeginSection[]{
  \begin{frame}
  \vfill
  \centering
  \begin{beamercolorbox}[sep=8pt,center,shadow=true,rounded=true]{title}
    \usebeamerfont{title}\insertsectionhead\par%
  \end{beamercolorbox}
  \vfill
  \end{frame}
}
% NOTA: \setbeamercolor{button}, \setbeamercovered{invisible} y
% \AtBeginSubsection NO están aquí porque sólo algunas clases los usaban
% en el original (afectan botones de hipervínculo, el efecto de \pause,
% y los separadores de \subsection). Se preservan como "ajustes locales"
% en los .tex que ya los tenían, justo después de \input{preambulo}.

% =====================================================================
% 9. Macros de Estout (tablas de resultados)
% =====================================================================
\newcommand{\sym}[1]{\rlap{#1}}% Thanks to David Carlisle
\let\estinput=\input% define a new input command so that we can still flatten the document
\newcommand{\estwide}[3]{
	\vspace{.75ex}{
		\begin{tabular*}
			{\textwidth}{@{\hskip\tabcolsep\extracolsep\fill}l*{#2}{#3}}
			\toprule
			\estinput{#1}
			\bottomrule
			\addlinespace[.75ex]
		\end{tabular*}
	}
}
\newcommand{\estauto}[3]{
	\vspace{.75ex}{
		\begin{tabular}{l*{#2}{#3}}
			\toprule
			\estinput{#1}
			\bottomrule
			\addlinespace[.75ex]
		\end{tabular}
	}
}
% Allow line breaks with \\ in specialcells
\newcommand{\specialcell}[2][c]{%
	\begin{tabular}[#1]{@{}c@{}}#2\end{tabular}}

% =====================================================================
% 10. Notas y fuentes al pie de figuras/tablas
% =====================================================================
% Note/Source/Text after Tables
\newcommand{\figtext}[1]{
	\vspace{-1.9ex}
	\captionsetup{justification=justified,font=footnotesize}
	\caption*{\hspace{6pt}\hangindent=1.5em #1}
}
\newcommand{\fignote}[1]{\figtext{\emph{Note:~}~#1}}
\newcommand{\figsource}[1]{\figtext{\emph{Source:~}~#1}}
% Add significance note with \starnote
\newcommand{\starnote}{\figtext{* p $<$ 0.10 ** p $<$ 0.05 *** p $<$ 0.01. }}

% =====================================================================
% 11. Entornos tipo teorema (en español)
% =====================================================================
\newtheorem{Prop}{Proposition}[section]
\newtheorem{Def}{Definition}
%\newtheorem{Assum}{Assumption}
%\newtheorem{Question}{Question}
%\newtheorem{Comment}{Comment}
%\theoremstyle{plain}
%\numberwithin{Assum}{section}
%\numberwithin{figure}{section}
%\numberwithin{equation}{section}
%\numberwithin{Def}{section}

\newenvironment{Teorema}[1][Teorema.]{\begin{trivlist}
		\item[\hskip \labelsep {\bfseries #1}]}{\end{trivlist}}
\newenvironment{Corolario}[1][Corolario.]{\begin{trivlist}
		\item[\hskip \labelsep {\bfseries #1}]}{\end{trivlist}}
\newenvironment{Definicion}[1][Definición.]{\begin{trivlist}
		\item[\hskip \labelsep {\bfseries #1}]}{\end{trivlist}}
\newenvironment{Lema}[1][Lema.]{\begin{trivlist}
		\item[\hskip \labelsep {\bfseries #1}]}{\end{trivlist}}
\newenvironment{Ejemplo}[1][Ejemplo.]{\begin{trivlist}
		\item[\hskip \labelsep {\bfseries #1}]}{\end{trivlist}}

\newenvironment{myindentpar}[1]%
{\begin{list}{}%
		{\setlength{\leftmargin}{#1}}%
		\item[]%
	}
	{\end{list}}

% =====================================================================
% 12. Notación matemática y econométrica
% =====================================================================
\newcommand{\plim}{\operatornamewithlimits{plim}}
\newcommand{\Var}{\operatorname{Var}}
\newcommand{\pderiv}[2]{\frac{\partial #1}{\partial #2}}
\newcommand{\spderiv}[2]{\frac{\partial^2 #1}{\partial #2^2}}
\newcommand{\xderiv}[3]{\frac{\partial^2 #1}{\partial #2 \partial #3}}
\newcommand{\argmin}{\operatornamewithlimits{arg\,min}}
\newcommand{\argmax}{\operatornamewithlimits{arg\,max}}
\newcommand{\amax}{\operatornamewithlimits{max}}
\newcommand{\amin}{\operatornamewithlimits{min}}
\newcommand*\circled[1]{\tikz[baseline=(char.base)]{\node[shape=circle,draw,inner sep=1pt] (char) {#1};}}
\newcommand{\indep}{\perp \!\!\! \perp}
\makeatletter
\newcommand{\distas}[1]{\mathbin{\overset{#1}{\kern\z@\sim}}}%
\newsavebox{\mybox}\newsavebox{\mysim}
\newcommand{\distras}[1]{%
	\savebox{\mybox}{\hbox{\kern3pt$\scriptstyle#1$\kern3pt}}%
	\savebox{\mysim}{\hbox{$\sim$}}%
	\mathbin{\overset{#1}{\kern\z@\resizebox{\wd\mybox}{\ht\mysim}{$\sim$}}}%
}
\makeatother

\DeclareMathOperator{\EX}{\mathbb{E}}% expected value
\newcommand{\Prob}[1]{\text{Pr}\left[#1\right]}
\newcommand{\RK}[1]{\mathbb{R}^{#1}}
\newcommand{\HT}[1]{\mathbb{H}_{#1}}

% Vectores y matrices en negrilla (romana)
\newcommand{\xB}{\textbf{x}}
\newcommand{\yB}{\textbf{y}}
\newcommand{\zB}{\textbf{z}}
\newcommand{\wB}{\textbf{w}}
\newcommand{\eB}{\textbf{e}}
\newcommand{\vB}{\textbf{v}}
\newcommand{\AB}{\textbf{A}}
\newcommand{\XB}{\textbf{X}}
\newcommand{\YB}{\textbf{Y}}
\newcommand{\ZB}{\textbf{Z}}
\newcommand{\QB}{\textbf{Q}}
\newcommand{\WB}{\textbf{W}}

% Vectores griegos en negrilla
\newcommand{\betaB}{\boldsymbol{\beta}}
\newcommand{\gammaB}{\boldsymbol{\gamma}}
\newcommand{\alphaB}{\boldsymbol{\alpha}}
\newcommand{\thetaB}{\boldsymbol{\theta}}
\newcommand{\deltaB}{\boldsymbol{\delta}}
\newcommand{\phiB}{\boldsymbol{\phi}}
\newcommand{\epsilonB}{\boldsymbol{\epsilon}}
\newcommand{\etaB}{\boldsymbol{\eta}}
\newcommand{\betaBH}{\boldsymbol{\hat{\beta}}}
\newcommand{\betaH}{\hat{\beta}}

% Colores de énfasis
\newcommand{\brown}[1]{\textcolor{brown}{#1}}
\newcommand{\blue}[1]{\textcolor{blue}{#1}}
\newcommand{\red}[1]{\textcolor{red}{#1}}

% Espaciados usados en itemize/enumerate
\newcommand{\smallspace}{\vspace{0.1cm}}
\newcommand{\mediumspace}{\vspace{0.3cm}}
\newcommand{\largespace}{\vspace{6cm}}

% Tamaños de fuente auxiliares
\newcommand\fontvia{\fontsize{4}{6}\selectfont}
\newcommand\fontvib{\fontsize{5}{8}\selectfont}
\newcommand\fontvic{\fontsize{7}{10}\selectfont}
\newcommand\fontvid{\fontsize{8}{11}\selectfont}
\newcommand\fontvie{\fontsize{10}{11}\selectfont}

% Celdas de mapa de calor estilizado para la intuición de shift-share
% Monotone (brown) cell: value 0..100 -> brown intensity
\newcommand{\sharecell}[3]{%
  \pgfmathtruncatemacro{\vv}{#3}%
  \fill[brown!\vv!white] (#1,#2) rectangle ({#1+1},{#2+1});%
  \draw[gray!40] (#1,#2) rectangle ({#1+1},{#2+1});%
}
% Diverging (blue/red) cell: value -50..50 -> blue (neg) or red (pos)
\newcommand{\divcell}[3]{%
  \pgfmathtruncatemacro{\absvv}{abs(#3)*2}%
  \pgfmathtruncatemacro{\negflag}{(#3)<0 ? 1 : 0}%
  \ifnum\negflag=1
    \fill[blue!\absvv!white] (#1,#2) rectangle ({#1+1},{#2+1});%
  \else
    \fill[red!\absvv!white] (#1,#2) rectangle ({#1+1},{#2+1});%
  \fi
  \draw[gray!40] (#1,#2) rectangle ({#1+1},{#2+1});%
}
% Uniform cell: same color regardless of (col,row)
\newcommand{\unifcell}[3]{%
  \pgfmathtruncatemacro{\vv}{#3}%
  \fill[red!\vv!white] (#1,#2) rectangle ({#1+1},{#2+1});%
  \draw[gray!40] (#1,#2) rectangle ({#1+1},{#2+1});%
}

% =====================================================================
% 13. Datos de portada (comunes a todas las clases)
% =====================================================================
\title{Econometría Avanzada}
\author{Manuel Fernández}
\date{\today}
%\titlegraphic{\hfill\includegraphics[height=1.5cm]{logo.pdf}}
